% Statistiques — Mathématiques 3ème — Cours signé OnBuch+
% Programme officiel MINESEC. Compiler avec : tectonic N14-statistiques.tex
\newcommand{\DOCMATIERE}{Mathématiques}
\newcommand{\DOCNIVEAU}{3ème}
\newcommand{\DOCMODULE}{Mathématiques – classe de 3ème}
\newcommand{\DOCLECON}{N14}
\newcommand{\DOCTITRE}{Statistiques}
% OnBuch+ — préambule « cours signés » ENRICHI (compilé avec tectonic / XeLaTeX)
% Système visuel « Pop » d'OnBuch+ : fond crème (jamais blanc), bordures encre
% épaisses, ombres « dures » décalées, titres Archivo Black en capitales.
% Version MATHÉMATIQUES : graphes (pgfplots/tikz), boîtes pédagogiques
% (définition, propriété, méthode, attention, le savais-tu, exercice résolu,
% exercice, TP, bilan), page de garde et signature OnBuch+ ancrée en bas.
%
% Macros attendues dans main.tex AVANT \input{preamble} :
%   \DOCMATIERE  \DOCNIVEAU  \DOCMODULE  \DOCLECON (numéro)  \DOCTITRE
\documentclass[11pt]{article}

\usepackage{fontspec}
\usepackage[french]{babel}
\usepackage[a4paper,margin=2cm,top=3.4cm,bottom=2.8cm,headheight=1.4cm,footskip=1.3cm]{geometry}
\usepackage{amsmath,amssymb,amsfonts}
\usepackage{mathtools} % \xmapsto, \coloneqq... souvent utilisés par les modèles
\usepackage{cancel} % simplifications barrées
% \abs{z} (module, valeur absolue) et \norm{u} (norme), utilisés par les modèles
\providecommand{\abs}{}\let\abs\relax\DeclarePairedDelimiter{\abs}{\lvert}{\rvert}
\providecommand{\norm}{}\let\norm\relax\DeclarePairedDelimiter{\norm}{\lVert}{\rVert}
\usepackage[table]{xcolor} % [table] : fournit aussi \rowcolors (alternance de lignes)
\usepackage{pagecolor}
\usepackage{tikz}
\usetikzlibrary{arrows.meta,positioning,calc,shapes.geometric,shapes.misc,shapes.arrows,decorations.pathmorphing,decorations.pathreplacing,decorations.markings,patterns,shadows,mindmap,trees,fit,backgrounds,matrix,intersections,angles,quotes}
\usepackage{pgfplots}
\usepackage{pgf-pie} % diagrammes circulaires \pie (statistiques)
\pgfplotsset{compat=1.18}
\usepgfplotslibrary{fillbetween,groupplots} % aires entre courbes (calcul intégral)
\usepackage{enumitem}
\usepackage{fancyhdr}
% [most] charge toutes les bibliothèques tcolorbox (listings, minted, raster,
% documentation...) et gonfle inutilement la mémoire TeX sur un document long
% (~300 boîtes) : on ne charge que ce qui est réellement utilisé.
\usepackage[skins,breakable]{tcolorbox}
\usepackage{graphicx}
\usepackage{colortbl}
\usepackage{booktabs}
\usepackage{tabularx}
\usepackage{diagbox} % case d'en-tête diagonale des tableaux à double entrée
% Colonnes extensibles centrée / gauche / droite (C, L, R), souvent utilisées par les modèles
\newcolumntype{C}{>{\centering\arraybackslash}X}
\newcolumntype{L}{>{\raggedright\arraybackslash}X}
\newcolumntype{R}{>{\raggedleft\arraybackslash}X}
\usepackage{array}
\usepackage{multicol}
\usepackage{siunitx}
% parse-numbers=false : accepte aussi \SI{2,5 \times 10^{-3}}{...}
\sisetup{locale=FR,output-decimal-marker={,},group-minimum-digits=4,parse-numbers=false}
\DeclareSIUnit\ohm{\ensuremath{\symup{\Omega}}} % U+2126 absent de la police
\usepackage{qrcode}
% \degree : commande fréquemment utilisée par les modèles pour un angle ou une
% température en dehors de \SI{}{} (non fournie par les paquets chargés).
\providecommand{\degree}{^{\circ}}
% \qed (fin de démonstration), utilisé par les modèles sans amsthm
\providecommand{\qed}{\hfill$\square$}
% \barre : double barre des valeurs interdites dans un tableau de variations (tkz-tab)
\providecommand{\barre}{\ensuremath{\|}}
% environnement proof (amsthm non chargé) utilisé par les modèles
\ifdefined\proof\else\newenvironment{proof}[1][Démonstration]{\par\noindent\textit{#1.}\ }{\qed\par}\fi
\usepackage{lastpage}
\usepackage{hyperref}
\hypersetup{hidelinks}

% ============================================================
% Palette « Pop » OnBuch+ — mêmes hex que la classe P (plus_ui.dart)
% ============================================================
\definecolor{poporange}{HTML}{F59321}
\definecolor{popdark}{HTML}{E07A0C}
\definecolor{popcream}{HTML}{FFF6E8}
\definecolor{popink}{HTML}{1C1714}
\definecolor{poppurple}{HTML}{7A5AE0}
\definecolor{popgreen}{HTML}{2E9E6A}
\definecolor{poppink}{HTML}{E0457B}
\definecolor{popblue}{HTML}{2F7DE1}
\definecolor{popgold}{HTML}{C98A12}
\definecolor{popmuted}{HTML}{8A7F76}
\definecolor{popline}{HTML}{EADFD2}
\definecolor{popgray}{HTML}{8A7F76}
\colorlet{popgrey}{popgray}
% Alias tolérants (le modèle nomme parfois les couleurs différemment)
\colorlet{popwhite}{white}
\colorlet{popblack}{popink}
\colorlet{popred}{poppink}
\colorlet{popyellow}{popgold}
% Teintes claires pour fonds de tableaux / zones
\colorlet{popblueL}{popblue!10!white}
\colorlet{poporangeL}{poporange!14!white}
\colorlet{popgreenL}{popgreen!12!white}
\colorlet{poppurpleL}{poppurple!12!white}
\colorlet{poppinkL}{poppink!12!white}
\colorlet{popgoldL}{popgold!14!white}

\pagecolor{popcream}

% ============================================================
% Typographie
% ============================================================
\defaultfontfeatures{Ligatures=TeX}
\setmainfont{PlusJakartaSans-Medium.ttf}[Path=../../../fonts/,BoldFont=PlusJakartaSans-Bold.ttf]
% Maths en sans-serif (Fira Math) avec chiffres et lettres droites de Plus
% Jakarta, pour que les formules se fondent dans le texte.
\usepackage[math-style=ISO,bold-style=ISO]{unicode-math}
\setmathfont{FiraMath-Regular.otf}[Path=../../../fonts/]
\setmathfont{PlusJakartaSans-Medium.ttf}[Path=../../../fonts/,range={up/num,up/{Latin,latin},\mathup}]
\usepackage{icomma} % virgule décimale sans espace en mode math
% Caractères mathématiques Unicode absents des polices de texte (Plus Jakarta,
% Archivo Black) : rendus par leur équivalent mathématique, en texte comme en
% formule — sinon ils s'affichent en carré vide (ex. « Arithmétique dans ℤ »).
\usepackage{newunicodechar}
\newunicodechar{ℝ}{\ensuremath{\mathbb{R}}}
\newunicodechar{ℤ}{\ensuremath{\mathbb{Z}}}
\newunicodechar{ℂ}{\ensuremath{\mathbb{C}}}
\newunicodechar{ℕ}{\ensuremath{\mathbb{N}}}
\newunicodechar{ℚ}{\ensuremath{\mathbb{Q}}}
\newunicodechar{𝔻}{\ensuremath{\mathbb{D}}}
\newunicodechar{α}{\ensuremath{\alpha}}
\newunicodechar{β}{\ensuremath{\beta}}
\newunicodechar{γ}{\ensuremath{\gamma}}
\newunicodechar{δ}{\ensuremath{\delta}}
\newunicodechar{ε}{\ensuremath{\varepsilon}}
\newunicodechar{θ}{\ensuremath{\theta}}
\newunicodechar{λ}{\ensuremath{\lambda}}
\newunicodechar{μ}{\ensuremath{\mu}}
\newunicodechar{σ}{\ensuremath{\sigma}}
\newunicodechar{τ}{\ensuremath{\tau}}
\newunicodechar{φ}{\ensuremath{\varphi}}
\newunicodechar{ψ}{\ensuremath{\psi}}
\newunicodechar{ω}{\ensuremath{\omega}}
\newunicodechar{Σ}{\ensuremath{\Sigma}}
% majuscules produites par \MakeUppercase dans les titres (α-aminés -> Α)
\newunicodechar{Α}{\ensuremath{\alpha}}
\newunicodechar{Β}{\ensuremath{\beta}}
\newunicodechar{Γ}{\ensuremath{\Gamma}}
\newunicodechar{Δ}{\ensuremath{\Delta}}
\newunicodechar{Ε}{\ensuremath{\varepsilon}}
\newunicodechar{Θ}{\ensuremath{\Theta}}
\newunicodechar{Λ}{\ensuremath{\Lambda}}
\newunicodechar{Μ}{\ensuremath{\mu}}
\newunicodechar{Τ}{\ensuremath{\tau}}
\newunicodechar{Φ}{\ensuremath{\Phi}}
\newunicodechar{Ψ}{\ensuremath{\Psi}}
\newunicodechar{Ω}{\ensuremath{\Omega}}
\newunicodechar{↦}{\ensuremath{\mapsto}}
\newunicodechar{∈}{\ensuremath{\in}}
\newunicodechar{∉}{\ensuremath{\notin}}
\newunicodechar{∧}{\ensuremath{\wedge}}
\newunicodechar{⇔}{\ensuremath{\Leftrightarrow}}
\newunicodechar{⟺}{\ensuremath{\Longleftrightarrow}}
\newunicodechar{⇒}{\ensuremath{\Rightarrow}}
\newunicodechar{⟹}{\ensuremath{\Longrightarrow}}
\newunicodechar{∘}{\ensuremath{\circ}}
\newunicodechar{∪}{\ensuremath{\cup}}
\newunicodechar{∩}{\ensuremath{\cap}}
\newunicodechar{≡}{\ensuremath{\equiv}}
\newunicodechar{⊂}{\ensuremath{\subset}}
\newunicodechar{⊥}{\ensuremath{\perp}}
\newunicodechar{∃}{\ensuremath{\exists}}
\newunicodechar{∀}{\ensuremath{\forall}}
\newunicodechar{∛}{\ensuremath{\sqrt[3]{\,}}}
\newunicodechar{∜}{\ensuremath{\sqrt[4]{\,}}}
\newunicodechar{⃗}{\ensuremath{{}^{\to}}}
\newunicodechar{ⁿ}{\ifmmode^{n}\else\textsuperscript{\ensuremath{n}}\fi}
\newunicodechar{ˣ}{\ifmmode^{x}\else\textsuperscript{\ensuremath{x}}\fi}
\newunicodechar{ᵇ}{\ifmmode^{b}\else\textsuperscript{\ensuremath{b}}\fi}
\newunicodechar{ʳ}{\ifmmode^{r}\else\textsuperscript{\ensuremath{r}}\fi}
\newunicodechar{ᵘ}{\ifmmode^{u}\else\textsuperscript{\ensuremath{u}}\fi}
\newunicodechar{ʸ}{\ifmmode^{y}\else\textsuperscript{\ensuremath{y}}\fi}
\newunicodechar{ᵃ}{\ifmmode^{a}\else\textsuperscript{\ensuremath{a}}\fi}
\newunicodechar{ᵉ}{\ifmmode^{e}\else\textsuperscript{\ensuremath{e}}\fi}
\newunicodechar{ᵏ}{\ifmmode^{k}\else\textsuperscript{\ensuremath{k}}\fi}
\newunicodechar{ᵖ}{\ifmmode^{p}\else\textsuperscript{\ensuremath{p}}\fi}
\newunicodechar{ᴺ}{\ifmmode^{N}\else\textsuperscript{\ensuremath{N}}\fi}
\newunicodechar{ᶜ}{\ifmmode^{c}\else\textsuperscript{\ensuremath{c}}\fi}
\newunicodechar{ʰ}{\ifmmode^{h}\else\textsuperscript{\ensuremath{h}}\fi}
\newunicodechar{ᵠ}{\ifmmode^{q}\else\textsuperscript{\ensuremath{q}}\fi}
\newunicodechar{ₙ}{\ifmmode_{n}\else\textsubscript{\ensuremath{n}}\fi}
\newunicodechar{ₐ}{\ifmmode_{a}\else\textsubscript{\ensuremath{a}}\fi}
\newunicodechar{ᵢ}{\ifmmode_{i}\else\textsubscript{\ensuremath{i}}\fi}
\newunicodechar{ᵣ}{\ifmmode_{r}\else\textsubscript{\ensuremath{r}}\fi}
\newunicodechar{ₖ}{\ifmmode_{k}\else\textsubscript{\ensuremath{k}}\fi}
\newunicodechar{ᵦ}{\ifmmode_{\beta}\else\textsubscript{\ensuremath{\beta}}\fi}
\newunicodechar{ₓ}{\ifmmode_{x}\else\textsubscript{\ensuremath{x}}\fi}
\newfontfamily\popdisplay{ArchivoBlack-Regular.ttf}[Path=../../../fonts/]
\newfontfamily\poplogo{SpaceGrotesk-Bold.ttf}[Path=../../../fonts/]
\newfontfamily\popsemi{PlusJakartaSans-SemiBold.ttf}[Path=../../../fonts/]
\newfontfamily\popxbold{PlusJakartaSans-ExtraBold.ttf}[Path=../../../fonts/]
\newfontfamily\popxboldls{PlusJakartaSans-ExtraBold.ttf}[Path=../../../fonts/,LetterSpace=3.0]

\setlist[enumerate]{leftmargin=1.7em,itemsep=5pt,topsep=4pt}
\setlist[itemize]{leftmargin=1.7em,itemsep=4pt,topsep=4pt,label={\color{poporange}\textbullet}}
\setlength{\parindent}{0pt}
\setlength{\parskip}{5pt}
\renewcommand{\arraystretch}{1.25}

% pgfplots : style Pop par défaut
\pgfplotsset{
  every axis/.append style={axis line style={popink,line width=1pt},tick style={popink},
    grid style={popline,line width=0.5pt},label style={font=\small},tick label style={font=\footnotesize}},
  popcurve/.style={poporange,line width=1.8pt,no markers,smooth},
}
% Même style disponible dans un \draw TikZ simple (hors pgfplots)
\tikzset{popcurve/.style={poporange,line width=1.8pt}}

% ============================================================
% Pill / squiggle
% ============================================================
\newcommand{\poppill}[3]{%
  \tikz[baseline=-0.55ex]\node[draw=popink,line width=1.2pt,rounded corners=2.6pt,%
    fill=#2,inner xsep=6.5pt,inner ysep=3pt]{%
    \popxboldls\fontsize{7}{7}\selectfont\color{#3}\MakeUppercase{#1}};%
}
\newcommand{\popsquiggle}[1]{%
  \tikz[baseline=-0.4ex]{
    \draw[#1,line width=2.6pt,line cap=round]
      plot [smooth] coordinates {(0,0.09) (0.34,0.01) (0.68,0.21) (1.02,0.01) (1.36,0.10)};
  }%
}
% Mot-clé surligné (marqueur orange sous le texte)
\newcommand{\cle}[1]{{\popxbold\color{popdark}#1}}

% ============================================================
% En-tête / pied
% ============================================================
\pagestyle{fancy}
\fancyhf{}
\renewcommand{\headrulewidth}{0pt}
\renewcommand{\footrulewidth}{0pt}
\lhead{\raisebox{-0.15ex}{\poplogo\fontsize{13}{13}\selectfont\color{popink}OnBuch\color{poporange}+}}
\rhead{\poppill{\DOCMATIERE\ \textbullet\ \DOCNIVEAU\ \textbullet\ Leçon \DOCLECON}{white}{popink}}
\lfoot{\popxbold\fontsize{7.6}{7.6}\selectfont\color{popmuted}\MakeUppercase{Cours signé OnBuch+}\ \textbullet\ {\popsemi\DOCTITRE}}
\rfoot{\tikz[baseline=-0.6ex]\node[rounded corners=8.5pt,fill=popink,minimum height=17pt,minimum width=17pt,inner xsep=5pt,inner sep=0]{\popxbold\fontsize{8}{8}\selectfont\color{popcream}\thepage\,/\,\pageref*{LastPage}};}

% ============================================================
% Titres
% ============================================================
\newcommand{\chaptitre}[2]{%
  \begin{tcolorbox}[enhanced,colback=poporange,colframe=popink,boxrule=2.6pt,arc=11pt,
      drop shadow=popink,left=15pt,right=15pt,top=12pt,bottom=11pt,before skip=3pt,after skip=18pt]
    \poppill{#2}{popink}{popcream}\par\vspace{9pt}
    {\raggedright\hyphenpenalty=10000\exhyphenpenalty=10000\popdisplay\fontsize{22}{24}\selectfont\color{popcream}\MakeUppercase{#1}\par}
  \end{tcolorbox}
}
% Section numérotée : pastille orange avec le numéro + titre Archivo
\newcounter{popsec}
\newcommand{\coursec}[1]{%
  \stepcounter{popsec}\setcounter{popsub}{0}%
  \par\vspace{16pt}\noindent
  \begin{minipage}{\linewidth}
  \tikz[baseline=-0.7ex]\node[draw=popink,line width=1.6pt,fill=poporange,rounded corners=4pt,minimum size=22pt,inner sep=0]{\popdisplay\fontsize{12}{12}\selectfont\color{popcream}\thepopsec};%
  \hspace{8pt}{\popdisplay\fontsize{15}{17}\selectfont\color{popink}\MakeUppercase{#1}}\par
  \vspace{4pt}\noindent{\color{poporange}\rule{36pt}{3.6pt}}
  \end{minipage}\par\nopagebreak\vspace{8pt}
}
\newcounter{popsub}
\newcommand{\courssub}[1]{%
  \stepcounter{popsub}%
  \par\vspace{9pt}\noindent{\popxbold\fontsize{11.5}{13}\selectfont\color{popink}{\color{poporange}\thepopsec.\thepopsub}\hspace{6pt}#1}\par\nopagebreak\vspace{4pt}
}

% ============================================================
% Boîtes pédagogiques — même recette : fond blanc, bordure épaisse colorée,
% ombre dure encre, badge plein. Titre optionnel [#1].
% ============================================================
\tcbset{popbase/.style={breakable,enhanced,colback=white,boxrule=2.3pt,arc=9pt,
  shadow={3pt}{-3pt}{0pt}{popink},left=14pt,right=14pt,top=11pt,bottom=11pt,before skip=11pt,after skip=15pt,
  fontupper=\popsemi\fontsize{10.6}{14.5}\selectfont\color{popink},
  fontlower=\popsemi\fontsize{10.6}{14.5}\selectfont\color{popink}}}
% Coche dessinée (le glyphe ✓ n'existe pas dans les polices du document)
\newcommand{\popcheck}[1]{\tikz[baseline=-0.5ex]\draw[#1,line width=1.6pt,line cap=round,line join=round](-0.13,0.0)--(-0.04,-0.09)--(0.13,0.1);}
\newcommand{\popboxhead}[3]{\poppill{#1}{#2}{white}\ifx\relax#3\relax\else\hspace{6pt}{\popxbold\fontsize{10.6}{12}\selectfont\color{popink}#3}\fi\par\vspace{7pt}}

\newtcolorbox{aretenir}[1][]{popbase,colframe=popgreen,before upper={\popboxhead{À retenir}{popgreen}{#1}}}
\newtcolorbox{exemplebox}[1][]{popbase,colframe=poporange,before upper={\popboxhead{Exemple}{poporange}{#1}}}
\newtcolorbox{definition}[1][]{popbase,colframe=popblue,before upper={\popboxhead{Définition}{popblue}{#1}}}
\newtcolorbox{propriete}[1][]{popbase,colframe=poppurple,before upper={\popboxhead{Propriété}{poppurple}{#1}}}
\newtcolorbox{methode}[1][]{popbase,colframe=popgold,before upper={\popboxhead{Méthode}{popgold}{#1}}}
\newtcolorbox{attention}[1][]{popbase,colframe=poppink,before upper={\popboxhead{Attention}{poppink}{#1}}}
\newtcolorbox{experience}[1][]{popbase,colframe=popblue,colback=popblueL,before upper={\popboxhead{Expérience}{popblue}{#1}}}
% « Le savais-tu ? » : carte violette pleine (contexte, culture, Cameroun)
\newtcolorbox{savaistu}[1][]{popbase,colback=poppurple,colframe=popink,
  fontupper=\popsemi\fontsize{10.4}{14.2}\selectfont\color{white},
  before upper={\poppill{Le savais-tu\,?}{white}{poppurple}\ifx\relax#1\relax\else\hspace{6pt}{\popxbold\color{white}#1}\fi\par\vspace{7pt}}}
% Exercice résolu : énoncé en haut, \tcblower puis la solution sur fond crème
\newcounter{popexo}\newcounter{popexr}
\newtcolorbox{exoresolu}[1][]{popbase,colframe=poporange,colbacklower=popcream,
  segmentation style={popink,line width=1.2pt,dashed},
  before upper={\stepcounter{popexr}\popboxhead{Exercice résolu \thepopexr}{poporange}{#1}},
  before lower={\poppill{Solution}{popink}{popcream}\par\vspace{6pt}}}
% Exercice (non résolu en place) — la correction est dans la partie Corrigés
\newtcolorbox{exercice}[1][]{popbase,colframe=popink,
  before upper={\stepcounter{popexo}\popboxhead{Exercice \thepopexo}{popink}{#1}}}
% Corrigé
\newtcolorbox{corrige}[1][]{popbase,colframe=popgreen,colback=popgreenL,
  before upper={\popboxhead{Corrigé}{popgreen}{#1}}}
% Objectifs / prérequis / bilan
\newtcolorbox{objectifs}{popbase,colframe=popink,colback=poporangeL,
  before upper={\popboxhead{Objectifs de la leçon}{popink}{}}}
\newtcolorbox{prerequis}{popbase,colframe=popink,colback=popblueL,
  before upper={\popboxhead{Prérequis}{popblue}{}}}
\newtcolorbox{bilan}{popbase,colframe=popink,colback=white,boxrule=2.8pt,
  before upper={\popboxhead{Fiche bilan}{poporange}{L'essentiel à connaître pour l'examen}}}
% Figure encadrée avec légende
\newtcolorbox{popfigure}[1][]{enhanced,breakable=false,colback=white,colframe=popline,boxrule=1.4pt,arc=8pt,
  left=10pt,right=10pt,top=10pt,bottom=8pt,before skip=10pt,after skip=12pt,halign=center,
  lower separated=false,halign lower=center,
  fontlower=\popsemi\fontsize{9}{11.5}\selectfont\color{popmuted},
  #1}
\newcommand{\legende}[1]{\par\vspace{6pt}{\popsemi\fontsize{9}{11.5}\selectfont\color{popmuted}#1}}

% ============================================================
% Signature OnBuch+
% ============================================================
\newcommand{\poplogoinline}[1]{{\poplogo\fontsize{#1}{#1}\selectfont\color{popink}OnBuch\color{poporange}+}}
\newcommand{\signatureonbuch}{%
  \begin{tcolorbox}[enhanced,colback=popink,colframe=popink,boxrule=2.2pt,arc=10pt,
      drop shadow=poporange,left=14pt,right=14pt,top=10pt,bottom=10pt,before skip=0pt,after skip=0pt]
    \tikz[baseline=-0.7ex]\node[circle,fill=popgreen,draw=popcream,line width=1.3pt,minimum size=22pt,inner sep=0]{\popcheck{white}};%
    \hspace{9pt}\begin{minipage}[c]{0.62\linewidth}
      {\popxbold\fontsize{12}{13}\selectfont\color{popcream}Signé\ \textbullet\ {\poplogo OnBuch\color{poporange}+}}\par\vspace{2pt}
      {\raggedright\popsemi\fontsize{8}{10.5}\selectfont\color{popline}\MakeUppercase{Équipe pédagogique OnBuch+}\\\MakeUppercase{Conforme au programme officiel MINESEC}\par}
    \end{minipage}\hfill
    {\poplogo\fontsize{22}{22}\selectfont\color{popcream}OnBuch\color{poporange}+}
  \end{tcolorbox}%
}

% ============================================================
% Page de garde
% #1 titre · #2 matière · #3 niveau · #4 module · #5 n° leçon · #6 durée ·
% #7 description / objectifs (texte ou itemize)
% ============================================================
\newcommand{\pagegarde}[7]{%
  \thispagestyle{empty}
  \newgeometry{a4paper,margin=1.7cm,top=1.6cm,bottom=1.4cm}%
  \begin{tikzpicture}[remember picture,overlay]
    \fill[poporange!18] ([xshift=-3.2cm,yshift=-3.6cm]current page.north east) circle (4.6cm);
    \fill[poppurple!14] ([xshift=2.4cm,yshift=7.6cm]current page.south west) circle (3.8cm);
  \end{tikzpicture}%
  {\poplogo\fontsize{30}{30}\selectfont\color{popink}OnBuch\color{poporange}+}\hfill
  \raisebox{0.6ex}{\poppill{Cours signé \textbullet\ Premium}{popink}{popcream}}\par
  \vspace{1pt}\popsquiggle{poppurple}\par
  \vspace{8pt}
  \begin{tcolorbox}[enhanced,colback=poporange,colframe=popink,boxrule=2.8pt,arc=13pt,
      drop shadow=popink,left=18pt,right=18pt,top=12pt,bottom=13pt,after skip=10pt]
    \poppill{#2 \textbullet\ #3}{popink}{popcream}\hspace{5pt}\poppill{#4}{white}{popink}\par\vspace{8pt}
    {\popdisplay\fontsize{14}{14}\selectfont\color{popink}LEÇON #5}\par\vspace{4pt}
    {\raggedright\hyphenpenalty=10000\exhyphenpenalty=10000\popdisplay\fontsize{25}{27}\selectfont\color{popcream}\MakeUppercase{#1}\par}
    \vspace{8pt}
    {\popxbold\fontsize{9.5}{9.5}\selectfont\color{popink}\MakeUppercase{Durée conseillée : #6}}
  \end{tcolorbox}
  \begin{tcolorbox}[enhanced,colback=white,colframe=popink,boxrule=2.2pt,arc=10pt,
      drop shadow=popink,left=16pt,right=16pt,top=10pt,bottom=10pt,after skip=8pt]
    \poppill{Dans cette leçon}{poppurple}{white}\par\vspace{5pt}
    \popsemi\fontsize{9.3}{12.6}\selectfont\color{popink}#7
  \end{tcolorbox}
  \noindent\begin{minipage}[t]{0.487\linewidth}
    \begin{tcolorbox}[enhanced,colback=popgreen,colframe=popink,boxrule=2.2pt,arc=10pt,
        drop shadow=popink,left=11pt,right=11pt,top=10pt,bottom=10pt,height=3.1cm,valign=center]
      \tcbox[colback=white,colframe=popink,boxrule=1.3pt,arc=5pt,boxsep=0pt,nobeforeafter,
        left=3pt,right=3pt,top=3pt,bottom=3pt,valign=center]{\qrcode[height=1.9cm]{https://play.google.com/store/apps/details?id=com.luvvix.onbuch&hl=fr}}%
      \hspace{8pt}\begin{minipage}[c]{0.5\linewidth}
        {\popdisplay\fontsize{11}{12.5}\selectfont\color{white}\MakeUppercase{Télécharge OnBuch}}\par\vspace{4pt}
        {\raggedright\popsemi\fontsize{8.4}{11}\selectfont\color{white}Gratuit \textbullet\ Google Play\\Tuteur IA, annales, concours, résultats\par}
      \end{minipage}
    \end{tcolorbox}
  \end{minipage}\hfill
  \begin{minipage}[t]{0.487\linewidth}
    \begin{tcolorbox}[enhanced,colback=poppurple,colframe=popink,boxrule=2.2pt,arc=10pt,
        drop shadow=popink,left=11pt,right=11pt,top=10pt,bottom=10pt,height=3.1cm,valign=center]
      \tikz[baseline=-0.7ex]\node[circle,fill=white,draw=popink,line width=1.3pt,minimum size=1.6cm,inner sep=0]{\popdisplay\fontsize{24}{24}\selectfont\color{poppurple}+};%
      \hspace{8pt}\begin{minipage}[c]{0.6\linewidth}
        {\popdisplay\fontsize{11}{12.5}\selectfont\color{white}\MakeUppercase{Passe à OnBuch+}}\par\vspace{4pt}
        {\raggedright\popsemi\fontsize{8.4}{11}\selectfont\color{white}Tous les cours signés, Tuteur IA illimité, fiches PDF — onglet OnBuch+ de l'app\par}
      \end{minipage}
    \end{tcolorbox}
  \end{minipage}
  \vspace{8pt}
  \popcontenu
  \vfill
  \signatureonbuch
  \restoregeometry
  \newpage
}

% Rangée « Ce cours contient » (page de garde)
\newcommand{\poptile}[3]{% #1 couleur #2 gros texte #3 légende
  \begin{tcolorbox}[enhanced,colback=#1,colframe=popink,boxrule=2pt,arc=9pt,shadow={2.5pt}{-2.5pt}{0pt}{popink},
      left=6pt,right=6pt,top=9pt,bottom=9pt,halign=center,valign=center,height=1.75cm,width=\linewidth,nobeforeafter]
    {\popdisplay\fontsize{12.5}{13}\selectfont\color{white}\MakeUppercase{#2}}\par\vspace{3pt}
    {\centering\popsemi\fontsize{7.8}{9.5}\selectfont\color{white}#3\par}
  \end{tcolorbox}}
\newcommand{\popcontenu}{%
  \par\noindent{\popxboldls\fontsize{7.5}{7.5}\selectfont\color{popmuted}CE COURS SIGNÉ CONTIENT}\par\vspace{7pt}
  \noindent\begin{minipage}[t]{0.235\linewidth}\poptile{popblue}{Cours}{complet, illustré, pas à pas}\end{minipage}\hfill
  \begin{minipage}[t]{0.235\linewidth}\poptile{poporange}{Méthodes}{et exercices résolus}\end{minipage}\hfill
  \begin{minipage}[t]{0.235\linewidth}\poptile{poppink}{Exercices}{type examen + corrigés}\end{minipage}\hfill
  \begin{minipage}[t]{0.235\linewidth}\poptile{popgreen}{Fiche bilan}{l'essentiel à retenir}\end{minipage}\par}

% Sommaire
\newtcolorbox{sommaire}{enhanced,colback=white,colframe=popink,boxrule=2.2pt,arc=10pt,
  drop shadow=popink,left=18pt,right=18pt,top=13pt,bottom=13pt,before skip=6pt,after skip=20pt,
  before upper={\poppill{Sommaire}{poppurple}{white}\par\vspace{9pt}\popsemi\fontsize{10.6}{14.5}\selectfont\color{popink}}}

% Signature de fin de cours — poussée tout en bas de la dernière page
\newcommand{\findecours}{%
  \par\vfill\nopagebreak
  \begin{center}\popsquiggle{poporange}\end{center}\vspace{4pt}
  \signatureonbuch
}
\AtBeginDocument{\def\llbracket{\mathopen{[\mkern-3.5mu[}}\def\rrbracket{\mathclose{]\mkern-3.5mu]}}} % intervalles d'entiers (glyphe ⟦ absent de la police)
% Glyphes absents de Fira Math : redéfinis à partir de glyphes disponibles
\AtBeginDocument{%
  \def\setminus{\mathbin{\backslash}}\let\smallsetminus\setminus
  \def\vdots{\mathord{\vbox{\baselineskip3.2pt\lineskiplimit0pt\kern4pt\hbox{.}\hbox{.}\hbox{.}}}}%
  \def\ddots{\mathinner{\mkern1mu\raise6.5pt\hbox{.}\mkern2mu\raise3.6pt\hbox{.}\mkern2mu\raise0.7pt\hbox{.}\mkern1mu}}%
  \def\triangle{\mathord{\scalebox{0.9}{$\symup{\Delta}$}}}%
  \def\nabla{\mathord{\raisebox{\depth}{\scalebox{1}[-1]{$\symup{\Delta}$}}}}%
  \def\checkmark{\popcheck{popdark}}%
  \def\star{\mathbin{\ast}}\def\bigstar{\mathord{\ast}}%
  \def\diamond{\mathbin{\scalebox{0.6}{$\Diamond$}}}%
  \def\bigcup{\mathop{\mathchoice{\raisebox{-0.3em}{\scalebox{1.7}{$\cup$}}}{\scalebox{1.25}{$\cup$}}{\cup}{\cup}}\displaylimits}%
  \def\bigcap{\mathop{\mathchoice{\raisebox{-0.3em}{\scalebox{1.7}{$\cap$}}}{\scalebox{1.25}{$\cap$}}{\cap}{\cap}}\displaylimits}%
  \def\lesseqqgtr{\mathrel{\lessgtr}}\def\bigoplus{\mathop{\oplus}\displaylimits}%
}
\providecommand{\ssi}{si et seulement si} % abréviation fréquente
\AtBeginDocument{\providecommand{\wideparen}{\overparen}} % arc de cercle

%% FIGFIX-BEGIN v3 — correctifs globaux des figures (docs/onbuch-generation/tools/figfix)
\usepackage{adjustbox}\usepackage{etoolbox}
% toute figure TikZ/pgfplots plus large que la ligne est réduite à la largeur disponible
\BeforeBeginEnvironment{tikzpicture}{\begin{adjustbox}{max width=\linewidth}}
\AfterEndEnvironment{tikzpicture}{\end{adjustbox}}
% légendes pgfplots lisibles par-dessus les courbes
\pgfplotsset{every axis/.append style={legend style={fill=white,fill opacity=0.92,text opacity=1,draw=black!15,inner sep=3pt}}}
% étiquettes d'arêtes (syntaxe edge["..."]) avec fond blanc
\tikzset{every edge quotes/.append style={fill=white,inner sep=1.2pt}}
% bulles de cartes mentales : texte à la ligne dans la bulle, bulle un peu plus grande
\tikzset{every concept/.append style={text width=2.0cm,align=center,minimum size=2.3cm,inner sep=2pt}}
%% FIGFIX-END
\begin{document}

\pagegarde{Statistiques}{Mathématiques}{3ème}{Mathématiques – classe de 3ème}{N14}{3 séances (fiche de progression harmonisée)}{Cette leçon poursuit l'étude des statistiques entamée en 4ème en traitant le cas des séries regroupées en classes. Les élèves apprennent à organiser des données nombreuses en classes, à déterminer la classe modale, à calculer fréquences et moyenne, et à construire des représentations graphiques adaptées (histogramme, polygone, diagrammes).
\vspace{4pt}
\begin{multicols}{2}\small
\begin{itemize}
\item À la fin de cette leçon, je sais regrouper une série statistique en classes d'amplitude donnée
\item À la fin de cette leçon, je sais construire un tableau d'effectifs et de fréquences d'une série regroupée en classes
\item À la fin de cette leçon, je sais déterminer la classe modale d'une série regroupée en classes
\item À la fin de cette leçon, je sais calculer la fréquence d'une classe et vérifier que la somme des fréquences vaut 1
\item À la fin de cette leçon, je sais calculer le centre d'une classe
\item À la fin de cette leçon, je sais calculer la moyenne d'une série regroupée en classes à l'aide des centres de classes
\end{itemize}\end{multicols}}

\begin{sommaire}
\begin{multicols}{2}
\begin{enumerate}
\item Regroupement en classes
\item Classe modale et fréquences
\item Moyenne d'une série regroupée en classes
\item Histogramme
\item Diagrammes circulaires et semi-circulaires
\item Synthèse et résolution de problèmes
\item Activité d'intégration
\item Méthodes et exercices résolus
\item Exercices
\item Corrigés des exercices
\item Fiche bilan
\end{enumerate}
\end{multicols}
\end{sommaire}

\begin{objectifs}
\begin{itemize}
\item À la fin de cette leçon, je sais regrouper une série statistique en classes d'amplitude donnée
\item À la fin de cette leçon, je sais construire un tableau d'effectifs et de fréquences d'une série regroupée en classes
\item À la fin de cette leçon, je sais déterminer la classe modale d'une série regroupée en classes
\item À la fin de cette leçon, je sais calculer la fréquence d'une classe et vérifier que la somme des fréquences vaut 1
\item À la fin de cette leçon, je sais calculer le centre d'une classe
\item À la fin de cette leçon, je sais calculer la moyenne d'une série regroupée en classes à l'aide des centres de classes
\item À la fin de cette leçon, je sais construire et interpréter un histogramme d'une série regroupée en classes
\item À la fin de cette leçon, je sais construire un diagramme circulaire ou semi-circulaire à partir des fréquences
\item À la fin de cette leçon, je sais rédiger et justifier une démarche statistique complète pour résoudre un problème de la vie courante
\end{itemize}
\end{objectifs}

\begin{prerequis}
\begin{itemize}
\item Effectif, effectif total et fréquence d'une valeur (4ème)
\item Moyenne d'une série statistique non regroupée (4ème)
\item Diagrammes en bâtons et diagrammes circulaires (4ème)
\item Calculs avec les fractions et les pourcentages
\item Lecture et construction de tableaux de données
\end{itemize}
\end{prerequis}

\newpage

\coursec{Regroupement en classes}

Imagine la scène : le proviseur du lycée de Bafoussam a relevé les tailles des 120 élèves de 3ème pour commander les uniformes scolaires. Il se retrouve avec une liste interminable de 120 nombres : 162, 158, 175, 160, 168, 155, 172… Face à cette « montagne » de données brutes, impossible de décider combien d'uniformes de chaque taille commander ! Comment regrouper ces données en classes pour y voir clair ? Quelle taille est la plus fréquente ? Quelle est la taille moyenne des élèves ? Comment présenter ces résultats au fournisseur sous forme de graphique lisible ? Cette section donne les premiers outils pour répondre : le regroupement en classes.

\courssub{Rappel : vocabulaire des séries statistiques}

Avant de regrouper, rappelons le vocabulaire de base. Une \textbf{série statistique} décrit un \textbf{caractère} (ce qu'on mesure ou observe) pour chaque \textbf{individu} d'une \textbf{population}.

\begin{definition}[Vocabulaire de base]
Soit une population étudiée.
\begin{itemize}
    \item \textbf{Individu} : chaque élément de la population (ici, un élève).
    \item \textbf{Caractère} : la propriété mesurée (ici, la taille en cm). Il peut être \emph{quantitatif} (nombre : taille, masse, âge) ou \emph{qualitatif} (catégorie : couleur, sexe, série).
    \item \textbf{Valeurs du caractère} : les différentes valeurs prises par le caractère.
    \item \textbf{Effectif} d'une valeur : nombre d'individus prenant cette valeur. Noté souvent $n_i$.
    \item \textbf{Effectif total} : somme de tous les effectifs. Noté $N = \sum n_i$.
\end{itemize}
\end{definition}

\begin{exemplebox}[Série brute : tailles de 10 élèves (extrait)]
\begin{center}
\begin{tabular}{|c|c|c|c|c|c|c|c|c|c|c|}\hline
Élève & 1 & 2 & 3 & 4 & 5 & 6 & 7 & 8 & 9 & 10 \\ \hline
Taille (cm) & 162 & 158 & 175 & 160 & 168 & 155 & 172 & 160 & 159 & 165 \\ \hline
\end{tabular}
\end{center}
Ici, la population est l'ensemble des 10 élèves, le caractère est la taille (quantitatif continu), les valeurs sont 155, 158, 159, 160, 162, 165, 168, 172, 175. L'effectif de la valeur 160 est 2. L'effectif total $N=10$.
\end{exemplebox}

\courssub{Pourquoi regrouper en classes ?}

Quand le caractère est \textbf{quantitatif continu} (comme la taille, la masse, le temps) ou prend \textbf{beaucoup de valeurs différentes} (ex : notes sur 20 de 200 élèves), le tableau « valeur / effectif » devient illisible et peu exploitable. On perd la vision d'ensemble.

\begin{aretenir}[Intérêt du regroupement]
Regrouper les valeurs en \textbf{classes} (intervalles) permet :
\begin{itemize}
    \item de condenser l'information (moins de lignes) ;
    \item de faire apparaître la \textbf{distribution} des données (où se concentrent-elles ?) ;
    \item de calculer une \textbf{moyenne approchée} et de tracer des \textbf{graphiques} (histogramme, diagramme circulaire).
\end{itemize}
\end{aretenir}

\courssub{Définition d'une classe et amplitude}

On découpe l'étendue des valeurs en intervalles contigus, sans trou ni chevauchement.

\begin{definition}[Classe d'une série statistique]
Une \textbf{classe} est un intervalle de la forme $[a ; b[$ qui regroupe toutes les valeurs du caractère $x$ telles que $a \le x < b$.
\begin{itemize}
    \item $a$ est la \textbf{borne inférieure} (incluse).
    \item $b$ est la \textbf{borne supérieure} (exclue).
    \item L'\textbf{amplitude} de la classe est la différence $b - a$.
\end{itemize}
\end{definition}

\begin{propriete}[Choix des classes]
Pour faciliter l'étude et le tracé de l'histogramme, on choisit généralement des \textbf{classes de même amplitude}. L'amplitude se choisit en fonction de l'étendue des données ($x_{\max} - x_{\min}$) et du nombre de classes souhaité (souvent entre 5 et 15 classes).
\end{propriete}

\courssub{Convention sur les bornes : la règle $[a ; b[$}

C'est un point crucial : \textbf{la borne inférieure est incluse, la borne supérieure est exclue}. Ainsi, une valeur exactement égale à une borne supérieure appartient à la classe \emph{suivante}, pas à la classe actuelle.

\begin{attention}[Piège mortel : le double comptage aux bornes]
\textbf{Erreur fréquente :} Compter une donnée située sur une borne dans les deux classes adjacentes.
\begin{itemize}
    \item Si une taille vaut exactement 160 cm, elle appartient à la classe $[160 ; 170[$ et \textbf{pas} à $[150 ; 160[$.
    \item Règle mnémotechnique : « \emph{La borne de gauche reste, la borne de droite part.} »
\end{itemize}
Vérifie toujours que la somme des effectifs des classes retrouve l'effectif total $N$.
\end{attention}

\begin{popfigure}
\begin{tikzpicture}[scale=0.8, >=Stealth]
% Axe principal
\draw[thick, ->] (0,0) -- (14,0) node[right] {Taille (cm)};
% Graduations principales
\foreach \x/\label in {0/150, 2/155, 4/160, 6/165, 8/170, 10/175, 12/180, 14/185} {
    \draw (\x, -0.15) -- (\x, 0.15) node[below=2pt] {\label};
}
% Classes : rectangles au-dessus de l'axe
% Classe 1 : [150; 160[
\draw[popblue, thick, fill=popblueL] (0, 0.5) rectangle (4, 1.2);
\node[popblue] at (2, 1.4) {\textbf{Classe 1 : $[150 ; 160[$}};
\node[popblue, font=\small] at (2, 0.85) {Amplitude = 10};
% Borne incluse 150
\fill[popblue] (0, 0.85) circle (2.5pt);
\node[popblue, left] at (0, 0.85) {\small $\bullet$ 150 \textbf{incluse}};
% Borne exclue 160
\draw[popblue, thick] (4, 0.5) -- (4, 1.2);
\draw[popblue, fill=white, thick] (4, 0.85) circle (2.5pt);
\node[popblue, right] at (4, 0.85) {\small 160 \textbf{exclue} $\circ$};
% Classe 2 : [160; 170[
\draw[poporange, thick, fill=poporangeL] (4, 0.5) rectangle (8, 1.2);
\node[poporange] at (6, 1.4) {\textbf{Classe 2 : $[160 ; 170[$}};
\node[poporange, font=\small] at (6, 0.85) {Amplitude = 10};
% Borne incluse 160
\fill[poporange] (4, 0.85) circle (2.5pt);
\node[poporange, left] at (4, 0.85) {\small $\bullet$ 160 \textbf{incluse}};
% Borne exclue 170
\draw[poporange, thick] (8, 0.5) -- (8, 1.2);
\draw[poporange, fill=white, thick] (8, 0.85) circle (2.5pt);
\node[poporange, right] at (8, 0.85) {\small 170 \textbf{exclue} $\circ$};
% Classe 3 : [170; 180[
\draw[popgreen, thick, fill=popgreenL] (8, 0.5) rectangle (12, 1.2);
\node[popgreen] at (10, 1.4) {\textbf{Classe 3 : $[170 ; 180[$}};
\node[popgreen, font=\small] at (10, 0.85) {Amplitude = 10};
% Flèches de continuité
\draw[popmuted, <->, thick] (0, -0.8) -- (4, -0.8) node[midway, below] {Amplitude 10};
\draw[popmuted, <->, thick] (4, -0.8) -- (8, -0.8) node[midway, below] {Amplitude 10};
\draw[popmuted, <->, thick] (8, -0.8) -- (12, -0.8) node[midway, below] {Amplitude 10};
% Légende convention
\node[align=center, font=\small, text=popdark] at (7, -1.5) {
    Convention : \textbf{borne inf. incluse ($\bullet$)}, \textbf{borne sup. exclue ($\circ$)}. \\
    Les classes se touchent sans se chevaucher : $160 \in [160;170[$, $160 \notin [150;160[$.
};
\end{tikzpicture}
\legende{Représentation des classes juxtaposées sur une demi-droite graduée. La convention $[a;b[$ évite tout chevauchement ni trou.}
\end{popfigure}

\courssub{Centre d'une classe}

Pour calculer une moyenne ou tracer un histogramme, on a besoin d'une valeur unique représentant toute la classe. On utilise le \textbf{centre} (ou milieu) de la classe.

\begin{definition}[Centre d'une classe]
Pour une classe $[a ; b[$, son \textbf{centre} (ou valeur représentative) est :
\[
c = \frac{a + b}{2}
\]
C'est le milieu de l'intervalle. On suppose que toutes les valeurs de la classe sont concentrées en ce point pour les calculs approchés.
\end{definition}

\begin{exemplebox}[Centres des classes de tailles]
\begin{itemize}
    \item Classe $[150 ; 160[$ : centre $c = \frac{150+160}{2} = 155$ cm.
    \item Classe $[160 ; 170[$ : centre $c = \frac{160+170}{2} = 165$ cm.
    \item Classe $[170 ; 180[$ : centre $c = \frac{170+180}{2} = 175$ cm.
\end{itemize}
\end{exemplebox}

\courssub{Construction complète du tableau de regroupement}

Voici la démarche pas à pas pour transformer une liste brute en tableau exploitable.

\begin{methode}[Regrouper une série statistique en classes]
\begin{enumerate}
    \item \textbf{Trier} les données (optionnel mais utile) et repérer $x_{\min}$ et $x_{\max}$.
    \item \textbf{Choisir l'amplitude} $h$ des classes (entier pratique : 1, 2, 5, 10, 20…). Calculer le nombre de classes $k \approx \frac{x_{\max}-x_{\min}}{h}$. Viser $5 \le k \le 15$.
    \item \textbf{Ecrire les bornes} des classes en commençant par un multiple de $h$ inférieur ou égal à $x_{\min}$. Respecter la convention $[a ; b[$.
    \item \textbf{Dénombrer} : pour chaque donnée, cocher la classe correspondante (une seule !). Compter les coches $\rightarrow$ effectif $n_i$ de la classe.
    \item \textbf{Vérifier} : $\sum n_i = N$ (effectif total).
    \item \textbf{Calculer les centres} $c_i = \frac{a_i+b_i}{2}$ pour chaque classe.
    \item \textbf{Compléter le tableau} : Classes | Effectifs $n_i$ | Centres $c_i$ | (optionnel : fréquences $f_i = n_i/N$).
\end{enumerate}
\end{methode}

\begin{exoresolu}[Regroupement des masses de 30 sacs de riz]
\textbf{Énoncé :} Un commerçant du marché central de Yaoundé pèse 30 sacs de riz (en kg) :
\[
52,\; 48,\; 55,\; 50,\; 47,\; 53,\; 51,\; 49,\; 54,\; 50,\;
46,\; 52,\; 56,\; 48,\; 50,\; 51,\; 49,\; 53,\; 50,\; 47,\;
54,\; 52,\; 48,\; 55,\; 50,\; 49,\; 51,\; 53,\; 46,\; 52.
\]
Regrouper ces masses en classes d'amplitude 5 kg.

\textbf{Solution détaillée :}
\begin{enumerate}
    \item $x_{\min} = 46$, $x_{\max} = 56$. Étendue $= 10$.
    \item Amplitude choisie $h = 5$ kg. Nombre de classes $\approx 10/5 = 2$ (un peu juste, on prend 3 classes pour mieux voir). On démarre à 45 (multiple de 5 $\le 46$).
    \item Classes : $[45 ; 50[$, $[50 ; 55[$, $[55 ; 60[$.
    \item Dénombrement (on coche) :
    \begin{itemize}
        \item $[45 ; 50[$ : 46, 48, 47, 49, 46, 48, 49, 47, 48, 49, 46 $\rightarrow$ \textbf{10} sacs. (Attention : les 50 vont dans la classe suivante !)
        \item $[50 ; 55[$ : 52, 50, 53, 51, 54, 50, 52, 51, 50, 53, 50, 52, 54, 52, 51, 53, 52 $\rightarrow$ \textbf{17} sacs.
        \item $[55 ; 60[$ : 55, 56, 55 $\rightarrow$ \textbf{3} sacs. (Note : 55 est $\ge 55$, donc dans $[55;60[$).
    \end{itemize}
    \item Vérification : $10 + 17 + 3 = 30$. \textbf{OK}.
    \item Centres :
    $c_1 = \frac{45+50}{2} = 47,5$ kg.
    $c_2 = \frac{50+55}{2} = 52,5$ kg.
    $c_3 = \frac{55+60}{2} = 57,5$ kg.
\end{enumerate}

\textbf{Tableau final :}
\begin{center}
\begin{tabular}{|c|c|c|}\hline
\textbf{Classes (kg)} & \textbf{Effectifs $n_i$} & \textbf{Centres $c_i$ (kg)} \\ \hline
$[45 ; 50[$ & 10 & 47,5 \\ \hline
$[50 ; 55[$ & 17 & 52,5 \\ \hline
$[55 ; 60[$ & 3 & 57,5 \\ \hline
\textbf{Total} & \textbf{30} & \\ \hline
\end{tabular}
\end{center}
La classe modale (plus grand effectif) est $[50 ; 55[$. La masse moyenne approchée sera calculée plus tard avec les centres.
\end{exoresolu}

\courssub{Illustration : Tailles des 40 élèves de 3ème A}

Appliquons la méthode sur l'exemple du proviseur (échantillon de 40 élèves). Données (cm) :
152, 165, 158, 172, 160, 155, 168, 175, 159, 162,
170, 157, 164, 178, 161, 156, 169, 173, 154, 166,
171, 153, 167, 176, 163, 151, 174, 179, 150, 160,
177, 158, 162, 170, 165, 159, 168, 172, 164, 161.

$x_{\min}=150$, $x_{\max}=179$. Amplitude $h=10$ cm. Classes : $[150;160[$, $[160;170[$, $[170;180[$, $[180;190[$.

\begin{center}
\begin{tabular}{|c|c|c|c|}\hline
\textbf{Classes (cm)} & \textbf{Effectifs $n_i$} & \textbf{Centres $c_i$ (cm)} & \textbf{Fréquences $f_i$} \\ \hline
$[150 ; 160[$ & 12 & 155 & $12/40 = 0,30$ \\ \hline
$[160 ; 170[$ & 16 & 165 & $16/40 = 0,40$ \\ \hline
$[170 ; 180[$ & 12 & 175 & $12/40 = 0,30$ \\ \hline
$[180 ; 190[$ & 0 & 185 & $0$ \\ \hline
\textbf{Total} & \textbf{40} & & \textbf{1} \\ \hline
\end{tabular}
\end{center}

\begin{popfigure}
\begin{tikzpicture}
\begin{axis}[
    ybar,
    bar width=12pt,
    width=\linewidth,
    height=8cm,
    xlabel={Taille (cm)},
    ylabel={Effectif},
    xtick={155,165,175,185},
    xticklabels={$[150;160[$,\quad$[160;170[$,\quad$[170;180[$,\quad$[180;190[$},
    ymin=0, ymax=18,
    enlarge x limits=0.15,
    enlarge y limits={upper=0.15},
    tick label style={font=\small},
    label style={font=\small},
    nodes near coords,
    nodes near coords align={vertical},
    every node near coord/.append style={font=\small, popdark},
    bar shift=0pt,
    fill=popblueL,
    draw=popblue,
]
\addplot coordinates {(155,12) (165,16) (175,12) (185,0)};
\end{axis}
\end{tikzpicture}
\legende{Histogramme des tailles (aperçu de la section 4). Les classes modales sont $[150;160[$ et $[170;180[$ (ex-æquo) ou $[160;170[$ selon l'effectif max. Ici $[160;170[$ est la classe modale unique.}
\end{popfigure}

\courssub{Le savais-tu ? : Le recensement au Cameroun}

\begin{savaistu}[Statistiques et BUCREP]
Au Cameroun, le \textbf{BUCREP} (Bureau Central des Recensements et des Études de Population) organise le Recensement Général de la Population et de l'Habitat (RGPH). Pour présenter la structure par âge de la population (pyramide des âges), les statisticiens ne donnent pas l'âge exact de chaque Camerounais (trop de valeurs !). Ils regroupent en \textbf{classes d'âge} : $[0;5[$, $[5;10[$, $[10;15[$, …, $[75;80[$, $[80;+\infty[$. L'amplitude est souvent de 5 ans pour les jeunes, parfois 10 ans pour les âges élevés. Cela permet de voir immédiatement le « poids » de la jeunesse (base large de la pyramide) et de planifier les écoles, les centres de santé, l'emploi. C'est exactement la méthode que tu viens d'apprendre, appliquée à 27 millions d'individus !
\end{savaistu}

\begin{exercice}[Entraînement : Regroupement des notes]
Un professeur a relevé les notes sur 20 de 25 élèves :
12, 8, 14, 10, 16, 9, 11, 13, 15, 7,
12, 10, 14, 18, 9, 11, 13, 6, 15, 12,
10, 14, 11, 9, 13.
\begin{enumerate}
    \item Regrouper ces notes en classes d'amplitude 2 en commençant par $[6;8[$.
    \item Compléter le tableau (classes, effectifs, centres, fréquences en \%). 
    \item Quelle est la classe modale ?
\end{enumerate}
\end{exercice}

\begin{corrige}[Exercice : Regroupement des notes]
\begin{enumerate}
    \item Classes d'amplitude 2 : $[6;8[$, $[8;10[$, $[10;12[$, $[12;14[$, $[14;16[$, $[16;18[$, $[18;20[$.
    \item Dénombrement rigoureux (respecter $[a;b[$) :
    \begin{center}
    \begin{tabular}{|c|c|c|c|}\hline
    \textbf{Classes} & \textbf{Effectifs $n_i$} & \textbf{Centres $c_i$} & \textbf{Fréquences $f_i(\%)$} \\ \hline
    $[6 ; 8[$ & 2 (6, 7) & 7 & $8\%$ \\ \hline
    $[8 ; 10[$ & 4 (8, 9, 9, 9) & 9 & $16\%$ \\ \hline
    $[10 ; 12[$ & 6 (10, 11, 10, 11, 10, 11) & 11 & $24\%$ \\ \hline
    $[12 ; 14[$ & 6 (12, 13, 12, 13, 12, 13) & 13 & $24\%$ \\ \hline
    $[14 ; 16[$ & 5 (14, 15, 14, 15, 14) & 15 & $20\%$ \\ \hline
    $[16 ; 18[$ & 1 (16) & 17 & $4\%$ \\ \hline
    $[18 ; 20[$ & 1 (18) & 19 & $4\%$ \\ \hline
    \textbf{Total} & \textbf{25} & & \textbf{100\%} \\ \hline
    \end{tabular}
    \end{center}
    \item Il y a deux classes modales (série bimodale) : $\mathbf{[10 ; 12[}$ et $\mathbf{[12 ; 14[}$ (effectif maximal 6).
\end{enumerate}
\end{corrige}

\coursec{Classe modale et fréquences}

\courssub{Classe modale}

\begin{definition}[Classe modale]
Dans une série statistique regroupée en classes, la \textbf{classe modale} est la classe qui possède le \textbf{plus grand effectif} (ou la plus grande fréquence).
\end{definition}

\begin{aretenir}
\begin{itemize}
    \item La classe modale indique l'intervalle où les données sont les plus concentrées.
    \item Il peut y avoir plusieurs classes modales (série bimodale, multimodale) si plusieurs classes partagent l'effectif maximal.
    \item Pour une série non regroupée, on parle de \textbf{valeur modale} (la valeur la plus fréquente).
\end{itemize}
\end{aretenir}

\courssub{Fréquence et fréquence en pourcentage}

\begin{definition}[Fréquence]
La \textbf{fréquence} d'une classe (ou d'une valeur) est le quotient de son effectif par l'effectif total :
\[
f_i = \frac{n_i}{N}
\]
La \textbf{fréquence en pourcentage} est $f_i \times 100$.
\end{definition}

\begin{propriete}[Propriétés des fréquences]
\begin{itemize}
    \item Une fréquence est un nombre compris entre 0 et 1 ($0 \le f_i \le 1$).
    \item La somme des fréquences de toutes les classes vaut 1 : $\sum f_i = 1$.
    \item La somme des fréquences en pourcentage vaut 100 \%. 
\end{itemize}
\end{propriete}

\begin{attention}[Fréquence vs Effectif]
Ne confonds pas \textbf{effectif} (nombre entier, comptage) et \textbf{fréquence} (quotient, proportion). L'histogramme (section 4) utilise les effectifs (ou densités), le diagramme circulaire (section 5) utilise les fréquences (angles proportionnels).
\end{attention}

\begin{exemplebox}[Fréquences pour les sacs de riz (Ex. précédent)]
$N=30$.
\begin{itemize}
    \item Classe $[45;50[$ : $n_1=10$, $f_1 = 10/30 \approx 0,333$ soit $33,3\%$.
    \item Classe $[50;55[$ : $n_2=17$, $f_2 = 17/30 \approx 0,567$ soit $56,7\%$.
    \item Classe $[55;60[$ : $n_3=3$, $f_3 = 3/30 = 0,10$ soit $10\%$.
\end{itemize}
Vérification : $0,333 + 0,567 + 0,10 = 1$.
\end{exemplebox}

\coursec{Moyenne d'une série regroupée en classes}

\courssub{Formule de la moyenne approchée}

Lorsque les données sont regroupées en classes, on ne connaît plus les valeurs exactes. On \textbf{approximé} la moyenne en supposant que toutes les valeurs d'une classe sont égales au centre de cette classe.

\begin{propriete}[Moyenne d'une série regroupée en classes]
Soit une série regroupée en $k$ classes d'effectifs $n_i$ et de centres $c_i$ ($i=1\dots k$). La \textbf{moyenne approchée} $\bar{x}$ est :
\[
\bar{x} = \frac{\sum_{i=1}^{k} n_i c_i}{\sum_{i=1}^{k} n_i} = \frac{\sum n_i c_i}{N}
\]
On peut aussi utiliser les fréquences : $\bar{x} = \sum f_i c_i$.
\end{propriete}

\begin{attention}[Moyenne approchée vs exacte]
Cette moyenne est une \textbf{estimation}. Elle est exacte seulement si toutes les valeurs de chaque classe sont effectivement égales au centre (ce qui est rare). Plus l'amplitude des classes est petite, plus l'approximation est bonne. Au Probatoire, on demande souvent : « Calculer la moyenne \textbf{approchée} ».
\end{attention}

\begin{methode}[Calculer la moyenne approchée]
\begin{enumerate}
    \item Vérifier que le tableau contient les colonnes : Classes, Effectifs $n_i$, Centres $c_i$.
    \item Ajouter une colonne $n_i \times c_i$ (produit effectif $\times$ centre).
    \item Calculer la somme $\sum n_i c_i$.
    \item Diviser par l'effectif total $N$.
    \item Conclure avec une phrase : « La moyenne approchée est ... » (avec unité).
\end{enumerate}
\end{methode}

\begin{exoresolu}[Moyenne approchée des masses de riz]
Reprenons le tableau des 30 sacs de riz.
\begin{center}
\begin{tabular}{|c|c|c|c|}\hline
\textbf{Classes (kg)} & \textbf{Effectifs $n_i$} & \textbf{Centres $c_i$} & \textbf{$n_i c_i$} \\ \hline
$[45 ; 50[$ & 10 & 47,5 & $10 \times 47,5 = 475$ \\ \hline
$[50 ; 55[$ & 17 & 52,5 & $17 \times 52,5 = 892,5$ \\ \hline
$[55 ; 60[$ & 3 & 57,5 & $3 \times 57,5 = 172,5$ \\ \hline
\textbf{Total} & \textbf{30} & & \textbf{1540} \\ \hline
\end{tabular}
\end{center}
\[
\bar{x} = \frac{1540}{30} \approx 51,33 \text{ kg}
\]
\textbf{Conclusion :} La masse moyenne approchée d'un sac de riz est de $51,3$ kg (arrondi au dixième).
\end{exoresolu}

\begin{exercice}[Moyenne des tailles]
Calculer la taille moyenne approchée des 40 élèves de 3ème A (tableau de la section 1).
\end{exercice}

\begin{corrige}[Exercice : Moyenne des tailles]
\begin{center}
\begin{tabular}{|c|c|c|c|}\hline
\textbf{Classes (cm)} & $n_i$ & $c_i$ & $n_i c_i$ \\ \hline
$[150 ; 160[$ & 12 & 155 & 1860 \\ \hline
$[160 ; 170[$ & 16 & 165 & 2640 \\ \hline
$[170 ; 180[$ & 12 & 175 & 2100 \\ \hline
$[180 ; 190[$ & 0 & 185 & 0 \\ \hline
\textbf{Total} & \textbf{40} & & \textbf{6600} \\ \hline
\end{tabular}
\end{center}
$\bar{x} = 6600 / 40 = 165$ cm.
La taille moyenne approchée est \textbf{165 cm}.
\end{corrige}

\coursec{Histogramme}

\courssub{Définition et construction}

L'histogramme est la représentation graphique par excellence d'une série quantitative continue regroupée en classes \textbf{d'amplitude constante}.

\begin{definition}[Histogramme]
C'est un diagramme en barres juxtaposées (sans espace) où :
\begin{itemize}
    \item L'axe des abscisses porte la variable (échelle numérique, bornes des classes).
    \item L'axe des ordonnées porte les \textbf{effectifs} (ou les fréquences).
    \item La \textbf{base} de chaque rectangle correspond exactement à l'intervalle de la classe $[a_i ; b_i[$.
    \item La \textbf{hauteur} du rectangle est proportionnelle à l'effectif $n_i$ (si amplitudes égales, hauteur $= n_i$).
    \item La \textbf{surface} du rectangle est proportionnelle à l'effectif (règle générale valable même si amplitudes différentes : surface $= n_i$).
\end{itemize}
\end{definition}

\begin{propriete}[Cas des amplitudes égales]
Si toutes les classes ont la même amplitude $h$, on peut prendre :
\[
\text{Hauteur du rectangle} = \text{Effectif } n_i \quad (\text{ou Fréquence } f_i)
\]
Les rectangles ont même largeur, leurs surfaces sont proportionnelles aux hauteurs.
\end{propriete}

\begin{methode}[Tracer un histogramme (amplitudes égales)]
\begin{enumerate}
    \item Choisir une échelle sur l'axe des abscisses pour placer les bornes des classes.
    \item Choisir une échelle sur l'axe des ordonnées pour les effectifs max.
    \item Pour chaque classe $[a_i ; b_i[$ d'effectif $n_i$ : tracer un rectangle de base $[a_i ; b_i[$ et de hauteur $n_i$.
    \item Les rectangles se touchent (pas d'espace). Légender les axes et titrer le graphique.
\end{enumerate}
\end{methode}

\begin{exemplebox}[Histogramme des tailles (40 élèves)]
\begin{popfigure}
\begin{tikzpicture}
\begin{axis}[
    ybar,
    bar width=18pt,
    width=\linewidth,
    height=9cm,
    xlabel={Taille (cm)},
    ylabel={Effectif},
    xtick={155,165,175},
    xticklabels={$[150;160[$,\quad$[160;170[$,\quad$[170;180[$},
    ymin=0, ymax=18,
    enlarge x limits=0.15,
    enlarge y limits={upper=0.15},
    tick label style={font=\small},
    label style={font=\small},
    nodes near coords,
    nodes near coords align={vertical},
    every node near coord/.append style={font=\small, popdark},
    fill=popblueL,
    draw=popblue,
]
\addplot coordinates {(155,12) (165,16) (175,12)};
\end{axis}
\end{tikzpicture}
\legende{Histogramme des tailles des 40 élèves. La classe modale $[160;170[$ se voit immédiatement (rectangle le plus haut).}
\end{popfigure}
\end{exemplebox}

\courssub{Polygone des fréquences}

\begin{definition}[Polygone des fréquences]
On marque le sommet milieu de la base supérieure de chaque rectangle de l'histogramme (coordonnées $(c_i ; n_i)$ ou $(c_i ; f_i)$) et on relie ces points par des segments. On ferme souvent le polygone en ajoutant des points fictifs de fréquence 0 avant la première et après la dernière classe.
\end{definition}

\begin{popfigure}
\begin{tikzpicture}
\begin{axis}[
    width=\linewidth,
    height=7cm,
    xlabel={Taille (cm)},
    ylabel={Effectif},
    xtick={145,155,165,175,185},
    xticklabels={$[140;150[$,\quad$[150;160[$,\quad$[160;170[$,\quad$[170;180[$,\quad$[180;190[$},
    ymin=0, ymax=18,
    enlarge x limits=0.1,
    enlarge y limits={upper=0.15},
    tick label style={font=\small},
    label style={font=\small},
]
% Histogramme en fond
\addplot[ybar, bar width=18pt, fill=popblueL, draw=popblue, opacity=0.3] coordinates {(155,12) (165,16) (175,12)};
% Polygone
\addplot[poporange, thick, mark=*, mark size=2pt] coordinates {(145,0) (155,12) (165,16) (175,12) (185,0)};
\end{axis}
\end{tikzpicture}
\legende{Polygone des fréquences (orange) superposé à l'histogramme (bleu transparent).}
\end{popfigure}

\courssub{Histogramme à amplitudes inégales (Culture)}

Si les classes n'ont pas la même amplitude, la hauteur n'est plus l'effectif. On utilise la \textbf{densité} (effectif divisé par amplitude) pour que la \textbf{surface} du rectangle représente l'effectif.
\[
\text{Hauteur (densité)} = \frac{n_i}{b_i - a_i}
\]
\textit{Ce cas est au programme de Première, mais il est bon de savoir qu'il existe. En 3ème, on travaille quasi exclusivement avec des classes de même amplitude.}

\coursec{Diagrammes circulaires et semi-circulaires}

\courssub{Diagramme circulaire (Camembert)}

Il représente la répartition des fréquences (parts d'un tout). Adapté aux caractères qualitatifs ou quantitatifs regroupés en peu de classes.

\begin{definition}[Diagramme circulaire]
Un cercle (360°) est divisé en secteurs angulaires. L'angle au centre du secteur $i$ est proportionnel à la fréquence $f_i$ :
\[
\alpha_i = 360^\circ \times f_i = 360^\circ \times \frac{n_i}{N}
\]
\end{definition}

\begin{methode}[Tracer un diagramme circulaire]
\begin{enumerate}
    \item Calculer l'angle de chaque secteur : $\alpha_i = \frac{n_i}{N} \times 360^\circ$.
    \item Vérifier : $\sum \alpha_i = 360^\circ$.
    \item Tracer un cercle au compas. Tracer un rayon horizontal (origine des angles).
    \item À l'aide du rapporteur, construire chaque angle successif dans le sens trigonométrique (anti-horaire).
    \item Colorier/hachurer les secteurs, ajouter une légende (classe, effectif, \%). 
\end{enumerate}
\end{methode}

\begin{exemplebox}[Diagramme circulaire : Masses des sacs de riz]
$N=30$.
\begin{itemize}
    \item $[45;50[$ : $n=10$, $\alpha = \frac{10}{30}\times 360 = 120^\circ$.
    \item $[50;55[$ : $n=17$, $\alpha = \frac{17}{30}\times 360 = 204^\circ$.
    \item $[55;60[$ : $n=3$, $\alpha = \frac{3}{30}\times 360 = 36^\circ$.
\end{itemize}
Vérification : $120 + 204 + 36 = 360^\circ$.
\begin{popfigure}
\begin{tikzpicture}
% Camembert
\fill[popblue] (0,0) -- (0.00:2.3) arc (0.00:120.00:2.3) -- cycle;
\node[white,font=\scriptsize\bfseries] at (60.00:1.55) {10\,\%};
\fill[popblue] (3.0,1.90) rectangle ++(0.28,0.28); \node[anchor=west,font=\scriptsize] at (3.35,2.04) {$[45;50[$ 33 (10\,\%)};
\fill[poporange] (0,0) -- (120.00:2.3) arc (120.00:324.00:2.3) -- cycle;
\node[white,font=\scriptsize\bfseries] at (222.00:1.55) {17\,\%};
\fill[poporange] (3.0,1.48) rectangle ++(0.28,0.28); \node[anchor=west,font=\scriptsize] at (3.35,1.62) {$[50;55[$ 56 (17\,\%)};
\fill[popgreen] (0,0) -- (324.00:2.3) arc (324.00:360.00:2.3) -- cycle;
\node[white,font=\scriptsize\bfseries] at (342.00:1.55) {3\,\%};
\fill[popgreen] (3.0,1.06) rectangle ++(0.28,0.28); \node[anchor=west,font=\scriptsize] at (3.35,1.20) {$[55;60[$ 10\% (3\,\%)};
\draw[white,thick] (0,0) circle (2.3);
\end{tikzpicture}
\legende{Diagramme circulaire des masses de riz. La part majoritaire (56,7\%) est la classe $[50;55[$.}
\end{popfigure}
\end{exemplebox}

\courssub{Diagramme semi-circulaire}

C'est un demi-cercle (180°). Utile pour comparer deux séries sur la même figure (ex : garçons / filles) ou gagner de la place.

\begin{definition}[Diagramme semi-circulaire]
Même principe, mais l'angle total est $180^\circ$.
\[
\alpha_i = 180^\circ \times f_i = 180^\circ \times \frac{n_i}{N}
\]
On trace un demi-cercle, le diamètre horizontal sert de base.
\end{definition}

\begin{exemplebox}[Comparaison Garçons / Filles (Semi-cercles dos à dos)]
On peut placer deux demi-cercles dos à dos (diamètres colinéaires) pour comparer la répartition des tailles des garçons (au-dessus) et des filles (en dessous).
\begin{popfigure}
\begin{tikzpicture}[scale=0.9]
% Garçons (haut)
\foreach \angle/\color/\label in {0/popblueL/$[150;160[$, 54/popblue/$[160;170[$, 126/popblue!70!black/$[170;180[$} {
    \pgfmathsetmacro{\nextangle}{\angle + (180*0.3)} % dummy
}
% Dessin manuel simplifié pour l'exemple
% Garçons : effectifs 5, 10, 5 (Total 20) -> angles 45, 90, 45
\fill[popblueL] (0,0) -- (2,0) arc (0:45:2) -- cycle;
\fill[popblue] (0,0) -- ++(45:2) arc (45:135:2) -- cycle;
\fill[popblue!70!black] (0,0) -- ++(135:2) arc (135:180:2) -- cycle;
\draw (0,0) -- (2,0) arc (0:180:2) -- cycle;
\draw (0,0) -- ++(45:2) (0,0) -- ++(135:2);
\node[above] at (0,2.3) {\textbf{Garçons (20)}};
% Filles (bas) : effectifs 7, 6, 7 (Total 20) -> angles 63, 54, 63
\begin{scope}[yshift=-0.5cm, rotate=180]
\fill[poppinkL] (0,0) -- (2,0) arc (0:63:2) -- cycle;
\fill[poppink] (0,0) -- ++(63:2) arc (63:117:2) -- cycle;
\fill[poppink!70!black] (0,0) -- ++(117:2) arc (117:180:2) -- cycle;
\draw (0,0) -- (2,0) arc (0:180:2) -- cycle;
\draw (0,0) -- ++(63:2) (0,0) -- ++(117:2);
\node[below] at (0,-2.8) {\textbf{Filles (20)}};
\end{scope}
\end{tikzpicture}
\legende{Diagrammes semi-circulaires dos à dos : comparaison garçons/filles.}
\end{popfigure}
\end{exemplebox}

\coursec{Synthèse et résolution de problèmes}

\courssub{Résumé des formules et repères}

\begin{aretenir}[Boîte à outils 3ème - Statistiques]
\begin{itemize}
    \item \textbf{Regroupement :} Classes $[a;b[$ même amplitude, borne inf. incluse, borne sup. exclue.
    \item \textbf{Centre :} $c = \frac{a+b}{2}$.
    \item \textbf{Effectif total :} $N = \sum n_i$.
    \item \textbf{Fréquence :} $f = \frac{n}{N}$ (somme = 1).
    \item \textbf{Classe modale :} Classe de plus grand effectif.
    \item \textbf{Moyenne approchée :} $\bar{x} = \frac{\sum n_i c_i}{N} = \sum f_i c_i$.
    \item \textbf{Histogramme :} Rectangles juxtaposés, base = classe, hauteur = effectif (amplitudes égales), surface $\propto$ effectif (général).
    \item \textbf{Diagramme circulaire :} Angle $\alpha = 360^\circ \times f$.
    \item \textbf{Diagramme semi-circulaire :} Angle $\alpha = 180^\circ \times f$.
\end{itemize}
\end{aretenir}

\courssub{Méthodologie : Résoudre un problème complet}

\begin{methode}[Problème statistique complet (Type BEPC)]
\begin{enumerate}
    \item \textbf{Lire et comprendre} : Identifier la population, le caractère, l'effectif total $N$. Noter $x_{\min}, x_{\max}$.
    \item \textbf{Regrouper} : Choisir amplitude $h$ (souvent donnée), construire les classes $[a;b[$, compter les effectifs $n_i$, vérifier $\sum n_i = N$. Calculer centres $c_i$ et fréquences $f_i$.
    \item \textbf{Analyser} : Identifier la classe modale. Calculer la moyenne approchée $\bar{x}$. Interpréter : « La valeur la plus fréquente se situe dans... », « En moyenne, le caractère vaut... ».
    \item \textbf{Représenter} : 
        \begin{itemize}
            \item Histogramme (si quantitatif continu, classes même amplitude).
            \item Diagramme circulaire / semi-circulaire (pour voir les parts, comparer).
        \end{itemize}
    \item \textbf{Conclure} : Répondre à la question posée (ex : « Quelle taille d'uniforme commander le plus ? » $\rightarrow$ Classe modale).
\end{enumerate}
\end{methode}

\begin{exoresolu}[Problème complet : Consommation d'électricité]
\textbf{Énoncé :} La consommation mensuelle (en kWh) de 50 abonnés ENEO dans un quartier de Douala est relevée :
\[
\begin{array}{cccccccccc}
85 & 120 & 95 & 110 & 130 & 90 & 105 & 125 & 100 & 115 \\
140 & 98 & 112 & 108 & 122 & 88 & 102 & 118 & 128 & 92 \\
135 & 103 & 117 & 123 & 97 & 113 & 107 & 127 & 93 & 108 \\
132 & 99 & 114 & 109 & 124 & 94 & 111 & 119 & 129 & 96 \\
138 & 101 & 116 & 110 & 121 & 91 & 106 & 126 & 104 & 131
\end{array}
\]
\begin{enumerate}
    \item Regrouper en classes d'amplitude 10 kWh en commençant par $[80;90[$. Compléter le tableau (effectifs, centres, fréquences \%). 
    \item Déterminer la classe modale. Interpréter.
    \item Calculer la consommation moyenne approchée.
    \item Tracer l'histogramme et le diagramme semi-circulaire.
\end{enumerate}

\textbf{Solution détaillée :}

\textbf{1. Regroupement}
$x_{\min}=85$, $x_{\max}=140$. Amplitude $h=10$. Classes : $[80;90[$, $[90;100[$, $[100;110[$, $[110;120[$, $[120;130[$, $[130;140[$, $[140;150[$.

Dénombrement (vérification somme = 50) :
\begin{center}
\begin{tabular}{|c|c|c|c|}\hline
\textbf{Classes (kWh)} & \textbf{Effectifs $n_i$} & \textbf{Centres $c_i$} & \textbf{Fréq. $f_i(\%)$} \\ \hline
$[80 ; 90[$ & 3 (85, 88, 88? non 88 une fois. 85, 88, 88? Liste: 85, 88. Attendez. 85, 88. Il n'y a que 85 et 88 dans les 80. Et 88? Non. 85, 88. Cherchons <90 : 85, 88. C'est 2. Recheck liste.) 
\end{tabular}
\end{center}
\emph{Comptage rigoureux des 50 valeurs fournies :}
Lignes :
1: 85, 120, 95, 110, 130, 90, 105, 125, 100, 115
2: 140, 98, 112, 108, 122, 88, 102, 118, 128, 92
3: 135, 103, 117, 123, 97, 113, 107, 127, 93, 108
4: 132, 99, 114, 109, 124, 94, 111, 119, 129, 96
5: 138, 101, 116, 110, 121, 91, 106, 126, 104, 131

Classes $[a;b[$ :
$[80;90[$ : 85, 88 $\rightarrow$ \textbf{2}
$[90;100[$ : 95, 90, 98, 92, 97, 93, 99, 94, 96, 91, 99? (99 ligne 4), 96. Liste: 95, 90, 98, 92, 97, 93, 99, 94, 96, 91. $\rightarrow$ \textbf{10}
$[100;110[$ : 105, 100, 108, 102, 103, 107, 108, 109, 101, 106, 104. $\rightarrow$ \textbf{11}
$[110;120[$ : 110, 115, 112, 118, 113, 117, 114, 111, 119, 110, 116, 110? (ligne 5), 111? (ligne 4). Liste: 110, 115, 112, 118, 113, 117, 114, 111, 119, 110, 116. $\rightarrow$ \textbf{11}
$[120;130[$ : 120, 125, 122, 128, 123, 127, 124, 129, 121, 126, 121? (non), 126. Liste: 120, 125, 122, 128, 123, 127, 124, 129, 121, 126. $\rightarrow$ \textbf{10}
$[130;140[$ : 130, 135, 132, 138, 131. $\rightarrow$ \textbf{5}
$[140;150[$ : 140. $\rightarrow$ \textbf{1}
Total : 2+10+11+11+10+5+1 = 50. OK.

Tableau :
\begin{center}
\begin{tabular}{|c|c|c|c|}\hline
\textbf{Classes (kWh)} & \textbf{$n_i$} & \textbf{$c_i$} & \textbf{$f_i(\%)$} \\ \hline
$[80 ; 90[$ & 2 & 85 & 4\% \\ \hline
$[90 ; 100[$ & 10 & 95 & 20\% \\ \hline
$[100 ; 110[$ & 11 & 105 & 22\% \\ \hline
$[110 ; 120[$ & 11 & 115 & 22\% \\ \hline
$[120 ; 130[$ & 10 & 125 & 20\% \\ \hline
$[130 ; 140[$ & 5 & 135 & 10\% \\ \hline
$[140 ; 150[$ & 1 & 145 & 2\% \\ \hline
\textbf{Total} & \textbf{50} & & \textbf{100\%} \\ \hline
\end{tabular}
\end{center}

\textbf{2. Classe modale}
Deux classes modales : $\mathbf{[100 ; 110[}$ et $\mathbf{[110 ; 120[}$ (effectif 11 chacune). La consommation la plus fréquente se situe entre 100 et 120 kWh.

\textbf{3. Moyenne approchée}
Calcul $\sum n_i c_i$ :
$2\times85 = 170$
$10\times95 = 950$
$11\times105 = 1155$
$11\times115 = 1265$
$10\times125 = 1250$
$5\times135 = 675$
$1\times145 = 145$
Somme $= 170+950+1155+1265+1250+675+145 = 5610$.
$\bar{x} = 5610 / 50 = 112,2$ kWh.
Consommation moyenne approchée : \textbf{112,2 kWh}.

\textbf{4. Représentations graphiques}
\begin{popfigure}
\begin{tikzpicture}
\begin{axis}[
    ybar,
    bar width=10pt,
    width=0.9\linewidth,
    height=7cm,
    xlabel={Conso (kWh)},
    ylabel={Effectif},
    xtick={85,95,105,115,125,135,145},
    xticklabels={$[80;90[$,\quad$[90;100[$,\quad$[100;110[$,\quad$[110;120[$,\quad$[120;130[$,\quad$[130;140[$,\quad$[140;150[$},
    ymin=0, ymax=13,
    enlarge x limits=0.08,
    enlarge y limits={upper=0.15},
    tick label style={font=\tiny, rotate=45, anchor=east},
    label style={font=\small},
    nodes near coords,
    nodes near coords align={vertical},
    every node near coord/.append style={font=\tiny, popdark},
    fill=popblueL, draw=popblue,
]
\addplot coordinates {(85,2) (95,10) (105,11) (115,11) (125,10) (135,5) (145,1)};
\end{axis}
\end{tikzpicture}
\legende{Histogramme de la consommation. Distribution symétrique autour de 110-115 kWh.}
\end{popfigure}

\begin{popfigure}
\begin{tikzpicture}[scale=0.9]
% Semi-cercle
\foreach \angle/\color/\label in {
    0/popblueL/2,
    7.2/popblue/10,
    46.8/poporangeL/11,
    86.4/poporange/11,
    126/popgreenL/10,
    162/popgreen/5,
    176.4/poppinkL/1} {
    \pgfmathsetmacro{\nextangle}{\angle + (180*\label/50)} % wrong logic
}
% Let's draw manually using pie chart logic for semi-circle
% Angles: 2->7.2°, 10->36°, 11->39.6°, 11->39.6°, 10->36°, 5->18°, 1->3.6°. Sum=180.
\begin{scope}
\fill[popblueL] (0,0) -- (3,0) arc (0:7.2:3) -- cycle;
\fill[popblue] (0,0) -- ++(7.2:3) arc (7.2:43.2:3) -- cycle;
\fill[poporangeL] (0,0) -- ++(43.2:3) arc (43.2:82.8:3) -- cycle;
\fill[poporange] (0,0) -- ++(82.8:3) arc (82.8:122.4:3) -- cycle;
\fill[popgreenL] (0,0) -- ++(122.4:3) arc (122.4:158.4:3) -- cycle;
\fill[popgreen] (0,0) -- ++(158.4:3) arc (158.4:176.4:3) -- cycle;
\fill[poppinkL] (0,0) -- ++(176.4:3) arc (176.4:180:3) -- cycle;
\draw (0,0) -- (3,0) arc (0:180:3) -- cycle;
\draw (0,0) -- ++(7.2:3) (0,0) -- ++(43.2:3) (0,0) -- ++(82.8:3) (0,0) -- ++(122.4:3) (0,0) -- ++(158.4:3) (0,0) -- ++(176.4:3);
\node[below] at (0,-0.5) {\textbf{Consommation (kWh) - Semi-cercle}};
\end{scope}
\end{tikzpicture}
\legende{Diagramme semi-circulaire. Les deux classes centrales ($[100;120[$) occupent 44\% de la surface.}
\end{popfigure}
\end{exoresolu}

\begin{exercice}[BEPC Type : Série complète]
Une enquête sur le temps de trajet (en minutes) domicile-école de 40 élèves donne le tableau suivant :
\begin{center}
\begin{tabular}{|c|c|}\hline
\textbf{Temps (min)} & \textbf{Effectif} \\ \hline
$[0 ; 10[$ & 5 \\ \hline
$[10 ; 20[$ & 12 \\ \hline
$[20 ; 30[$ & 15 \\ \hline
$[30 ; 40[$ & 6 \\ \hline
$[40 ; 50[$ & 2 \\ \hline
\end{tabular}
\end{center}
\begin{enumerate}
    \item Vérifier l'effectif total. Calculer les centres et les fréquences en \%.
    \item Déterminer la classe modale. 
    \item Calculer le temps de trajet moyen approché.
    \item Tracer l'histogramme et le polygone des fréquences sur le même repère.
    \item Tracer le diagramme semi-circulaire.
\end{enumerate}
\end{exercice}

\begin{corrige}[Exercice : Temps de trajet]
\begin{enumerate}
    \item $N = 5+12+15+6+2 = 40$.
    \begin{center}
    \begin{tabular}{|c|c|c|c|}\hline
    \textbf{Classes} & $n_i$ & $c_i$ & $f_i(\%)$ \\ \hline
    $[0;10[$ & 5 & 5 & 12,5\% \\ \hline
    $[10;20[$ & 12 & 15 & 30\% \\ \hline
    $[20;30[$ & 15 & 25 & 37,5\% \\ \hline
    $[30;40[$ & 6 & 35 & 15\% \\ \hline
    $[40;50[$ & 2 & 45 & 5\% \\ \hline
    \textbf{Total} & \textbf{40} & & \textbf{100\%} \\ \hline
    \end{tabular}
    \end{center}
    \item Classe modale : $\mathbf{[20 ; 30[}$ (effectif max 15). La plupart des élèves mettent 20 à 30 min.
    \item $\sum n_i c_i = 5\times5 + 12\times15 + 15\times25 + 6\times35 + 2\times45 = 25 + 180 + 375 + 210 + 90 = 880$.
    $\bar{x} = 880 / 40 = 22$ min. Temps moyen approché : \textbf{22 min}.
    \item \textbf{Histogramme + Polygone} :
    \begin{popfigure}
    \begin{tikzpicture}
    \begin{axis}[
        ybar, bar width=14pt, width=\linewidth, height=7cm,
        xlabel={Temps (min)}, ylabel={Effectif},
        xtick={5,15,25,35,45},
        xticklabels={$[0;10[$,$[10;20[$,$[20;30[$,$[30;40[$,$[40;50[$},
        ymin=0, ymax=17, enlarge x limits=0.1, enlarge y limits={upper=0.15},
        tick label style={font=\small}, label style={font=\small},
        fill=popblueL, draw=popblue,
    ]
    \addplot coordinates {(5,5) (15,12) (25,15) (35,6) (45,2)};
    \addplot[poporange, thick, mark=*, mark size=2pt] coordinates {( -5,0) (5,5) (15,12) (25,15) (35,6) (45,2) (55,0)};
    \end{axis}
    \end{tikzpicture}
    \legende{Histogramme (bleu) et polygone des fréquences (orange).}
    \end{popfigure}
    \item \textbf{Semi-cercle} : Angles = $180 \times f_i$.
    $5 \to 22,5^\circ$, $12 \to 54^\circ$, $15 \to 67,5^\circ$, $6 \to 27^\circ$, $2 \to 9^\circ$. Somme $= 180^\circ$.
    \begin{popfigure}
    \begin{tikzpicture}[scale=0.8]
    \fill[popblueL] (0,0) -- (3,0) arc (0:22.5:3) -- cycle;
    \fill[popblue] (0,0) -- ++(22.5:3) arc (22.5:76.5:3) -- cycle;
    \fill[poporangeL] (0,0) -- ++(76.5:3) arc (76.5:144:3) -- cycle;
    \fill[poporange] (0,0) -- ++(144:3) arc (144:171:3) -- cycle;
    \fill[popgreenL] (0,0) -- ++(171:3) arc (171:180:3) -- cycle;
    \draw (0,0) -- (3,0) arc (0:180:3) -- cycle;
    \draw (0,0) -- ++(22.5:3) (0,0) -- ++(76.5:3) (0,0) -- ++(144:3) (0,0) -- ++(171:3);
    \node[below] at (0,-0.5) {\textbf{Temps de trajet (min)}};
    \end{tikzpicture}
    \legende{Diagramme semi-circulaire du temps de trajet.}
    \end{popfigure}
\end{enumerate}
\end{corrige}

\begin{attention}[Check-list avant de rendre ta copie (BEPC)]
\begin{itemize}
    \item \textbf{Classes} : Bien écrites $[a;b[$ ? Amplitudes égales ? Bornes cohérentes avec $x_{\min}, x_{\max}$ ?
    \item \textbf{Effectifs} : Somme $= N$ ? Pas de double comptage aux bornes (ex: 20 va dans $[20;30[$) ?
    \item \textbf{Centres} : $(a+b)/2$ ? Unités notées ?
    \item \textbf{Moyenne} : Formule $\frac{\sum n_i c_i}{N}$ écrite ? Calcul détaillé ? Résultat arrondi raisonnable (souvent 0,1 ou entier) ? Unité ?
    \item \textbf{Graphiques} : Axes gradués, légendés, unités ? Histogramme : rectangles joints, bases sur bornes exactes. Circulaire : angles calculés, somme 360° (ou 180°), légende claire.
    \item \textbf{Phrases} : Réponses aux questions posées (classe modale, interprétation moyenne).
\end{itemize}
\end{attention}

\coursec{Classe modale et fréquences}

Dans la section précédente, nous avons appris à \textbf{regrouper} une série statistique brute en classes d'amplitude égale pour la rendre lisible. Nous avons construit un tableau avec des classes, leurs effectifs et parfois les valeurs caractéristiques (milieu de classe). Mais un tableau, même bien fait, reste une liste de nombres. Pour \emph{analyser} la série, il faut en extraire des indicateurs synthétiques : quelle est la classe qui revient le plus souvent ? Quelle part représente chaque classe dans l'ensemble ? C'est l'objet de cette section : la \cle{classe modale} et les \cle{fréquences}.

\courssub{Rappel : le mode d'une série non regroupée}

Avant de parler de classes, rappelons la notion de \cle{mode} pour une série de valeurs brutes (non regroupées).

\begin{definition}[Mode d'une série statistique]
Pour une série statistique non regroupée, le \textbf{mode} est la valeur qui possède le \textbf{plus grand effectif}. C'est la valeur la plus fréquente, celle que l'on observe le plus souvent.
\end{definition}

\begin{exemplebox}[Mode sur notes brutes]
Un élève de 3ème a obtenu les notes suivantes sur 20 aux devoirs de mathématiques :
\[ 12,\; 14,\; 12,\; 15,\; 12,\; 13,\; 14,\; 12,\; 16,\; 11 \]
L'effectif de la note 12 est 4. C'est le plus grand effectif. Donc le \textbf{mode est 12}.
\end{exemplebox}

\begin{attention}[Mode vs Moyenne]
Ne confonds pas le \textbf{mode} (valeur la plus fréquente) et la \textbf{moyenne} (valeur d'équilibre). Dans l'exemple ci-dessus, la moyenne est $\frac{12+14+12+15+12+13+14+12+16+11}{10} = 13,1$, alors que le mode est 12. Ils donnent des informations différentes.
\end{attention}

\courssub{La classe modale : la classe « championne »}

Lorsque la série est regroupée en classes, on ne connaît plus les valeurs exactes, seulement l'intervalle dans lequel elles se trouvent. On ne peut donc plus parler de \emph{valeur} modale, mais de \emph{classe} modale.

\begin{definition}[Classe modale d'une série regroupée en classes]
Soit une série statistique regroupée en $k$ classes $C_1, C_2, \dots, C_k$ d'effectifs respectifs $n_1, n_2, \dots, n_k$.
La \textbf{classe modale} est la classe (ou les classes) qui possède(nt) le \textbf{plus grand effectif}.
\end{definition}

\begin{propriete}[Série unimodale, bimodale, multimodale]
\begin{itemize}
    \item Si une \textbf{seule} classe a l'effectif maximal, la série est \textbf{unimodale}.
    \item Si \textbf{deux} classes ont le même effectif maximal (et strictement supérieur aux autres), la série est \textbf{bimodale}. Il y a deux classes modales.
    \item Si \textbf{plus de deux} classes partagent l'effectif maximal, la série est \textbf{multimodale}.
\end{itemize}
\end{propriete}

\begin{exemplebox}[Recherche de la classe modale]
Reprenons les notes de mathématiques d'une classe de 3ème (effectif total $N=40$) regroupées en classes d'amplitude 5 :

\begin{center}
\begin{tabular}{|c|c|c|}\hline
\textbf{Classes (notes)} & \textbf{Effectifs $n_i$} & \textbf{Observation} \\ \hline
$[0 ; 5[$ & 2 & \\ \hline
$[5 ; 10[$ & 8 & \\ \hline
$[10 ; 15[$ & \textbf{18} & $\leftarrow$ \textbf{Plus grand effectif} \\ \hline
$[15 ; 20]$ & 12 & \\ \hline
\textbf{Total} & \textbf{40} & \\ \hline
\end{tabular}
\end{center}

La classe $[10 ; 15[$ a l'effectif 18, qui est le maximum. C'est la \textbf{classe modale}. La série est unimodale.
\emph{Interprétation :} La majorité des élèves ont une note comprise entre 10 et 15 (10 inclus, 15 exclu).
\end{exemplebox}

\begin{attention}[Piège majeur : Classe modale $\neq$ Mode]
\emph{Après regroupement en classes, on ne connaît plus la valeur exacte la plus fréquente (le mode).}
Dans l'exemple ci-dessus, la classe modale est $[10 ; 15[$. Le mode \emph{aurait pu} être 12, ou 13, ou 14... On ne le sait plus. \textbf{On ne doit jamais écrire : « Le mode est $[10 ; 15[$ »}. On écrit : « La classe modale est $[10 ; 15[$ ».
\end{attention}

\begin{savaistu}[Pourquoi « modal » ?]
Le terme vient du latin \emph{modus} qui signifie « manière », « mesure », « mode ». En statistiques, le mode est la valeur « à la mode », celle qui est la plus à la mode (la plus populaire) dans la série. La classe modale est donc la classe « à la mode ».
\end{savaistu}

\courssub{La fréquence d'une classe : la part du tout}

L'effectif brut (ex: 18 élèves) dépend de la taille totale de la population (ex: 40 élèves). Pour comparer des séries de tailles différentes ou pour comprendre l'importance relative d'une classe, on utilise la \cle{fréquence}.

\begin{definition}[Fréquence d'une classe]
Soit une classe $C_i$ d'effectif $n_i$ dans une série de total $N = \sum n_i$.
La \textbf{fréquence} de la classe $C_i$, notée $f_i$, est le quotient :
\[ f_i = \frac{n_i}{N} \]
\end{definition}

\begin{propriete}[Écritures équivalentes d'une fréquence]
Une fréquence est un nombre compris entre 0 et 1. On peut l'exprimer de trois façons équivalentes :
\begin{itemize}
    \item Sous forme \textbf{fractionnaire} : $f_i = \frac{n_i}{N}$ (ex: $\frac{18}{40}$).
    \item Sous forme \textbf{décimale} : $f_i = 0,45$ (arrondi éventuel).
    \item Sous forme de \textbf{pourcentage} : $f_i = 45\%$ (on multiplie le décimal par 100).
\end{itemize}
\end{propriete}

\begin{propriete}[Somme des fréquences]
La somme des fréquences de toutes les classes est toujours égale à 1 (ou 100\%).
\[ \sum_{i=1}^{k} f_i = \sum_{i=1}^{k} \frac{n_i}{N} = \frac{1}{N}\sum_{i=1}^{k} n_i = \frac{N}{N} = 1 \]
\end{propriete}

\begin{exoresolu}[Calcul des fréquences — Notes de 3ème]
Reprenons le tableau précédent ($N=40$) et calculons les fréquences sous les trois formes.

\begin{center}
\begin{tabular}{|c|c|c|c|c|c|}\hline
\textbf{Classes} & \textbf{Effectifs $n_i$} & \textbf{Fréquence $f_i = \frac{n_i}{40}$} & \textbf{Décimal} & \textbf{Pourcentage} & \textbf{Fréq. cumulée} \\ \hline
$[0 ; 5[$ & 2 & $\frac{2}{40} = \frac{1}{20}$ & 0,05 & 5\% & 0,05 (5\%) \\ \hline
$[5 ; 10[$ & 8 & $\frac{8}{40} = \frac{1}{5}$ & 0,20 & 20\% & 0,25 (25\%) \\ \hline
$[10 ; 15[$ & 18 & $\frac{18}{40} = \frac{9}{20}$ & \textbf{0,45} & \textbf{45\%} & 0,70 (70\%) \\ \hline
$[15 ; 20]$ & 12 & $\frac{12}{40} = \frac{3}{10}$ & 0,30 & 30\% & 1,00 (100\%) \\ \hline
\textbf{Total} & \textbf{40} & \textbf{1} & \textbf{1,00} & \textbf{100\%} & \\ \hline
\end{tabular}
\end{center}

\textbf{Analyse :}
\begin{itemize}
    \item La classe modale $[10 ; 15[$ représente \textbf{45\%} de l'effectif total. Près d'un élève sur deux a une note dans cet intervalle.
    \item La somme des fréquences (dernière ligne) vaut bien 1 (ou 100\%).
    \item La colonne « Fréquence cumulée » (vue en détail après) permet de dire : « 70\% des élèves ont une note strictement inférieure à 15 ».
\end{itemize}
\end{exoresolu}

\courssub{Effectifs cumulés et fréquences cumulées croissantes}

Ces notions, bien que parfois implicites, sont indispensables pour lire les tableaux officiels (bulletins, recensements, résultats d'examens comme le BEPC) et pour construire l'histogramme cumulé ou la courbe des effectifs cumulés (polygone des fréquences cumulées).

\begin{definition}[Effectif cumulé croissant]
L'\textbf{effectif cumulé croissant} d'une classe $C_i$ est la somme des effectifs de toutes les classes \textbf{inférieures ou égales} à $C_i$ (en suivant l'ordre naturel des classes).
On le note souvent $N_i$.
\[ N_i = n_1 + n_2 + \dots + n_i \]
Le dernier effectif cumulé est l'effectif total $N$.
\end{definition}

\begin{definition}[Fréquence cumulée croissante]
La \textbf{fréquence cumulée croissante} d'une classe $C_i$ est la somme des fréquences de toutes les classes inférieures ou égales à $C_i$.
On la note souvent $F_i$.
\[ F_i = f_1 + f_2 + \dots + f_i = \frac{N_i}{N} \]
La dernière fréquence cumulée vaut 1 (ou 100\%).
\end{definition}

\begin{methode}[Compléter un tableau avec effectifs et fréquences cumulés]
\begin{enumerate}
    \item Vérifier que les classes sont ordonnées (de la plus petite borne à la plus grande).
    \item Calculer l'effectif total $N = \sum n_i$.
    \item Calculer la fréquence $f_i = n_i / N$ pour chaque classe (garder la fraction pour exactitude).
    \item \textbf{Effectif cumulé :} Première classe : $N_1 = n_1$. Classe suivante : $N_{i} = N_{i-1} + n_i$.
    \item \textbf{Fréquence cumulée :} $F_i = N_i / N$ (ou $F_{i} = F_{i-1} + f_i$).
    \item Vérifier : le dernier $N_k = N$ et le dernier $F_k = 1$.
\end{enumerate}
\end{methode}

\begin{exemplebox}[Lecture d'un bulletin officiel — Fréquences cumulées]
Au BEPC, les notes sont souvent publiées par tranches. Imaginons un extrait pour une académie ($N=10\,000$ candidats) :

\begin{center}
\begin{tabular}{|c|c|c|c|}\hline
\textbf{Tranche de notes} & \textbf{Effectifs} & \textbf{Fréquence} & \textbf{Fréq. cumulée} \\ \hline
$[0 ; 5[$ & 500 & 5\% & 5\% \\ \hline
$[5 ; 10[$ & 2\,000 & 20\% & 25\% \\ \hline
$[10 ; 12[$ & 3\,500 & 35\% & 60\% \\ \hline
$[12 ; 14[$ & 2\,500 & 25\% & 85\% \\ \hline
$[14 ; 20]$ & 1\,500 & 15\% & 100\% \\ \hline
\end{tabular}
\end{center}

\textbf{Interprétations concrètes :}
\begin{itemize}
    \item La classe modale est $[10 ; 12[$ (35\% des candidats).
    \item « 60\% des candidats ont obtenu une note \textbf{strictement inférieure à 12} » (fréquence cumulée de la classe $[10 ; 12[$).
    \item « 85\% des candidats ont une note \textbf{inférieure à 14} ».
    \item « 15\% des candidats ont une note \textbf{supérieure ou égale à 14} » (complément à 1 de 85\%).
\end{itemize}
\end{exemplebox}

\courssub{Représentation graphique : Diagramme en bâtons des effectifs}

Pour visualiser immédiatement la classe modale, le \textbf{diagramme en bâtons} (ou diagramme en barres) des effectifs est idéal. La hauteur de chaque barre est proportionnelle à l'effectif de la classe. La barre la plus haute correspond à la classe modale.

\begin{popfigure}
\begin{tikzpicture}
\begin{axis}[
    width=\linewidth,
    height=8cm,
    ybar,
    bar width=18pt,
    ymin=0, ymax=22,
    ylabel={Effectifs},
    xlabel={Classes de notes},
    xtick={1,2,3,4},
    xticklabels={$[0;5[$, $[5;10[$, $[10;15[$, $[15;20]$},
    xticklabel style={rotate=0, anchor=north, yshift=-2pt},
    ymajorgrids=true,
    grid style={dashed, gray!50},
    axis lines*=left,
    enlarge x limits=0.15,
    title style={align=center},
    title={Diagramme en bâtons des effectifs — Notes de 3ème},
]
\addplot[fill=popblueL, draw=popblue, thick] coordinates {(1,2)};
\addplot[fill=popblueL, draw=popblue, thick] coordinates {(2,8)};
\addplot[fill=poporange, draw=popdark, very thick] coordinates {(3,18)}; % Classe modale en orange
\addplot[fill=popblueL, draw=popblue, thick] coordinates {(4,12)};
% Étiquette pour la classe modale
\node[above, font=\bfseries\small, popdark] at (axis cs:3,18) {Classe modale : 18};
\end{axis}
\end{tikzpicture}
\legende{Diagramme en bâtons mettant en évidence la classe modale $[10;15[$ (barre orange, effectif 18).}
\end{popfigure}

\begin{attention}[Diagramme en bâtons vs Histogramme]
\begin{itemize}
    \item \textbf{Diagramme en bâtons (barres) :} Les barres ont \textbf{toutes la même largeur}. La hauteur = effectif (ou fréquence). Utilisé pour variables discrètes \textbf{ou} classes qualitatives/quantitatives regroupées quand on veut comparer les effectifs bruts visuellement. Les barres sont souvent séparées par un espace.
    \item \textbf{Histogramme :} La \textbf{largeur} des barres est proportionnelle à l'amplitude de la classe. L'\textbf{aire} de la barre est proportionnelle à l'effectif (ou fréquence). Indispensable si les classes n'ont \textbf{pas la même amplitude}. Nous verrons l'histogramme rigoureux dans la section 4.
\end{itemize}
Ici, comme les classes ont même amplitude (5), le diagramme en bâtons et l'histogramme ont le même aspect visuel, mais leur définition mathématique diffère.
\end{attention}

\courssub{Synthèse et exercice d'application}

\begin{aretenir}[Bilan : Classe modale et Fréquences]
\begin{itemize}
    \item \textbf{Classe modale} : Classe de plus grand effectif. Peut y en avoir plusieurs (bimodale).
    \item \textbf{Fréquence d'une classe} : $f = \frac{\text{effectif de la classe}}{\text{effectif total}}$. Valeur entre 0 et 1.
    \item \textbf{Écritures} : Fraction $\frac{n_i}{N}$, Décimal (ex: 0,45), Pourcentage (ex: 45\%).
    \item \textbf{Somme des fréquences} = 1 (ou 100\%).
    \item \textbf{Effectif cumulé croissant} : Somme des effectifs jusqu'à la classe. Dernier = $N$.
    \item \textbf{Fréquence cumulée croissante} : Somme des fréquences jusqu'à la classe. Dernière = 1.
    \item \textbf{Interprétation \%} : « $p\%$ de la population se trouve dans cette classe » ou « $p\%$ ont une valeur inférieure à la borne sup de cette classe » (fréquence cumulée).
\end{itemize}
\end{aretenir}

\begin{exercice}[Application : Hauteurs de plants de maïs]
Un agriculteur de la région de l'Ouest mesure la hauteur (en cm) de 50 plants de maïs après 3 mois. Il obtient le tableau suivant :

\begin{center}
\begin{tabular}{|c|c|}\hline
\textbf{Hauteur $h$ (en cm)} & \textbf{Effectifs} \\ \hline
$[50 ; 60[$ & 5 \\ \hline
$[60 ; 70[$ & 12 \\ \hline
$[70 ; 80[$ & 18 \\ \hline
$[80 ; 90[$ & 10 \\ \hline
$[90 ; 100]$ & 5 \\ \hline
\end{tabular}
\end{center}

\begin{enumerate}
    \item Vérifier l'effectif total $N$.
    \item Compléter le tableau en ajoutant les colonnes : Fréquence (fraction, décimal, \%), Effectif cumulé, Fréquence cumulée.
    \item Quelle est la classe modale ? Interpréter sa fréquence en pourcentage.
    \item Quelle proportion (en \%) de plants a une hauteur \textbf{inférieure à 80 cm} ?
    \item L'agriculteur veut sélectionner les 20\% des plants les plus grands pour la reproduction. À partir de quelle hauteur (borne de classe) doit-il couper ?
\end{enumerate}
\end{exercice}

\begin{corrige}[Exercice : Hauteurs de plants de maïs]
\begin{enumerate}
    \item $N = 5 + 12 + 18 + 10 + 5 = 50$ plants.
    \item Tableau complété :
    \begin{center}
    \begin{tabular}{|c|c|c|c|c|c|c|}\hline
    \textbf{Classes} & \textbf{Eff. $n_i$} & \textbf{Fréq. $f_i$ (frac)} & \textbf{Fréq. (déc)} & \textbf{Fréq. (\%)} & \textbf{Eff. cum. $N_i$} & \textbf{Fr. cum. $F_i$} \\ \hline
    $[50 ; 60[$ & 5 & $\frac{5}{50}=\frac{1}{10}$ & 0,10 & 10\% & 5 & 0,10 (10\%) \\ \hline
    $[60 ; 70[$ & 12 & $\frac{12}{50}=\frac{6}{25}$ & 0,24 & 24\% & 17 & 0,34 (34\%) \\ \hline
    $[70 ; 80[$ & 18 & $\frac{18}{50}=\frac{9}{25}$ & \textbf{0,36} & \textbf{36\%} & 35 & 0,70 (70\%) \\ \hline
    $[80 ; 90[$ & 10 & $\frac{10}{50}=\frac{1}{5}$ & 0,20 & 20\% & 45 & 0,90 (90\%) \\ \hline
    $[90 ; 100]$ & 5 & $\frac{5}{50}=\frac{1}{10}$ & 0,10 & 10\% & 50 & 1,00 (100\%) \\ \hline
    \textbf{Total} & \textbf{50} & \textbf{1} & \textbf{1,00} & \textbf{100\%} & & \\ \hline
    \end{tabular}
    \end{center}
    \item La classe modale est $\boldsymbol{[70 ; 80[}$ (effectif 18, fréquence 36\%). \emph{Interprétation :} Plus d'un plant sur trois (36\%) a une hauteur comprise entre 70 cm (inclus) et 80 cm (exclu). C'est la tranche de hauteur la plus représentée.
    \item « Hauteur inférieure à 80 cm » $\rightarrow$ Fréquence cumulée de la classe $[70 ; 80[$ : $\boldsymbol{70\%}$.
    \item Les 20\% les plus grands correspondent aux 20\% de la queue de la distribution (fréquences cumulées les plus hautes).
    $100\% - 20\% = 80\%$. On cherche la classe où la fréquence cumulée dépasse 80\% (ou atteint 80\%).
    Fréquence cumulée à la fin de $[70 ; 80[$ : 70\%.
    Fréquence cumulée à la fin de $[80 ; 90[$ : 90\%.
    Le seuil des 20\% les plus grands se situe \textbf{dans la classe $[80 ; 90[$}. L'agriculteur doit sélectionner les plants de hauteur \textbf{supérieure ou égale à 80 cm} (borne inférieure de cette classe). Il en aura 15 (10+5), soit 30\%, ce qui est le plus proche possible par classes entières des 20\% visés (10 plants).
\end{enumerate}
\end{corrige}

\begin{savaistu}[Statistiques et décisions agricoles au Cameroun]
L'exemple du maïs n'est pas anodin. Au Cameroun, l'IRAD (Institut de Recherche Agricole pour le Développement) et les services du MINADER utilisent exactement ces outils (classes modales, fréquences cumulées, percentiles) pour :
\begin{itemize}
    \item Définir les \textbf{variétés améliorées} (ex: maïs CMS 8704, CMS 9015) en sélectionnant les plants les plus productifs (queue de distribution).
    \item Suivre la \textbf{sécurité alimentaire} : proportion de ménages en insécurité (fréquence), classes modales de revenus.
    \item Calibrer les \textbf{engrais} et intrants selon la distribution des hauteurs ou rendements observés dans les champs-écoles.
\end{itemize}
La statistique, c'est l'outil pour transformer des observations en décisions éclairées.
\end{savaistu}

\coursec{Moyenne d'une série regroupée en classes}

Dans la section précédente, nous avons appris à repérer la \cle{classe modale} d'une série regroupée en classes et à calculer les \cle{fréquences}. La classe modale nous dit \emph{où se concentrent} les données. Mais elle ne répond pas à une question tout aussi naturelle : quelle est la valeur \emph{moyenne} de la série ? Quelle est, par exemple, la taille moyenne des 40 élèves de notre classe de 3ème ? C'est l'objet de cette section. Nous allons voir qu'une difficulté nouvelle apparaît : le regroupement en classes a fait disparaître les valeurs exactes. Il faudra donc ruser intelligemment.

\courssub{Le problème posé par le regroupement en classes}

Reprenons la situation de la section 1. On a mesuré les tailles des 40 élèves d'une classe de 3ème, puis on les a regroupées dans le tableau suivant :

\begin{center}
\begin{tabularx}{\linewidth}{|l|X|X|X|X|X|}
\hline
\rowcolor{popblueL} Classe (taille en cm) & $[150 ; 155[$ & $[155 ; 160[$ & $[160 ; 165[$ & $[165 ; 170[$ & $[170 ; 175[$ \\
\hline
Effectif $n_i$ & 3 & 7 & 12 & 11 & 7 \\
\hline
\end{tabularx}
\end{center}

Pour calculer une moyenne \emph{à la manière habituelle}, il faudrait additionner les 40 tailles exactes, puis diviser par 40. Or le tableau ne nous donne plus les tailles exactes : nous savons seulement que 12 élèves mesurent entre 160 cm et 165 cm, sans connaître la taille précise de chacun. Les valeurs exactes sont \textbf{perdues}.

\begin{attention}[Ce qui change par rapport aux séries non regroupées]
Pour une série donnée valeur par valeur, la moyenne est $\bar{x} = \dfrac{\text{somme des valeurs}}{\text{effectif total}}$, et le résultat est \emph{exact}. Après un regroupement en classes, cette formule n'est plus directement applicable : il faut d'abord \textbf{reconstruire une valeur représentative pour chaque classe}.
\end{attention}

\courssub{Le principe : remplacer chaque classe par son centre}

L'idée est simple et naturelle. Puisque nous ne connaissons pas la taille exacte de chacun des 12 élèves de la classe $[160 ; 165[$, nous décidons de les compter \emph{tous} comme s'ils mesuraient la valeur du milieu de la classe, c'est-à-dire $\dfrac{160+165}{2} = 162,5$ cm. Cette valeur s'appelle le \cle{centre de la classe}.

\begin{definition}[Centre d'une classe]
Le \cle{centre} de la classe $[a ; b[$ est le nombre
\[ c = \frac{a+b}{2}. \]
C'est le milieu de l'intervalle : il est à égale distance des deux bornes.
\end{definition}

Pourquoi ce choix est-il raisonnable ? Les élèves de la classe $[160 ; 165[$ sont répartis un peu partout dans l'intervalle : certains font 161 cm, d'autres 164 cm. En les remplaçant tous par 162,5 cm, on fait des petites erreurs, mais celles-ci \emph{se compensent} globalement : les élèves plus petits que le centre compensent les élèves plus grands. Le centre est donc le meilleur représentant possible de la classe.

\begin{popfigure}
\begin{center}
\begin{tikzpicture}[scale=1.1]
\draw[popline, thick] (0,0) -- (10.5,0);
\foreach \x/\lab in {2/{$a$}, 8/{$b$}} {
  \draw[popdark, thick] (\x,-0.15) -- (\x,0.15);
  \node[below, popdark] at (\x,-0.15) {\lab};
}
\draw[popblue, very thick] (2,0.35) -- (8,0.35);
\fill[popblue!25] (2,0) rectangle (8,0.35);
\draw[popblue, thick] (2,0) -- (2,0.35);
\draw[popblue, thick] (8,0) -- (8,0.35);
\fill[poporange] (5,0) circle (2.5pt);
\draw[poporange, thick, ->] (5,1.4) -- (5,0.15);
\node[above, poporange] at (5,1.4) {centre $c = \dfrac{a+b}{2}$};
\draw[<->, popmuted] (2,-0.9) -- (5,-0.9) node[midway, below] {\small même distance};
\draw[<->, popmuted] (5,-0.9) -- (8,-0.9) node[midway, below] {\small même distance};
\node[popblue, align=center] at (5,0.75) {\small toutes les données de la classe\\ \small sont remplacées par $c$};
\end{tikzpicture}
\end{center}
\legende{Principe de substitution : chaque donnée de la classe $[a ; b[$ est approchée par le centre $c$, milieu de l'intervalle.}
\end{popfigure}

\courssub{La formule de la moyenne}

\begin{propriete}[Moyenne d'une série regroupée en classes]
Soit une série statistique regroupée en $k$ classes. On note, pour chaque classe numéro $i$ :
\begin{itemize}
\item $c_i$ le \cle{centre} de la classe ;
\item $n_i$ l'\cle{effectif} de la classe ;
\item $N = n_1 + n_2 + \dots + n_k$ l'\cle{effectif total}.
\end{itemize}
La moyenne de la série est alors :
\[ \bar{x} = \frac{n_1 c_1 + n_2 c_2 + \dots + n_k c_k}{N} = \frac{\text{somme des } n_i \times c_i}{N}. \]
\textbf{Justification guidée.} On approche chacune des $n_i$ données de la classe numéro $i$ par le centre $c_i$. La somme des données de cette classe vaut donc environ $n_i \times c_i$. En additionnant les contributions de toutes les classes, on obtient une valeur approchée de la somme totale des données : $n_1 c_1 + n_2 c_2 + \dots + n_k c_k$. Il reste à diviser par l'effectif total $N$, comme dans toute moyenne.
\end{propriete}

\begin{attention}[La moyenne obtenue est une valeur approchée]
Comme on a remplacé chaque donnée réelle par le centre de sa classe, le résultat n'est pas la moyenne exacte des 40 tailles, mais une \textbf{valeur approchée}, en général très proche de la vraie moyenne. Dans la rédaction, on écrira donc : \og la moyenne est \emph{environ} égale à 164 cm \fg{} ou \og $\bar{x} \approx 164$ cm \fg. Il ne faut \textbf{jamais} écrire \og la moyenne exacte est 164 cm \fg{} : c'est une erreur de vocabulaire sanctionnée au BEPC.
\end{attention}

\courssub{Organisation du calcul en tableau}

Le calcul se mène proprement en ajoutant des colonnes au tableau des effectifs : une colonne pour les centres $c_i$, une colonne pour les produits $n_i \times c_i$, et une ligne de total.

\begin{methode}[Calculer la moyenne d'une série regroupée en classes en 4 étapes]
\begin{enumerate}
\item \textbf{Centres} : calculer le centre $c_i = \dfrac{a_i + b_i}{2}$ de chaque classe.
\item \textbf{Produits} : multiplier chaque centre par l'effectif de sa classe : $n_i \times c_i$.
\item \textbf{Somme} : additionner tous les produits $n_i \times c_i$ et vérifier l'effectif total $N$.
\item \textbf{Division} : diviser la somme obtenue par $N$ : $\bar{x} = \dfrac{\text{somme des } n_i c_i}{N}$.
\end{enumerate}
\end{methode}

\begin{exemplebox}[Taille moyenne des 40 élèves]
Appliquons la méthode à notre série des tailles. On complète le tableau :

\begin{center}
\begin{tabularx}{\linewidth}{|l|X|X|X|X|X|X|}
\hline
\rowcolor{popblueL} Classe (cm) & $[150 ; 155[$ & $[155 ; 160[$ & $[160 ; 165[$ & $[165 ; 170[$ & $[170 ; 175[$ & \textbf{Total} \\
\hline
Centre $c_i$ & $152,5$ & $157,5$ & $162,5$ & $167,5$ & $172,5$ & \\
\hline
Effectif $n_i$ & 3 & 7 & 12 & 11 & 7 & $N = 40$ \\
\hline
$n_i \times c_i$ & $457,5$ & $1102,5$ & $1950$ & $1842,5$ & $1207,5$ & $6560$ \\
\hline
\end{tabularx}
\end{center}

Détail des produits : $3 \times 152,5 = 457,5$ ; $7 \times 157,5 = 1102,5$ ; $12 \times 162,5 = 1950$ ; $11 \times 167,5 = 1842,5$ ; $7 \times 172,5 = 1207,5$. La somme vaut $457,5 + 1102,5 + 1950 + 1842,5 + 1207,5 = 6560$. D'où :
\[ \bar{x} = \frac{6560}{40} = 164 \text{ cm}. \]
\textbf{Interprétation} : la taille moyenne des élèves de cette classe est d'\emph{environ} 164 cm. Cela signifie que si les 40 élèves avaient tous la même taille tout en conservant la même somme totale, chacun mesurerait environ 164 cm.
\end{exemplebox}

\begin{attention}[Erreur fréquente : diviser par le nombre de classes]
L'erreur la plus répandue consiste à diviser la somme des produits par le \textbf{nombre de classes} (ici 5) au lieu de l'\textbf{effectif total} $N$ (ici 40). C'est faux : chaque produit $n_i \times c_i$ représente la contribution de $n_i$ élèves, donc on divise par le nombre total d'élèves. Autre erreur du même type : calculer la moyenne des centres $\dfrac{152,5 + 157,5 + 162,5 + 167,5 + 172,5}{5}$, qui oublie que les classes n'ont pas le même effectif. Une classe de 12 élèves doit \og peser \fg{} plus lourd qu'une classe de 3 élèves.
\end{attention}

\courssub{Moyenne et classe modale : deux informations différentes}

Il est important de ne pas confondre ces deux indicateurs :
\begin{itemize}
\item la \cle{classe modale} indique \emph{où se trouvent le plus grand nombre de données} : ici $[160 ; 165[$, avec 12 élèves ;
\item la \cle{moyenne} indique la \emph{valeur d'équilibre} de la série : ici environ 164 cm.
\end{itemize}
Ces deux nombres ne coïncident pas forcément. Ici, la moyenne (164 cm) appartient à la classe modale $[160 ; 165[$, mais ce n'est pas une règle générale : si les grandes tailles avaient été plus nombreuses, la moyenne aurait été \og tirée \fg{} vers le haut, hors de la classe modale. La moyenne est \emph{sensible} à toutes les valeurs, tandis que la classe modale ne dépend que de l'effectif le plus grand. Les deux indicateurs se complètent pour décrire la série.

\begin{exoresolu}[Note moyenne à un devoir de mathématiques]
Lors d'un devoir de mathématiques dans un collège de Yaoundé, les notes des 50 élèves de 3ème ont été regroupées ainsi : classe $[0 ; 5[$ : 4 élèves ; classe $[5 ; 10[$ : 14 élèves ; classe $[10 ; 15[$ : 26 élèves ; classe $[15 ; 20]$ : 6 élèves. Calculer la note moyenne de la classe à ce devoir, puis interpréter le résultat.
\tcblower
\textbf{Étape 1 : les centres.} $c_1 = \dfrac{0+5}{2} = 2,5$ ; $c_2 = \dfrac{5+10}{2} = 7,5$ ; $c_3 = \dfrac{10+15}{2} = 12,5$ ; $c_4 = \dfrac{15+20}{2} = 17,5$.

\textbf{Étapes 2 et 3 : produits et sommes} (on organise en tableau) :
\begin{center}
\begin{tabularx}{\linewidth}{|l|X|X|X|X|X|}
\hline
\rowcolor{popblueL} Classe & $[0 ; 5[$ & $[5 ; 10[$ & $[10 ; 15[$ & $[15 ; 20]$ & \textbf{Total} \\
\hline
$c_i$ & $2,5$ & $7,5$ & $12,5$ & $17,5$ & \\
\hline
$n_i$ & 4 & 14 & 26 & 6 & $N = 50$ \\
\hline
$n_i \times c_i$ & $10$ & $105$ & $325$ & $105$ & $545$ \\
\hline
\end{tabularx}
\end{center}

\textbf{Étape 4 : division.} $\bar{x} = \dfrac{545}{50} = 10,9$.

\textbf{Interprétation} : la note moyenne de la classe est d'\emph{environ} 10,9 sur 20. La classe modale est $[10 ; 15[$ (26 élèves) : la moyenne 10,9 est proche de cette classe, ce qui confirme que le devoir a été plutôt réussi, avec une majorité d'élèves autour de la moyenne.
\end{exoresolu}

\begin{savaistu}[Un premier pas vers l'espérance mathématique]
La formule $\bar{x} = \dfrac{\text{somme des } n_i c_i}{N}$ peut se réécrire avec les fréquences : $\bar{x} = f_1 c_1 + f_2 c_2 + \dots + f_k c_k$, où $f_i = \dfrac{n_i}{N}$. On multiplie chaque valeur par son \og poids \fg{} (sa fréquence), puis on additionne. Cette idée de \emph{moyenne pondérée} est très puissante : tu la retrouveras dans les classes supérieures, où elle deviendra la notion d'\cle{espérance mathématique}, utilisée en probabilités pour prévoir, par exemple, le gain moyen à un jeu ou le chiffre d'affaires attendu d'un commerçant. Ce que tu apprends ici en 3ème est donc la première marche d'un outil fondamental des mathématiques et de l'économie.
\end{savaistu}

\begin{exercice}[À toi de jouer]
Un opérateur téléphonique a relevé la durée des appels passés un soir par 60 abonnés de Douala : durée $[0 ; 2[$ min : 6 abonnés ; $[2 ; 4[$ : 18 abonnés ; $[4 ; 6[$ : 24 abonnés ; $[6 ; 8[$ : 12 abonnés. Calculer la durée moyenne d'un appel, puis comparer avec la classe modale.
\end{exercice}

\begin{corrige}[Exercice]
Centres : $c_1 = 1$, $c_2 = 3$, $c_3 = 5$, $c_4 = 7$. Produits : $6 \times 1 = 6$ ; $18 \times 3 = 54$ ; $24 \times 5 = 120$ ; $12 \times 7 = 84$. Somme : $6 + 54 + 120 + 84 = 264$. Effectif total : $N = 6 + 18 + 24 + 12 = 60$. Moyenne : $\bar{x} = \dfrac{264}{60} = 4,4$ min. La durée moyenne d'un appel est d'environ 4,4 minutes. La classe modale est $[4 ; 6[$ (24 abonnés) ; la moyenne 4,4 min appartient à cette classe, ce qui montre une série assez équilibrée autour de sa valeur centrale.
\end{corrige}

\coursec{Histogramme}

Après avoir appris à calculer la moyenne d'une série statistique regroupée en classes, nous allons maintenant voir comment la \emph{représenter graphiquement}. Un tableau de nombres est précis, mais une figure parle souvent plus vite à l'intuition : on voit d'un coup d'œil où se concentre la population, quelle classe domine, comment les effectifs décroissent. C'est tout l'intérêt de l'histogramme.

\courssub{Définition et intuition : des rectangles pour voir la répartition}

Imagine que tu mesures la taille de tous les élèves de ton établissement pour commander des tenues de sport. Tu as regroupé les tailles en classes (par exemple $[1,40 ; 1,50[$, $[1,50 ; 1,60[$, etc.) et tu as compté l'effectif de chaque classe. Si tu traces un rectangle pour chaque classe, dont la \textbf{base} est l'intervalle de la classe et dont la \textbf{hauteur} est proportionnelle à l'effectif, tu obtiens une figure où \emph{l'aire de chaque rectangle est proportionnelle à l'effectif de la classe}. C'est exactement un histogramme.

\begin{definition}[Histogramme d'une série regroupée en classes]
Soit une série statistique quantitative continue regroupée en classes d'amplitude $a_i$ (longueur de l'intervalle) et d'effectif $n_i$ (ou de fréquence $f_i$).
L'\textbf{histogramme} est la figure formée par des rectangles \textbf{juxtaposés} (sans espace entre eux) tels que :
\begin{itemize}
    \item La base de chaque rectangle est l'intervalle de la classe (sur l'axe des abscisses).
    \item La hauteur de chaque rectangle est proportionnelle à l'effectif $n_i$ (ou à la fréquence $f_i$) \textbf{si toutes les classes ont la même amplitude}.
    \item Dans le cas général (classes d'amplitudes différentes), l'\textbf{aire} de chaque rectangle est proportionnelle à l'effectif (ou à la fréquence) de la classe. La hauteur est alors la \textbf{densité} : $h_i = \frac{n_i}{a_i}$ (ou $\frac{f_i}{a_i}$).
\end{itemize}
\end{definition}

\begin{aretenir}[Cas particulier du programme de 3ème]
En classe de 3ème, nous nous concentrons sur le cas le plus courant : \textbf{les classes ont toutes la même amplitude}.
Dans ce cas, la hauteur des rectangles est simplement proportionnelle aux effectifs (ou aux fréquences). On peut donc prendre directement l'effectif (ou la fréquence) comme hauteur, à condition de choisir une échelle commode sur l'axe des ordonnées.
\end{aretenir}

\courssub{Construction pas à pas : la méthode rigoureuse}

Construire un histogramme, ce n'est pas juste « tracer des barres ». C'est une démarche géométrique précise qui garantit que la figure respecte les données.

\begin{methode}[Construire un histogramme à classes de même amplitude]
\textbf{Étape 1 : Préparer les axes.}
\begin{itemize}
    \item Axe horizontal (abscisses) : gradué avec les \textbf{bornes des classes} (valeurs extrêmes de chaque intervalle). L'origine n'est pas nécessairement à 0 ; on peut faire une coupure (zigzag) si la première borne est loin de 0.
    \item Axe vertical (ordonnées) : gradué pour représenter les \textbf{effectifs} (ou les fréquences). Choisis une unité de longueur (ex: 1 cm pour 5 élèves) pour que le plus grand rectangle tienne sur la feuille.
\end{itemize}

\textbf{Étape 2 : Tracer les rectangles.}
Pour chaque classe $[a ; b[$ d'effectif $n$ :
\begin{itemize}
    \item Repère $a$ et $b$ sur l'axe des abscisses.
    \item Trace un rectangle de base $[a ; b]$ et de hauteur $n$ (selon l'échelle choisie).
    \item Les rectangles se touchent : le côté droit du rectangle de la classe $i$ coïncide avec le côté gauche du rectangle de la classe $i+1$. \textbf{Pas d'espace vide !}
\end{itemize}

\textbf{Étape 3 : Légender et titrer.}
Indique le titre de l'histogramme (ex: « Répartition des tailles des élèves de 3ème »), les unités sur les axes (ex: « Taille en m », « Effectif »), et l'échelle si nécessaire.
\end{methode}

\begin{attention}[Piège n°1 : Les espaces entre les rectangles]
\emph{Erreur très fréquente :} Dessiner des rectangles séparés par des espaces, comme un diagramme en bâtons.
\emph{Rappel :} L'histogramme représente une variable \textbf{continue} (tailles, masses, temps, distances...). Les classes se suivent sans trou : la fin d'une classe est le début de la suivante. Les rectangles doivent être \textbf{juxtaposés} (collés les uns aux autres). Laisser un espace suggère à tort qu'il y a des valeurs « interdites » entre les classes.
\end{attention}

\begin{attention}[Piège n°2 : Confondre histogramme et diagramme en bâtons]
\begin{itemize}
    \item \textbf{Histogramme :} Variable \textbf{continue} (ou discrète regroupée en classes). Rectangles \textbf{juxtapossés}. La base a un sens (intervalle de valeurs). L'aire compte.
    \item \textbf{Diagramme en bâtons :} Variable \textbf{discrète} non groupée (ex: nombre d'enfants par famille, notes sur 20). Bâtons \textbf{séparés} par des espaces. La base est une étiquette (la valeur exacte), pas un intervalle. Seule la hauteur compte.
\end{itemize}
\end{attention}

\courssub{Exemple chiffré : distances parcourues pour venir au lycée}

Prenons une situation concrète. Dans un lycée de Yaoundé, on a interrogé 50 élèves sur la distance (en km) qu'ils parcourent chaque matin pour venir en cours. Voici le tableau de regroupement en classes d'amplitude 2 km.

\begin{exemplebox}[Distances domicile–lycée (en km)]
\begin{center}
\begin{tabular}{|c|c|c|c|c|c|c|}\hline
\textbf{Classes (km)} & $[0 ; 2[$ & $[2 ; 4[$ & $[4 ; 6[$ & $[6 ; 8[$ & $[8 ; 10[$ & \textbf{Total} \\ \hline
\textbf{Effectifs} & 8 & 15 & 12 & 10 & 5 & 50 \\ \hline
\textbf{Fréquences} & 0,16 & 0,30 & 0,24 & 0,20 & 0,10 & 1 \\ \hline
\end{tabular}
\end{center}
\end{exemplebox}

\begin{enumerate}
    \item \textbf{Classes de même amplitude :} $a = 2$ km pour toutes. On peut tracer l'histogramme des effectifs directement.
    \item \textbf{Axe des abscisses :} Bornes $0, 2, 4, 6, 8, 10$. Échelle : 1 cm pour 1 km (ou 2 km selon la feuille).
    \item \textbf{Axe des ordonnées :} Effectifs max = 15. Échelle : 1 cm pour 2 élèves (ou 1 pour 1).
    \item \textbf{Rectangles :}
    \begin{itemize}
        \item $[0;2[$ : hauteur 8
        \item $[2;4[$ : hauteur 15 $\leftarrow$ \textbf{Classe modale} (rectangle le plus haut)
        \item $[4;6[$ : hauteur 12
        \item $[6;8[$ : hauteur 10
        \item $[8;10[$ : hauteur 5
    \end{itemize}
\end{enumerate}

\begin{popfigure}
\begin{tikzpicture}
\begin{axis}[
    ybar interval,
    width=\linewidth,
    height=8cm,
    xlabel={Distance (km)},
    ylabel={Effectif},
    xtick={0,2,4,6,8,10},
    xticklabels={0,2,4,6,8,10},
    ytick={0,5,10,15},
    ymin=0, ymax=17,
    xmin=0, xmax=10,
    axis lines*=left,
    enlarge x limits=false,
    popblue!70,
    tick label style={font=\small},
    label style={font=\small},
    title={Histogramme des distances domicile-lycée},
    title style={font=\normalsize, yshift=0.5cm}
]
\addplot coordinates {
    (0,8) (2,15) (4,12) (6,10) (8,5) (10,0)
};
\end{axis}
\end{tikzpicture}
\legende{Histogramme des effectifs. La classe modale $[2;4[$ apparaît clairement comme le rectangle le plus haut.}
\end{popfigure}

\begin{aretenir}[Lecture de l'histogramme]
\begin{itemize}
    \item La \textbf{classe modale} est visuellement le rectangle de plus grande hauteur (ici $[2;4[$).
    \item La \textbf{répartition} se voit : concentration autour de 3 km, décroissance progressive pour les grandes distances.
    \item On peut estimer visuellement la médiane (classe où se trouve le 25ème élève) : effectifs cumulés $8 + 15 = 23 < 25$, $23 + 12 = 35 \ge 25$. La médiane est dans la classe $[4;6[$.
\end{itemize}
\end{aretenir}

\courssub{Le polygone des effectifs : une vue d'ensemble}

Souvent, on superpose à l'histogramme une ligne brisée qui relie les milieux des sommets des rectangles. C'est le \textbf{polygone des effectifs} (ou polygone des fréquences). Il permet de visualiser la \emph{tendance} générale de la distribution, un peu comme une courbe lissée.

\begin{definition}[Polygone des effectifs]
Pour chaque classe d'effectif $n_i$ et de milieu $m_i = \frac{\text{borne inf} + \text{borne sup}}{2}$, on place le point $M_i(m_i ; n_i)$. Le polygone des effectifs est la ligne brisée $M_1 M_2 \dots M_k$.
On ajoute souvent deux points fictifs aux extrémités (milieu de la classe précédente et suivante, d'effectif 0) pour « fermer » le polygone sur l'axe des abscisses.
\end{definition}

\begin{popfigure}
\begin{tikzpicture}
\begin{axis}[
    ybar interval,
    width=\linewidth,
    height=8cm,
    xlabel={Distance (km)},
    ylabel={Effectif},
    xtick={0,2,4,6,8,10},
    xticklabels={0,2,4,6,8,10},
    ytick={0,5,10,15},
    ymin=0, ymax=17,
    xmin=0, xmax=10,
    axis lines*=left,
    enlarge x limits=false,
    popblue!40,
    tick label style={font=\small},
    label style={font=\small},
    title={Histogramme + Polygone des effectifs},
    title style={font=\normalsize, yshift=0.5cm}
]
% Histogramme
\addplot coordinates {
    (0,8) (2,15) (4,12) (6,10) (8,5) (10,0)
};
% Polygone : milieux des classes = 1, 3, 5, 7, 9. Effectifs = 8, 15, 12, 10, 5.
% Points fictifs : (-1,0) et (11,0) pour fermer.
\addplot[poporange, very thick, mark=*, mark size=2pt, mark options={fill=poporange}] coordinates {
    (-1,0) (1,8) (3,15) (5,12) (7,10) (9,5) (11,0)
};
\addlegendimage{popblue!40, area legend}
\addlegendimage{poporange, very thick, mark=*}
\legend{Effectifs (Histogramme), Polygone des effectifs}
\end{axis}
\end{tikzpicture}
\legende{Superposition de l'histogramme (bleu) et du polygone des effectifs (orange). Le polygone relie les milieux des classes (1, 3, 5, 7, 9) et se ferme sur l'axe des abscisses.}
\end{popfigure}

\begin{attention}[Polygone et classes de même amplitude]
Le polygone des effectifs n'a de sens visuel direct (liaison des sommets des rectangles) que si les classes ont la \textbf{même amplitude}. Si les amplitudes diffèrent, les milieux ne sont plus au centre des rectangles de l'histogramme (qui ont des largeurs différentes), et le polygone ne « colle » plus aux rectangles. Au programme de 3ème, on reste dans le cas des classes de même amplitude.
\end{attention}

\courssub{Vérification de cohérence : tableau $\leftrightarrow$ histogramme}

Avant de rendre une copie ou de présenter un rapport, il faut toujours vérifier que le graphique correspond exactement au tableau. C'est une habitude de rigueur.

\begin{methode}[Vérifier la cohérence Tableau / Histogramme]
\begin{enumerate}
    \item \textbf{Nombre de rectangles} = Nombre de classes du tableau.
    \item \textbf{Bornes sur l'axe horizontal} : Les graduations principales correspondent exactement aux bornes inférieures et supérieures du tableau. Pas de décalage !
    \item \textbf{Hauteurs (effectifs)} : Pour chaque classe, la hauteur du rectangle (lu sur l'axe vertical) correspond \textbf{exactement} à l'effectif du tableau (en tenant compte de l'échelle).
    \item \textbf{Juxtaposition} : Aucun espace entre les rectangles. Le trait vertical entre deux rectangles tombe pile sur la borne commune.
    \item \textbf{Total des effectifs} : La somme des hauteurs (si échelle 1:1) ou la somme des aires correspond au total $N$.
    \item \textbf{Classe modale} : Le rectangle le plus haut correspond à la classe de plus grand effectif dans le tableau.
\end{enumerate}
\end{methode}

\begin{exoresolu}[Vérification d'un histogramme]
\textbf{Énoncé :} Un élève a tracé l'histogramme ci-contre pour représenter la série suivante :
\[
\begin{array}{|c|c|c|c|c|}\hline
\text{Classes} & [10;15[ & [15;20[ & [20;25[ & [25;30[ \\ \hline
\text{Effectifs} & 4 & 10 & 7 & 3 \\ \hline
\end{array}
\]
L'axe vertical est gradué de 1 en 1. L'axe horizontal est gradué de 10 en 10 (10, 20, 30). Les rectangles ont des bases : $[10;15]$, $[15;20]$, $[20;25]$, $[25;30]$. Hauteurs mesurées : 4, 10, 7, 3.
\textbf{Question :} Y a-t-il des erreurs ? Justifie.

\tcblower
\textbf{Solution :}
\begin{enumerate}
    \item \textbf{Nombre de rectangles :} 4 rectangles pour 4 classes. \checkmark
    \item \textbf{Bornes horizontales :} Le tableau a des bornes 10, 15, 20, 25, 30. L'axe est gradué 10, 20, 30. \textbf{Erreur !} Les bornes 15 et 25 ne sont pas graduées (ou pas visibles). Impossible de placer correctement les limites des classes $[10;15[$ et $[15;20[$ etc. sans les graduations 15 et 25.
    \item \textbf{Hauteurs :} 4, 10, 7, 3 correspondent au tableau. \checkmark
    \item \textbf{Juxtaposition :} Les rectangles se touchent. \checkmark
    \item \textbf{Conclusion :} L'histogramme est \textbf{incorrect} car l'axe des abscisses ne porte pas toutes les bornes des classes (manquent 15 et 25). Il faut graduer 10, 15, 20, 25, 30.
\end{enumerate}
\end{exoresolu}

\courssub{Exercices d'application}

\begin{exercice}[Construction d'histogramme]
Dans une plantation de palmiers à huile près de Douala, on a mesuré la hauteur (en mètres) de 60 palmiers. Résultats regroupés :
\[
\begin{array}{|c|c|c|c|c|c|}\hline
\text{Hauteur (m)} & [4;6[ & [6;8[ & [8;10[ & [10;12[ & [12;14[ \\ \hline
\text{Effectif} & 5 & 12 & 20 & 15 & 8 \\ \hline
\end{array}
\]
\begin{enumerate}
    \item Construis l'histogramme de cette série (choisis tes échelles).
    \item Trace le polygone des effectifs sur le même graphique.
    \item Quelle est la classe modale ? Interprète-la dans le contexte.
\end{enumerate}
\end{exercice}

\begin{corrige}[Exercice 1]
\begin{enumerate}
    \item \textbf{Axe horizontal :} Bornes 4, 6, 8, 10, 12, 14. Échelle : 1 cm pour 1 m (ou 2 m).
    \textbf{Axe vertical :} Effectif max = 20. Échelle : 1 cm pour 2 palmiers (ou 1 pour 1).
    \textbf{Rectangles :} Bases $[4;6]$ (h=5), $[6;8]$ (h=12), $[8;10]$ (h=20), $[10;12]$ (h=15), $[12;14]$ (h=8). Rectangles juxtaposés.
    \item \textbf{Polygone :} Milieux : 5, 7, 9, 11, 13. Points : $(5;5), (7;12), (9;20), (11;15), (13;8)$. Points fictifs : $(3;0)$ et $(15;0)$. Tracer la ligne brisée.
    \item \textbf{Classe modale :} $[8;10[$ (effectif 20, rectangle le plus haut). \textbf{Interprétation :} La majorité des palmiers de cette plantation ont une hauteur comprise entre 8 et 10 mètres. C'est la classe la plus représentée.
\end{enumerate}
\end{corrige}

\begin{exercice}[Lecture et analyse critique]
On a relevé les notes sur 20 d'un devoir de mathématiques pour une classe de 30 élèves. Un élève propose l'histogramme ci-dessous (description textuelle) :
\begin{itemize}
    \item Classes : $[0;5[$, $[5;10[$, $[10;15[$, $[15;20[$.
    \item Axe horizontal gradué : 0, 5, 10, 15, 20.
    \item Rectangles : Hauteurs 4, 10, 12, 4. \textbf{Espaces de 0,5 cm entre chaque rectangle.}
\end{itemize}
\textbf{Questions :}
\begin{enumerate}
    \item Reconstitue le tableau d'effectifs correspondant.
    \item Cet histogramme est-il correctement tracé ? Justifie précisément.
    \item Calcule la fréquence de la classe modale.
\end{enumerate}
\end{exercice}

\begin{corrige}[Exercice 2]
\begin{enumerate}
    \item \textbf{Tableau :}
    \[
    \begin{array}{|c|c|c|c|c|}\hline
    \text{Classes} & [0;5[ & [5;10[ & [10;15[ & [15;20[ \\ \hline
    \text{Effectifs} & 4 & 10 & 12 & 4 \\ \hline
    \end{array}
    \]
    \item \textbf{Non, l'histogramme n'est pas correct.}
    \textbf{Erreur majeure :} Il y a des \textbf{espaces entre les rectangles}. Les notes sont une variable quantitative continue (même si notées sur 20, on les regroupe en classes continues). Les classes se suivent : $[0;5[$, $[5;10[$, etc. La borne 5 est la fin de la première et le début de la deuxième. Les rectangles doivent être \textbf{juxtaposés} (collés). L'espace suggère une discontinuité qui n'existe pas.
    \item \textbf{Classe modale :} $[10;15[$ (effectif 12). Fréquence $f = \frac{12}{30} = 0,4$ soit $40\%$.
\end{enumerate}
\end{corrige}

\begin{savaistu}[L'histogramme en photographie : l'exposition]
Tu as sûrement vu un graphique bizarre sur l'écran de ton appareil photo ou ton smartphone (mode « Pro » ou retouche) : un histogramme \textbf{RVB} (Rouge, Vert, Bleu) ou \textbf{Luminosité}.
L'axe horizontal représente les \textbf{niveaux de luminosité} (de 0 = noir pur à 255 = blanc pur). L'axe vertical représente le \textbf{nombre de pixels} (effectif) ayant ce niveau.
\begin{itemize}
    \item Histogramme « collé à gauche » : Photo sous-exposée (trop sombre, beaucoup de pixels noirs).
    \item Histogramme « collé à droite » : Photo surexposée (trop claire, beaucoup de pixels blancs).
    \item Histogramme étalé au centre : Photo bien exposée, bon contraste.
\end{itemize}
C'est exactement le même outil mathématique qu'en cours : une variable continue (luminosité) regroupée en classes (les 256 niveaux), représentée par des barres juxtaposées (aire proportionnelle au nombre de pixels). Les photographes et les monteurs vidéo (pour la colorimétrie) l'utilisent en permanence pour corriger l'image sans même regarder la photo ! Les mathématiques sont partout, même dans tes selfies.
\end{savaistu}

\begin{aretenir}[Synthèse : L'histogramme en 3ème]
\begin{itemize}
    \item Représente une série \textbf{regroupée en classes de même amplitude}.
    \item Rectangles \textbf{juxtapossés} (pas d'espace).
    \item Base = intervalle de la classe (bornes sur l'axe horizontal).
    \item Hauteur = effectif (ou fréquence) de la classe (échelle sur l'axe vertical).
    \item \textbf{Classe modale} = Rectangle le plus haut.
    \item \textbf{Polygone des effectifs} = Ligne brisée reliant les milieux des sommets (utile pour voir la forme de la distribution).
    \item \textbf{Vérification systématique :} Bornes graduées, hauteurs exactes, pas d'espaces, total cohérent.
\end{itemize}
\end{aretenir}

\coursec{Diagrammes circulaires et semi-circulaires}

Dans la section précédente, nous avons appris à représenter une série regroupée en classes par un \cle{histogramme} : l'aire de chaque rectangle est proportionnelle à l'effectif de la classe. L'histogramme est parfait pour comparer les effectifs entre eux et visualiser la répartition des valeurs le long d'un axe. Mais il a une limite : il montre mal la \emph{part de chaque classe dans le total}. Or, dans la vie courante, on veut souvent répondre à des questions comme : « Quelle part des élèves prend le taxi pour venir au \textbf{collège} ? », « Quelle fraction du budget familial part dans la nourriture ? », « Quel pourcentage de candidats a réussi au BEPC ? ». Pour ce type de questions, on utilise un graphique en forme de disque : le \cle{diagramme circulaire}, surnommé « camembert ». C'est l'objet de cette section, avec son cousin le \cle{diagramme semi-circulaire}.

\courssub{Principe : un angle proportionnel à la fréquence}

\textbf{L'intuition.} Imagine une galette partagée entre les membres d'une famille : celui qui a le plus faim reçoit la plus grande part. Un diagramme circulaire fonctionne exactement ainsi : le disque entier représente l'\emph{effectif total} (ou 100 \%), et chaque classe reçoit un « morceau » — un \cle{secteur} — dont la taille est proportionnelle à son importance dans la série.

Or, la « taille » d'un morceau de disque se mesure par l'\cle{angle au centre} de son secteur. Le disque entier correspond à un angle de $360^\circ$. Une classe qui représente la moitié de l'effectif total aura donc un secteur de $180^\circ$ ; une classe qui représente le quart aura un secteur de $90^\circ$, et ainsi de suite.

\begin{propriete}[Angle d'un secteur circulaire]
Dans un diagramme circulaire, l'angle au centre du secteur associé à une classe est \textbf{proportionnel à la fréquence} de cette classe :
\[ \text{angle du secteur} = \text{fréquence de la classe} \times 360^\circ. \]
Pour un diagramme \textbf{semi-circulaire} (le total est représenté par un demi-disque de $180^\circ$) :
\[ \text{angle du secteur} = \text{fréquence de la classe} \times 180^\circ. \]
On peut aussi calculer l'angle directement avec l'effectif : $\text{angle} = \dfrac{\text{effectif de la classe}}{\text{effectif total}} \times 360^\circ$.
\end{propriete}

\begin{experience}[Pourquoi la fréquence fois $360^\circ$ ?]
Justifions la formule pas à pas, en utilisant uniquement la proportionnalité.
\begin{enumerate}
\item Le disque entier représente l'effectif total $N$. Il correspond à un tour complet, c'est-à-dire à un angle de $360^\circ$.
\item Une classe d'effectif $n_i$ représente la fraction $\dfrac{n_i}{N}$ de l'effectif total : c'est précisément sa \cle{fréquence} $f_i$.
\item Par proportionnalité, le secteur de cette classe doit représenter la \emph{même fraction} du tour complet :
\[ \text{angle}_i = \frac{n_i}{N} \times 360^\circ = f_i \times 360^\circ. \]
\item Vérification : la somme des fréquences vaut $1$, donc la somme des angles vaut $1 \times 360^\circ = 360^\circ$ : les secteurs remplissent exactement le disque, sans trou ni chevauchement. $\square$
\end{enumerate}
Pour le diagramme semi-circulaire, le raisonnement est identique, en remplaçant le tour complet ($360^\circ$) par un demi-tour ($180^\circ$).
\end{experience}

\begin{attention}[Piège n°1 : ne jamais utiliser les effectifs comme angles]
L'erreur la plus fréquente au BEPC consiste à reporter directement les effectifs comme des angles : « la classe a un effectif de 12, donc je trace un secteur de $12^\circ$ ». C'est \textbf{faux} ! Un effectif n'est pas un angle. Il faut \textbf{toujours} passer par la fréquence (ou par le quotient $\frac{n_i}{N}$) avant de calculer l'angle. Deuxième piège : oublier de vérifier que la somme des angles calculés vaut bien $360^\circ$ (ou $180^\circ$) ; si ce n'est pas le cas, il y a une erreur de calcul ou d'arrondi.
\end{attention}

\courssub{Tableau de conversion : effectifs, fréquences, angles}

En pratique, avant de tracer quoi que ce soit, on complète un \cle{tableau de conversion} à quatre lignes : classes, effectifs, fréquences, angles. Reprenons l'exemple des notes de mathématiques d'une classe de 3ème de $40$ élèves, regroupées en classes d'amplitude $4$.

\begin{center}
\begin{tabularx}{\linewidth}{|l|X|X|X|X|X|}
\hline
\rowcolor{popblueL} \textbf{Note} & $[0;4[$ & $[4;8[$ & $[8;12[$ & $[12;16[$ & $[16;20]$ \\
\hline \textbf{Effectif} $n_i$ & $4$ & $8$ & $16$ & $9$ & $3$ \\
\hline \textbf{Fréquence} $f_i$ & $0{,}10$ & $0{,}20$ & $0{,}40$ & $0{,}225$ & $0{,}075$ \\
\hline \textbf{Angle} $= f_i \times 360^\circ$ & $36^\circ$ & $72^\circ$ & $144^\circ$ & $81^\circ$ & $27^\circ$ \\
\hline
\end{tabularx}
\end{center}

Détail des calculs pour deux classes :
\begin{align*}
\text{Classe } [0;4[ : \quad & f_1 = \frac{4}{40} = 0{,}10 \quad ; \quad \text{angle} = 0{,}10 \times 360^\circ = 36^\circ. \\
\text{Classe } [8;12[ : \quad & f_3 = \frac{16}{40} = 0{,}40 \quad ; \quad \text{angle} = 0{,}40 \times 360^\circ = 144^\circ.
\end{align*}

Vérification : $36 + 72 + 144 + 81 + 27 = 360$. Les secteurs rempliront exactement le disque.

\begin{methode}[Construire un diagramme circulaire en 3 étapes]
\begin{enumerate}
\item \textbf{Calculer les angles.} Complète le tableau de conversion : effectifs $\rightarrow$ fréquences $\rightarrow$ angles ($f_i \times 360^\circ$). Vérifie que la somme des angles vaut $360^\circ$.
\item \textbf{Tracer au rapporteur.} Trace un cercle et un premier rayon (par exemple vertical, vers le haut). Place le rapporteur sur ce rayon, mesure le premier angle et trace le deuxième rayon. Reprends le rapporteur sur ce nouveau rayon pour mesurer l'angle suivant, et ainsi de suite, en tournant toujours dans le même sens.
\item \textbf{Légender.} Colorie chaque secteur d'une couleur différente, puis écris à l'intérieur (ou dans une légende à côté) la classe et son pourcentage. N'oublie pas le \textbf{titre} du diagramme.
\end{enumerate}
Pour un diagramme semi-circulaire, la démarche est identique : on part d'un demi-cercle et on multiplie les fréquences par $180^\circ$.
\end{methode}

\begin{popfigure}
\begin{center}
\begin{tikzpicture}[scale=1.6]
% angles cumulés : 0, 36, 108, 252, 333, 360 (sens trigonométrique)
\fill[popblue!70] (0,0) -- (90:1.6) arc (90:126:1.6) -- cycle;
\fill[popgreen!70] (0,0) -- (126:1.6) arc (126:198:1.6) -- cycle;
\fill[poporange!80] (0,0) -- (198:1.6) arc (198:342:1.6) -- cycle;
\fill[poppurple!60] (0,0) -- (342:1.6) arc (342:423:1.6) -- cycle;
\fill[poppink!70] (0,0) -- (423:1.6) arc (423:450:1.6) -- cycle;
\draw[thick] (0,0) circle (1.6);
\draw (0,0) -- (90:1.6); \draw (0,0) -- (126:1.6); \draw (0,0) -- (198:1.6);
\draw (0,0) -- (342:1.6); \draw (0,0) -- (423:1.6);
\node at (108:1.05) {\small $10\,\%$};
\node at (162:1.05) {\small $20\,\%$};
\node at (270:1.05) {\small $40\,\%$};
\node at (22:1.05) {\small $22{,}5\,\%$};
\node at (436:1.15) {\small $7{,}5\,\%$};
% légende
\node[anchor=west] at (2.1,1.1) {\small \textcolor{popblue!70}{$\blacksquare$} $[0;4[$ : $36^\circ$};
\node[anchor=west] at (2.1,0.7) {\small \textcolor{popgreen!70}{$\blacksquare$} $[4;8[$ : $72^\circ$};
\node[anchor=west] at (2.1,0.3) {\small \textcolor{poporange!80}{$\blacksquare$} $[8;12[$ : $144^\circ$};
\node[anchor=west] at (2.1,-0.1) {\small \textcolor{poppurple!60}{$\blacksquare$} $[12;16[$ : $81^\circ$};
\node[anchor=west] at (2.1,-0.5) {\small \textcolor{poppink!70}{$\blacksquare$} $[16;20]$ : $27^\circ$};
\end{tikzpicture}
\end{center}
\legende{Diagramme circulaire de la répartition des notes de la classe de 3ème ($40$ élèves). Chaque secteur est proportionnel à la fréquence de sa classe.}
\end{popfigure}

\courssub{Le diagramme semi-circulaire}

Le diagramme semi-circulaire repose sur la même idée, mais le total est représenté par un \emph{demi-disque}. Il est très utilisé pour les répartitions en petit nombre de catégories (résultats d'un vote, moyens de transport, choix d'orientation), car il se lit facilement et prend moins de place en hauteur.

\begin{exemplebox}[Moyens de transport des élèves d'un collège de Douala]
Dans un collège de Douala, on a demandé aux $120$ élèves de 3ème comment ils viennent au collège. Résultats : $48$ à pied, $36$ en taxi, $24$ en moto-taxi, $12$ en voiture familiale. Construisons le tableau de conversion pour un diagramme \textbf{semi-circulaire} (multiplicateur : $180^\circ$).
\begin{center}
\begin{tabularx}{\linewidth}{|l|X|X|X|X|}
\hline
\rowcolor{popgreenL} \textbf{Transport} & À pied & Taxi & Moto-taxi & Voiture \\
\hline \textbf{Effectif} & $48$ & $36$ & $24$ & $12$ \\
\hline \textbf{Fréquence} & $0{,}40$ & $0{,}30$ & $0{,}20$ & $0{,}10$ \\
\hline \textbf{Angle} $= f_i \times 180^\circ$ & $72^\circ$ & $54^\circ$ & $36^\circ$ & $18^\circ$ \\
\hline
\end{tabularx}
\end{center}
Vérification : $72 + 54 + 36 + 18 = 180$. Parfait, le demi-disque sera exactement rempli.
\end{exemplebox}

\begin{popfigure}
\begin{center}
\begin{tikzpicture}[scale=1.5]
% demi-cercle de 0° à 180° : secteurs 72, 54, 36, 18
\fill[popblue!70] (0,0) -- (180:2) arc (180:108:2) -- cycle;
\fill[poporange!80] (0,0) -- (108:2) arc (108:54:2) -- cycle;
\fill[poppurple!60] (0,0) -- (54:2) arc (54:18:2) -- cycle;
\fill[poppink!70] (0,0) -- (18:2) arc (18:0:2) -- cycle;
\draw[thick] (-2,0) arc (180:0:2) -- cycle;
\draw (0,0) -- (108:2); \draw (0,0) -- (54:2); \draw (0,0) -- (18:2);
\node at (144:1.3) {\small À pied};
\node at (144:0.95) {\small $40\,\%$};
\node at (81:1.35) {\small Taxi};
\node at (81:1.0) {\small $30\,\%$};
\node at (36:1.35) {\small Moto};
\node at (36:1.0) {\small $20\,\%$};
\node at (9:1.45) {\small Voiture};
\node at (9:1.05) {\small $10\,\%$};
\draw[<->] (180:2.35) arc (180:108:2.35);
\node at (144:2.6) {\small $72^\circ$};
\end{tikzpicture}
\end{center}
\legende{Diagramme semi-circulaire de la répartition des moyens de transport des élèves de 3ème. L'angle total est $180^\circ$.}
\end{popfigure}

\courssub{Un exemple complet : le budget d'une famille de Yaoundé}

\begin{exoresolu}[Budget mensuel d'une famille de Yaoundé]
Une famille de Yaoundé dépense chaque mois \num{200000} FCFA, répartis ainsi : nourriture \num{80000} FCFA ; loyer \num{50000} FCFA ; transport \num{30000} FCFA ; scolarité des enfants \num{24000} FCFA ; autres dépenses \num{16000} FCFA. Construire le tableau de conversion et interpréter le diagramme circulaire obtenu.
\tcblower
\textbf{Solution détaillée.} L'effectif total est $N = \num{200000}$ FCFA. Pour chaque poste, on calcule la fréquence puis l'angle ($f_i \times 360^\circ$).
\begin{center}
\begin{tabularx}{\linewidth}{|l|X|X|X|}
\hline
\rowcolor{poporangeL} \textbf{Poste} & \textbf{Dépense (FCFA)} & \textbf{Fréquence} & \textbf{Angle} \\
\hline Nourriture & \num{80000} & $0{,}40$ & $144^\circ$ \\
\hline Loyer & \num{50000} & $0{,}25$ & $90^\circ$ \\
\hline Transport & \num{30000} & $0{,}15$ & $54^\circ$ \\
\hline Scolarité & \num{24000} & $0{,}12$ & $43{,}2^\circ$ \\
\hline Autres & \num{16000} & $0{,}08$ & $28{,}8^\circ$ \\
\hline \textbf{Total} & \num{200000} & $1$ & $360^\circ$ \\
\hline
\end{tabularx}
\end{center}
Exemple de calcul pour le loyer : $f = \dfrac{\num{50000}}{\num{200000}} = 0{,}25$, donc angle $= 0{,}25 \times 360^\circ = 90^\circ$ (un quart de tour : le loyer représente exactement un quart du budget).

\textbf{Interprétation.} Le plus grand secteur ($144^\circ$) correspond à la nourriture : c'est le poste dominant, avec $40\,\%$ du budget. Nourriture et loyer réunis occupent $144^\circ + 90^\circ = 234^\circ$, soit plus de la moitié du disque ($234^\circ > 180^\circ$) : ces deux postes absorbent à eux seuls $65\,\%$ des dépenses. On trace ensuite le diagramme au rapporteur en suivant la méthode en 3 étapes.
\end{exoresolu}

\courssub{Lire et comparer deux diagrammes circulaires}

Savoir \emph{lire} un diagramme est aussi important que savoir le construire. Deux règles d'or :

\begin{itemize}
\item \textbf{Comparer les parts au sein d'un même diagramme} : le plus grand secteur correspond à la classe de plus forte fréquence. On peut comparer directement les angles ou les pourcentages affichés.
\item \textbf{Comparer deux diagrammes différents} : on compare les \emph{fréquences} (pourcentages ou angles), \textbf{jamais les effectifs}, car les deux séries n'ont pas forcément le même effectif total.
\end{itemize}

\begin{attention}[Piège n°2 : comparer des effectifs à partir de deux camemberts]
Supposons que le collège A ($200$ élèves) ait $30\,\%$ d'élèves venant à pied, et le collège B ($500$ élèves) $20\,\%$. Le secteur « à pied » est plus grand sur le diagramme de A, mais en \emph{nombre d'élèves}, c'est B qui en compte le plus : $0{,}20 \times 500 = 100$ élèves contre $0{,}30 \times 200 = 60$. Un diagramme circulaire donne des \textbf{proportions}, pas des \textbf{effectifs} : pour retrouver un effectif, il faut multiplier la fréquence par l'effectif total de la série concernée.
\end{attention}

\courssub{Quel graphique choisir : histogramme ou diagramme circulaire ?}

Le choix du graphique dépend du \emph{message} que l'on veut faire passer :

\begin{center}
\begin{tabularx}{\linewidth}{|l|X|X|}
\hline
\rowcolor{popblueL} & \textbf{Histogramme} & \textbf{Diagramme circulaire} \\
\hline \textbf{Message} & Comparer les effectifs, voir la forme de la répartition le long d'un axe gradué & Comparer les parts de chaque catégorie dans le total \\
\hline \textbf{Ce qui est proportionnel} & L'aire des rectangles & L'angle des secteurs \\
\hline \textbf{Idéal pour} & Une variable quantitative regroupée en classes (notes, tailles, âges) & Des catégories ou des classes en petit nombre, avec des pourcentages \\
\hline
\end{tabularx}
\end{center}

\begin{aretenir}[L'essentiel de la section]
\begin{itemize}
\item Dans un diagramme circulaire, l'angle de chaque secteur est \textbf{proportionnel à la fréquence} : $\text{angle} = f_i \times 360^\circ$ (et $f_i \times 180^\circ$ pour un semi-circulaire).
\item On prépare toujours un \textbf{tableau de conversion} : classes $\rightarrow$ effectifs $\rightarrow$ fréquences $\rightarrow$ angles, et on vérifie que la somme des angles vaut $360^\circ$ (ou $180^\circ$).
\item La construction se fait au \textbf{rapporteur}, secteur après secteur, puis on légende avec les pourcentages.
\item Un camembert montre des \textbf{proportions}, pas des effectifs : attention aux comparaisons entre deux diagrammes.
\end{itemize}
\end{aretenir}

\begin{savaistu}[Les camemberts du BEPC dans la presse camerounaise]
Chaque année, lorsque l'Office du Baccalauréat du Cameroun (OBC) publie les résultats du BEPC, les journaux et les sites d'information illustrent leurs articles par des diagrammes circulaires : part des admis, des ajournés, répartition par région ou par académie. Ces « camemberts » permettent à tout lecteur, même sans formation mathématique, de visualiser d'un coup d'œil qu'une région représente, par exemple, le quart des admis du pays. Le diagramme circulaire a été inventé en 1801 par le statisticien anglais William Playfair : plus de deux siècles plus tard, il reste l'un des graphiques les plus utilisés au monde... et l'un des plus rentables à maîtriser pour ton épreuve de mathématiques !
\end{savaistu}

\begin{exercice}[À toi de jouer]
Dans un village de la région de l'Ouest, on a recensé les cultures pratiquées par $60$ exploitations agricoles : $24$ cultivent le maïs, $18$ le manioc, $12$ le plantain et $6$ le maraîchage.
\begin{enumerate}
\item Complète le tableau de conversion (effectifs, fréquences, angles pour un diagramme circulaire).
\item Construis le diagramme circulaire au rapporteur sur ta copie.
\item Quelle culture représente exactement un cinquième du disque ? Justifie.
\end{enumerate}
\end{exercice}

\begin{corrige}[Exercice : cultures de l'Ouest]
\begin{enumerate}
\item Fréquences : maïs $\frac{24}{60} = 0{,}4$ ; manioc $\frac{18}{60} = 0{,}3$ ; plantain $\frac{12}{60} = 0{,}2$ ; maraîchage $\frac{6}{60} = 0{,}1$. Angles : maïs $0{,}4 \times 360^\circ = 144^\circ$ ; manioc $108^\circ$ ; plantain $72^\circ$ ; maraîchage $36^\circ$. Somme : $144 + 108 + 72 + 36 = 360^\circ$. Correct.
\item On trace un cercle, un premier rayon, puis on reporte successivement $144^\circ$, $108^\circ$, $72^\circ$ et $36^\circ$ au rapporteur, en coloriant chaque secteur et en indiquant les pourcentages ($40\,\%$, $30\,\%$, $20\,\%$, $10\,\%$).
\item Un cinquième du disque correspond à $\frac{360^\circ}{5} = 72^\circ$ : c'est le secteur du \textbf{plantain} (fréquence $0{,}2 = \frac{1}{5}$).
\end{enumerate}
\end{corrige}

\coursec{Synthèse et résolution de problèmes}

Dans les sections précédentes, tu as appris à regrouper des données en classes, à déterminer la classe modale, à calculer les fréquences et la moyenne, puis à représenter la série par un histogramme, un diagramme circulaire ou semi-circulaire. Il te reste maintenant à assembler toutes ces pièces : c'est exactement ce que l'on te demandera au BEPC, où un problème de statistiques se présente presque toujours comme une \emph{étude complète}, de la collecte des données jusqu'à la conclusion. Cette dernière section te donne la méthode, les formules et les réflexes de rédaction pour réussir ce type d'exercice à coup sûr.

\courssub{La démarche complète d'une étude statistique}

\textbf{Les cinq étapes}\par

Faire une étude statistique, c'est raconter une histoire avec des nombres : on part d'une réalité désordonnée (des dizaines de mesures, de ventes, d'âges...) et on aboutit à une phrase claire que tout le monde comprend. Cette histoire suit toujours le même chemin, en cinq étapes.

\begin{methode}[Démarche complète d'une étude statistique]
\begin{enumerate}
\item \textbf{Collecter les données} : recenser les valeurs du caractère étudié sur la population (enquête, mesures, registre de ventes...).
\item \textbf{Regrouper en classes et construire le tableau} : choisir des classes d'amplitude égale, compter les effectifs $n_i$, calculer les centres $c_i$, les fréquences $f_i$ et, si besoin, les angles des secteurs.
\item \textbf{Calculer les indicateurs} : classe modale, moyenne $\bar{x}$, fréquences remarquables.
\item \textbf{Tracer un graphique} : histogramme (effectifs) ou diagramme circulaire (fréquences), avec titre, axes et légende.
\item \textbf{Rédiger une conclusion} : interpréter les indicateurs dans le contexte concret du problème.
\end{enumerate}
\end{methode}

\begin{popfigure}
\begin{center}
\begin{tikzpicture}[
  etape/.style={rectangle, rounded corners=4pt, draw=popblue, fill=popblueL, text=popdark, align=center, minimum width=2.6cm, minimum height=1.2cm, font=\small},
  fleche/.style={-{Stealth[length=3mm]}, thick, poporange}
]
\node[etape, align=center] (e1) at (0,0){\textbf{1. Données}\\ collecte};
\node[etape, align=center] (e2) at (3.4,0){\textbf{2. Tableau}\\ classes, effectifs};
\node[etape, align=center] (e3) at (6.8,0){\textbf{3. Indicateurs}\\ moyenne, classe modale};
\node[etape, align=center] (e4) at (10.2,0){\textbf{4. Graphique}\\ histogramme, secteurs};
\node[etape, draw=popgreen, fill=popgreenL, align=center] (e5) at (13.6,0){\textbf{5. Conclusion}\\ interprétation};
\draw[fleche] (e1) -- (e2);
\draw[fleche] (e2) -- (e3);
\draw[fleche] (e3) -- (e4);
\draw[fleche] (e4) -- (e5);
\end{tikzpicture}
\end{center}
\legende{Les cinq étapes de la démarche statistique : chaque étape s'appuie sur la précédente.}
\end{popfigure}

\begin{attention}[Une étape oubliée = des points perdus]
Au BEPC, le barème récompense chaque étape : le tableau complet (avec les centres et les fréquences), le calcul de la moyenne \emph{avec la formule écrite}, le graphique soigné, et la conclusion. Ne saute jamais une étape, même si elle te semble facile : c'est souvent là que se cachent les points les plus simples à gagner.
\end{attention}

\courssub{Rédiger et justifier une conclusion statistique}

C'est l'étape que les élèves négligent le plus... et que les correcteurs surveillent le plus ! Une conclusion statistique n'est pas une impression (« il y en a beaucoup ») : c'est une affirmation \emph{appuyée sur un indicateur calculé}.

\begin{methode}[Rédiger une conclusion en trois temps]
Pour chaque affirmation de ta conclusion, respecte le schéma suivant :
\begin{enumerate}
\item \textbf{Citer l'indicateur} : « la classe modale est... », « la moyenne vaut... », « la fréquence de la classe... est... » ;
\item \textbf{Donner sa valeur} exacte, avec l'unité ou en pourcentage ;
\item \textbf{Interpréter dans le contexte} : dire ce que cette valeur signifie concrètement pour la situation (le commerçant, le médecin, l'agriculteur...).
\end{enumerate}
\emph{Modèle de phrase :} « La classe modale est $[20\,;30[$, donc la majorité des clients dépense entre \num{20000} et \num{30000} FCFA : la boutique doit surtout stocker des articles de cette gamme de prix. »
\end{methode}

\begin{attention}[Erreur fréquente : conclure « au feeling »]
Écrire « les ventes sont bonnes » ou « les jeunes sont nombreux » sans citer \textbf{aucun} indicateur calculé ne rapporte aucun point. Toute conclusion doit être \emph{justifiée} par un nombre issu de tes calculs : une moyenne, une classe modale, une fréquence. Pas de calcul, pas de conclusion !
\end{attention}

\courssub{Exemple chiffré : étude complète des ventes d'une boutique de téléphonie à Douala}

\begin{exoresolu}[Ventes journalières d'une boutique à Douala]
\textbf{Énoncé.} Une boutique de téléphonie mobile du quartier Deïdo à Douala a relevé le montant des ventes (en milliers de FCFA) réalisées chaque jour pendant 30 jours :
\begin{center}
12 ; 25 ; 8 ; 31 ; 45 ; 18 ; 22 ; 38 ; 15 ; 27 ; 41 ; 9 ; 33 ; 24 ; 19 ; 36 ; 28 ; 14 ; 48 ; 21 ; 30 ; 16 ; 39 ; 26 ; 11 ; 34 ; 23 ; 44 ; 17 ; 29
\end{center}
\begin{enumerate}
\item Regrouper ces données en classes d'amplitude 10 (à partir de $[5\,;15[$) et dresser le tableau complet (effectifs, centres, fréquences).
\item Déterminer la classe modale.
\item Calculer le montant moyen des ventes journalières.
\item Calculer l'angle du secteur de la classe $[25\,;35[$ dans un diagramme circulaire.
\item Rédiger une conclusion pour la gérante.
\end{enumerate}
\tcblower
\textbf{Solution détaillée.}

\textbf{1. Tableau.} On compte les valeurs classe par classe : $[5\,;15[$ contient 12, 8, 9, 14, 11 (5 valeurs) ; $[15\,;25[$ contient 18, 22, 15, 24, 19, 21, 16, 23, 17 (9 valeurs) ; $[25\,;35[$ contient 25, 31, 27, 33, 28, 30, 26, 34, 29 (9 valeurs) ; $[35\,;45[$ contient 45, 38, 41, 36, 39, 44 (6 valeurs) ; $[45\,;55[$ contient 48 (1 valeur). Vérification : $5+9+9+6+1=30$. On complète avec les centres $c_i=\dfrac{a+b}{2}$ et les fréquences $f_i=\dfrac{n_i}{N}$ :

\begin{center}
\begin{tabularx}{\linewidth}{|l|X|X|X|X|X|}
\hline
\rowcolor{popblueL} Classe (en milliers de FCFA) & $[5\,;15[$ & $[15\,;25[$ & $[25\,;35[$ & $[35\,;45[$ & $[45\,;55[$ \\
\hline
Effectif $n_i$ & 5 & 9 & 9 & 6 & 1 \\
\hline
Centre $c_i$ & 10 & 20 & 30 & 40 & 50 \\
\hline
Fréquence $f_i$ & $\dfrac{5}{30}\approx 0{,}167$ & $0{,}3$ & $0{,}3$ & $0{,}2$ & $\approx 0{,}033$ \\
\hline
$n_i \times c_i$ & 50 & 180 & 270 & 240 & 50 \\
\hline
\end{tabularx}
\end{center}

\textbf{2. Classe modale.} Les classes $[15\,;25[$ et $[25\,;35[$ ont le même effectif maximal (9) : la série admet \textbf{deux classes modales}, $[15\,;25[$ et $[25\,;35[$.

\textbf{3. Moyenne.} On applique la formule de la moyenne d'une série regroupée en classes :
\[
\bar{x}=\frac{\sum n_i c_i}{N}=\frac{50+180+270+240+50}{30}=\frac{790}{30}\approx 26{,}3.
\]
Le montant moyen des ventes est d'environ \textbf{\num{26300} FCFA par jour}.

\textbf{4. Angle du secteur.} Pour la classe $[25\,;35[$, la fréquence est $f=\dfrac{9}{30}=0{,}3$, donc :
\[
\text{angle}=f \times 360^\circ = 0{,}3 \times 360^\circ = 108^\circ.
\]

\textbf{5. Conclusion rédigée.} « Les classes modales sont $[15\,;25[$ et $[25\,;35[$ : la plupart des journées (18 jours sur 30, soit une fréquence de $0{,}6$), la boutique vend entre \num{15000} et \num{35000} FCFA. La moyenne est d'environ \num{26300} FCFA par jour. Les grosses journées (plus de \num{45000} FCFA) sont rares (fréquence $\approx 3{,}3\,\%$). La gérante peut donc prévoir sa trésorerie sur une base d'environ \num{26000} FCFA quotidiens. »
\end{exoresolu}

\courssub{Lecture critique : détecter un graphique ou une conclusion erronée}

Au BEPC comme dans la vie, on te montrera parfois des graphiques trompeurs ou des conclusions hâtives. Voici les trois pièges classiques à savoir repérer — et à éviter dans tes propres copies.

\begin{attention}[Les trois erreurs de lecture les plus fréquentes]
\begin{itemize}
\item \textbf{Histogramme à classes d'amplitudes inégales} : si les classes n'ont pas la même amplitude, comparer les \emph{hauteurs} des rectangles est trompeur ; au programme de 3ème, on exige des classes d'\emph{amplitude égale}, donc des hauteurs proportionnelles aux effectifs. Vérifie toujours les bornes des classes avant de conclure.
\item \textbf{Confusion effectif / fréquence} : dire « la classe $[35\,;45[$ représente $6\,\%$ » alors que 6 est l'\emph{effectif} est faux ; la fréquence est $\dfrac{6}{30}=0{,}2$, soit $20\,\%$.
\item \textbf{Conclusion non justifiée} : affirmer « les ventes augmentent » à partir d'une seule série statique est un abus : une moyenne ou une classe modale décrit une situation à un instant donné, pas une évolution.
\end{itemize}
\end{attention}

\begin{exemplebox}[Critique d'une conclusion erronée]
Un élève affirme, à partir de l'exemple de la boutique : « La classe $[45\,;55[$ a un effectif de 1, donc la boutique ne fait jamais de grosses ventes. »

\textbf{Critique :} l'effectif 1 sur 30 correspond à une fréquence de $\dfrac{1}{30}\approx 3{,}3\,\%$ : les grosses ventes sont \emph{rares}, pas \emph{inexistantes}. De plus, un seul mois d'observation ne permet pas d'affirmer ce qui se passe « jamais ». La conclusion correcte est : « les ventes supérieures à \num{45000} FCFA sont rares (environ $3{,}3\,\%$ des journées observées) ».
\end{exemplebox}

\courssub{Problèmes de la vie courante}

La même démarche s'applique à tous les domaines. Entraîne-toi à reconnaître, dans chaque situation, la population, le caractère, les classes et l'indicateur pertinent.

\begin{exercice}[Rendements agricoles à Bafoussam]
Une coopérative de la région de l'Ouest a mesuré le rendement en maïs (en tonnes par hectare) de 40 parcelles. Les résultats, regroupés en classes, sont :
\begin{center}
\begin{tabularx}{\linewidth}{|l|X|X|X|X|}
\hline
\rowcolor{popgreenL} Rendement (t/ha) & $[1\,;2[$ & $[2\,;3[$ & $[3\,;4[$ & $[4\,;5[$ \\
\hline
Nombre de parcelles & 6 & 14 & 12 & 8 \\
\hline
\end{tabularx}
\end{center}
\begin{enumerate}
\item Calculer les fréquences de chaque classe.
\item Déterminer la classe modale et l'interpréter.
\item Calculer le rendement moyen par parcelle.
\item La coopérative estime un rendement « satisfaisant » à partir de 3 t/ha. Quelle proportion des parcelles atteint ce seuil ?
\end{enumerate}
\end{exercice}

\begin{corrige}[Exercice : rendements agricoles]
\textbf{1.} Fréquences : $\dfrac{6}{40}=0{,}15$ ; $\dfrac{14}{40}=0{,}35$ ; $\dfrac{12}{40}=0{,}3$ ; $\dfrac{8}{40}=0{,}2$. La somme vaut bien $1$.

\textbf{2.} La classe modale est $[2\,;3[$ (effectif maximal 14) : le rendement le plus fréquent se situe entre 2 et 3 tonnes par hectare.

\textbf{3.} Centres : $1{,}5$ ; $2{,}5$ ; $3{,}5$ ; $4{,}5$. Moyenne :
\[
\bar{x}=\frac{6\times 1{,}5+14\times 2{,}5+12\times 3{,}5+8\times 4{,}5}{40}=\frac{9+35+42+36}{40}=\frac{122}{40}=3{,}05.
\]
Le rendement moyen est de $3{,}05$ t/ha.

\textbf{4.} Les classes $[3\,;4[$ et $[4\,;5[$ dépassent le seuil : $12+8=20$ parcelles, soit une fréquence de $\dfrac{20}{40}=0{,}5$, c'est-à-dire $50\,\%$. Conclusion : la moitié des parcelles atteint le rendement satisfaisant, et la moyenne de $3{,}05$ t/ha confirme un niveau global correct, tiré vers le haut par les bonnes parcelles.
\end{corrige}

\begin{exercice}[Classes d'âge des patients d'un centre de santé]
Un centre de santé de Yaoundé a enregistré l'âge de ses 200 patients une semaine donnée :
\begin{center}
\begin{tabularx}{\linewidth}{|l|X|X|X|X|X|}
\hline
\rowcolor{poppinkL} Âge (ans) & $[0\,;15[$ & $[15\,;30[$ & $[30\,;45[$ & $[45\,;60[$ & $[60\,;75[$ \\
\hline
Effectif & 50 & 70 & 40 & 26 & 14 \\
\hline
\end{tabularx}
\end{center}
\begin{enumerate}
\item Calculer l'âge moyen des patients.
\item Calculer l'angle du secteur des moins de 15 ans dans un diagramme circulaire.
\item Le médecin-chef affirme : « plus de la moitié de nos patients ont moins de 30 ans ». A-t-il raison ? Justifier par un calcul.
\end{enumerate}
\end{exercice}

\begin{corrige}[Exercice : classes d'âge]
\textbf{1.} Centres : $7{,}5$ ; $22{,}5$ ; $37{,}5$ ; $52{,}5$ ; $67{,}5$.
\[
\bar{x}=\frac{50\times 7{,}5+70\times 22{,}5+40\times 37{,}5+26\times 52{,}5+14\times 67{,}5}{200}=\frac{375+1575+1500+1365+945}{200}=\frac{5760}{200}=28{,}8.
\]
L'âge moyen est de $28{,}8$ ans.

\textbf{2.} Fréquence des moins de 15 ans : $\dfrac{50}{200}=0{,}25$. Angle : $0{,}25 \times 360^\circ = 90^\circ$.

\textbf{3.} Patients de moins de 30 ans : $50+70=120$, soit une fréquence de $\dfrac{120}{200}=0{,}6=60\,\%$. Or $60\,\% > 50\,\%$ : le médecin-chef a \textbf{raison}, et l'affirmation est justifiée par le calcul de la fréquence cumulée.
\end{corrige}

\courssub{Bilan des formules de la leçon}

\begin{aretenir}[Formulaire de la leçon « Statistiques »]
Pour une série regroupée en classes, d'effectif total $N$, de classes d'effectifs $n_i$ et de centres $c_i$ :
\begin{center}
\begin{tabularx}{\linewidth}{|l|X|}
\hline
\rowcolor{popblueL} \textbf{Notion} & \textbf{Formule / définition} \\
\hline
Centre d'une classe $[a\,;b[$ & $c=\dfrac{a+b}{2}$ \\
\hline
Fréquence d'une classe & $f_i=\dfrac{n_i}{N}$ \quad (nombre entre 0 et 1 ; en \% : $f_i \times 100$) \\
\hline
Somme des fréquences & $f_1+f_2+\dots+f_k=1$ \quad (soit $100\,\%$) \\
\hline
Classe modale & classe d'\textbf{effectif maximal} (ou de fréquence maximale) \\
\hline
Moyenne de la série & $\bar{x}=\dfrac{n_1 c_1+n_2 c_2+\dots+n_k c_k}{N}=\dfrac{\sum n_i c_i}{N}$ \\
\hline
Angle d'un secteur (diagramme circulaire) & $\alpha_i = f_i \times 360^\circ$ \\
\hline
Angle d'un secteur (diagramme semi-circulaire) & $\alpha_i = f_i \times 180^\circ$ \\
\hline
\end{tabularx}
\end{center}
\end{aretenir}

\begin{attention}[Dernières vérifications avant de rendre ta copie]
\begin{itemize}
\item La somme des effectifs est-elle égale à l'effectif total $N$ ?
\item La somme des fréquences vaut-elle bien $1$ (ou $100\,\%$) ? La somme des angles vaut-elle $360^\circ$ (ou $180^\circ$) ?
\item As-tu utilisé les \emph{centres} des classes (et non les bornes) pour calculer la moyenne ?
\item Ta conclusion cite-t-elle au moins un indicateur avec sa valeur et son interprétation ?
\end{itemize}
\end{attention}

\begin{savaistu}[Les statistiques au BEPC]
Au BEPC, les exercices de statistiques demandent presque toujours, \emph{dans le même problème}, la construction d'un tableau (effectifs, centres, fréquences), le calcul de la moyenne et le tracé d'un histogramme ou d'un diagramme circulaire — parfois complété par une question de conclusion ouverte. C'est l'un des chapitres où la méthode paie le plus : un élève qui applique rigoureusement les cinq étapes de cette section y gagne des points réguliers, même sans être un génie du calcul. Les statistiques sont aussi omniprésentes au Cameroun : recensements de la population (BUCREP), bulletins météo, statistiques scolaires du MINESEC, sondages d'opinion... Maîtriser cette leçon, c'est apprendre à lire le monde avec des nombres !
\end{savaistu}

\begin{exercice}[Pour réviser toute la leçon]
Un commerçant du marché Mokolo a relevé le montant de ses 50 ventes d'une journée, regroupées par tranche de prix (en FCFA) :
\begin{center}
\begin{tabularx}{\linewidth}{|l|X|X|X|X|}
\hline
\rowcolor{popgoldL} Montant (FCFA) & $[0\,;1000[$ & $[1000\,;2000[$ & $[2000\,;3000[$ & $[3000\,;4000[$ \\
\hline
Nombre de ventes & 10 & 20 & 15 & 5 \\
\hline
\end{tabularx}
\end{center}
\begin{enumerate}
\item Compléter le tableau avec les centres et les fréquences.
\item Déterminer la classe modale.
\item Calculer le montant moyen d'une vente.
\item Calculer les angles des quatre secteurs d'un diagramme circulaire, puis le construire.
\item Rédiger une conclusion de deux phrases pour le commerçant.
\end{enumerate}
\end{exercice}

\begin{corrige}[Exercice : révisions]
\textbf{1.} Centres : 500 ; \num{1500} ; \num{2500} ; \num{3500}. Fréquences : $0{,}2$ ; $0{,}4$ ; $0{,}3$ ; $0{,}1$ (somme : 1).

\textbf{2.} Classe modale : $[1000\,;2000[$ (effectif 20).

\textbf{3.} $\bar{x}=\dfrac{10\times 500+20\times 1500+15\times 2500+5\times 3500}{50}=\dfrac{5000+30000+37500+17500}{50}=\dfrac{90000}{50}=\num{1800}$ FCFA.

\textbf{4.} Angles : $0{,}2\times 360^\circ=72^\circ$ ; $0{,}4\times 360^\circ=144^\circ$ ; $0{,}3\times 360^\circ=108^\circ$ ; $0{,}1\times 360^\circ=36^\circ$. Vérification : $72+144+108+36=360^\circ$. Le diagramme se trace au rapporteur avec ces quatre secteurs.

\textbf{5.} Conclusion : « La classe modale est $[1000\,;2000[$ : la plupart des ventes (fréquence $0{,}4$, soit $40\,\%$) portent sur des montants entre \num{1000} et \num{2000} FCFA. La vente moyenne est de \num{1800} FCFA : le commerçant peut prévoir sa journée sur cette base et privilégier les articles de cette gamme de prix. »
\end{corrige}

\newpage

\coursec{Activité d'intégration}

\courssub{Objectifs de l'activité}
Cette activité d'intégration vise à mobiliser l'ensemble des savoir-faire acquis au cours de la leçon sur les statistiques (regroupement en classes, recherche de la classe modale, calcul des fréquences et de la moyenne, représentations graphiques par histogramme et diagramme circulaire) dans une situation authentique de collecte et d'exploitation de données réelles. Elle permet à l'élève de :
\begin{itemize}
    \item Mettre en œuvre une démarche statistique complète : de la collecte brute à la communication des résultats.
    \item Manipuler des effectifs importants justifiant le regroupement en classes.
    \item Interpréter les indicateurs de position (classe modale, moyenne) pour prendre une décision argumentée.
    \item Rédiger un rapport structuré, chiffré et adressé à un destinataire identifié (le proviseur), dans un souci de rigueur langagière et mathématique.
\end{itemize}

\courssub{Situation}
Au lycée bilingue de Nkolbisson (Yaoundé), l'administration souhaite mettre en place un \emph{ticket de transport scolaire subventionné} pour soulager les familles. Pour fixer le montant de la participation familiale, le proviseur a besoin de connaître la réalité des dépenses hebdomadaires de transport des élèves.
Une enquête anonyme est réalisée auprès de \textbf{40 élèves} de la classe de 3ème A, tirés au sort. Chaque élève indique sa dépense hebdomadaire (aller-retour du lundi au samedi) en francs CFA. Voici les données brutes recueillies (en FCFA) :

\begin{center}
\begin{tabular}{|c|c|c|c|c|c|c|c|c|c|c|}\hline
1200 & 1500 & 800 & 2200 & 1800 & 950 & 1100 & 2500 & 1300 & 1600 \\\hline
1400 & 1900 & 1000 & 1700 & 2100 & 1250 & 1550 & 900 & 1850 & 2000 \\\hline
1350 & 1650 & 1150 & 2300 & 1450 & 1050 & 1750 & 1950 & 1200 & 1500 \\\hline
1800 & 2400 & 1300 & 1600 & 1100 & 2050 & 1400 & 1700 & 1900 & 2200 \\\hline
\end{tabular}
\end{center}

\courssub{Consignes}
Travailler par groupes de 4 à 5 élèves. Chaque groupe rendra un rapport écrit structuré répondant aux questions suivantes :

\begin{enumerate}
    \item \textbf{Organisation des données :} Construire un tableau de regroupement en classes d'amplitude 500 FCFA, en commençant par la classe $[800 ; 1300[$. Compléter le tableau avec les colonnes : \emph{Classes}, \emph{Effectifs}, \emph{Centres des classes}, \emph{Fréquences (en \%)}, \emph{Angles (en degrés pour le diagramme circulaire)}.
    \item \textbf{Analyse des effectifs :} Identifier la \cle{classe modale}. Quelle proportion d'élèves (en \%) se trouve dans cette classe ? Interpréter ce résultat dans le contexte.
    \item \textbf{Calcul de la moyenne :} Calculer la dépense moyenne hebdomadaire de transport pour cet échantillon. Arrondir au FCFA le plus proche. Comparer cette moyenne aux bornes de la classe modale.
    \item \textbf{Représentations graphiques :}
    \begin{itemize}
        \item Tracer l'\cle{histogramme} de la série statistique (sur papier millimétré, en choisissant une échelle adaptée : 1 cm pour 500 FCFA en abscisses, 1 cm pour 2 élèves en ordonnées).
        \item Tracer le \cle{diagramme circulaire} (ou semi-circulaire) des fréquences.
    \end{itemize}
    \item \textbf{Prise de décision et rédaction :} En vous basant \emph{uniquement} sur les résultats chiffrés obtenus (classe modale, moyenne, lecture des graphiques), proposer au proviseur une \emph{tranche de prix} (intervalle) pour le ticket de transport scolaire subventionné. Justifier votre proposition par un argumentaire écrit d'une dizaine de lignes, en précisant si vous recommandez un montant unique ou une grille tarifaire par tranche de distance.
\end{enumerate}

\courssub{Ressources à mobiliser}
\begin{itemize}
    \item \textbf{Regroupement en classes :} Définir des classes d'amplitude constante $a = 500$. La première classe est $[800 ; 1300[$. Une valeur $x$ appartient à la classe $[a ; b[$ si $a \le x < b$. La dernière classe inclut sa borne supérieure.
    \item \textbf{Centre d'une classe :} $c = \frac{\text{borne inf} + \text{borne sup}}{2}$.
    \item \textbf{Fréquence d'une classe :} $f = \frac{\text{effectif de la classe}}{\text{effectif total}} \times 100$ (en \%). Somme des fréquences $= 100\%$.
    \item \textbf{Angle pour diagramme circulaire :} $\alpha = \frac{\text{effectif de la classe}}{\text{effectif total}} \times 360^\circ$. Somme des angles $= 360^\circ$.
    \item \textbf{Classe modale :} Classe de plus grand effectif.
    \item \textbf{Moyenne d'une série regroupée :} $\bar{x} = \frac{\sum (n_i \times c_i)}{\sum n_i}$ où $n_i$ est l'effectif et $c_i$ le centre de la classe $i$.
    \item \textbf{Histogramme :} Aire du rectangle = effectif de la classe. Hauteur (densité) $= \frac{\text{effectif}}{\text{amplitude}}$. Ici amplitude constante $\Rightarrow$ hauteur proportionnelle à l'effectif.
\end{itemize}

\begin{exoresolu}[Démarche et corrigé commenté]
\underline{\textbf{1. Construction du tableau de regroupement}}

L'amplitude des classes est $a = 500$. La première classe est $[800 ; 1300[$. Les classes successives sont :
$[800 ; 1300[$, $[1300 ; 1800[$, $[1800 ; 2300[$, $[2300 ; 2800[$.
L'effectif total est $N = 40$.

Comptage des effectifs (vérification : somme $= 40$) :
\begin{itemize}
    \item $[800 ; 1300[$ : valeurs $\ge 800$ et $< 1300$.
        Données : 800, 900, 950, 1000, 1050, 1100, 1100, 1150, 1200, 1200, 1250. $\rightarrow$ \textbf{11 élèves}.
    \item $[1300 ; 1800[$ : valeurs $\ge 1300$ et $< 1800$.
        Données : 1300, 1300, 1350, 1400, 1400, 1450, 1500, 1500, 1550, 1600, 1600, 1650, 1700, 1700, 1750. $\rightarrow$ \textbf{15 élèves}.
    \item $[1800 ; 2300[$ : valeurs $\ge 1800$ et $< 2300$.
        Données : 1800, 1800, 1850, 1900, 1900, 1950, 2000, 2050, 2100, 2200, 2200. $\rightarrow$ \textbf{11 élèves}.
    \item $[2300 ; 2800[$ : valeurs $\ge 2300$.
        Données : 2300, 2400, 2500. $\rightarrow$ \textbf{3 élèves}. (Note : 2300 appartient à cette classe car $2300 \ge 2300$).
\end{itemize}
Vérification : $11 + 15 + 11 + 3 = 40$. Correct.

Calcul des centres, fréquences et angles :
\begin{center}
\begin{tabular}{|c|c|c|c|c|}\hline
\textbf{Classes (FCFA)} & \textbf{Effectifs $n_i$} & \textbf{Centres $c_i$ (FCFA)} & \textbf{Fréquences $f_i$ (\%)} & \textbf{Angles $\alpha_i$ ($^\circ$)} \\\hline
$[800 ; 1300[$ & 11 & $\frac{800+1300}{2} = 1050$ & $\frac{11}{40}\times 100 = 27,5\%$ & $\frac{11}{40}\times 360 = 99^\circ$ \\\hline
$[1300 ; 1800[$ & 15 & $\frac{1300+1800}{2} = 1550$ & $\frac{15}{40}\times 100 = 37,5\%$ & $\frac{15}{40}\times 360 = 135^\circ$ \\\hline
$[1800 ; 2300[$ & 11 & $\frac{1800+2300}{2} = 2050$ & $\frac{11}{40}\times 100 = 27,5\%$ & $\frac{11}{40}\times 360 = 99^\circ$ \\\hline
$[2300 ; 2800[$ & 3 & $\frac{2300+2800}{2} = 2550$ & $\frac{3}{40}\times 100 = 7,5\%$ & $\frac{3}{40}\times 360 = 27^\circ$ \\\hline
\textbf{Total} & \textbf{40} & & \textbf{100\%} & \textbf{360$^\circ$} \\\hline
\end{tabular}
\end{center}

\underline{\textbf{2. Classe modale et interprétation}}
La classe modale est la classe de plus grand effectif : \textbf{$[1300 ; 1800[$} avec un effectif de 15 élèves.
Sa fréquence est de \textbf{37,5\%}.
\textit{Interprétation :} 37,5\% des élèves interrogés dépensent entre 1300 et 1800 FCFA par semaine pour leur transport. C'est la tranche de dépense la plus représentée dans l'échantillon, elle correspond au « profil type » de l'élève de cet établissement.

\underline{\textbf{3. Calcul de la moyenne}}
On utilise la formule $\bar{x} = \frac{\sum n_i c_i}{N}$.
\begin{align*}
\sum n_i c_i &= 11 \times 1050 + 15 \times 1550 + 11 \times 2050 + 3 \times 2550 \\
&= 11\,550 + 23\,250 + 22\,550 + 7\,650 \\
&= 65\,000
\end{align*}
\[
\bar{x} = \frac{65\,000}{40} = \textbf{1625 FCFA}
\]
\textit{Comparaison :} La moyenne (1625 FCFA) se situe \emph{à l'intérieur} de la classe modale $[1300 ; 1800[$. Elle est proche du centre de cette classe (1550 FCFA). Cela indique une relative homogénéité de la distribution autour de cette classe centrale (pas de forte asymétrie tirant la moyenne vers les valeurs extrêmes).

\underline{\textbf{4. Représentations graphiques}}

\textbf{Histogramme :}
Amplitude constante $= 500$ FCFA $\rightarrow$ 1 cm en abscisses.
Échelle en ordonnées : 1 cm pour 2 élèves $\Rightarrow$ hauteur en cm $= \frac{n_i}{2}$.
\begin{itemize}
    \item Classe $[800;1300[$ : largeur 1 cm, hauteur 5,5 cm (eff. 11).
    \item Classe $[1300;1800[$ : largeur 1 cm, hauteur 7,5 cm (eff. 15). $\leftarrow$ Plus grand rectangle.
    \item Classe $[1800;2300[$ : largeur 1 cm, hauteur 5,5 cm (eff. 11).
    \item Classe $[2300;2800[$ : largeur 1 cm, hauteur 1,5 cm (eff. 3).
\end{itemize}
\begin{popfigure}
\begin{tikzpicture}[scale=0.7]
% Axes
\draw[->, thick] (0,0) -- (5,0) node[right] {Dépense (FCFA)};
\draw[->, thick] (0,0) -- (0,9) node[above] {Effectif};
% Graduations abscisses (échelle: 1cm = 500 FCFA)
\foreach \x/\label in {0/800, 1/1300, 2/1800, 3/2300, 4/2800} {
    \draw (\x, -0.1) -- (\x, 0.1) node[below] {\label};
}
% Graduations ordonnées (1cm = 2 élèves)
\foreach \y/\label in {0/0, 1/2, 2/4, 3/6, 4/8, 5/10, 6/12, 7/14, 8/16} {
    \draw (-0.1, \y) -- (0.1, \y) node[left] {\label};
}
% Rectangles (Histogramme)
\fill[popblue!50] (0,0) rectangle (1,5.5);   % Eff 11 -> h=5.5
\fill[popblue!70] (1,0) rectangle (2,7.5);   % Eff 15 -> h=7.5
\fill[popblue!50] (2,0) rectangle (3,5.5);   % Eff 11 -> h=5.5
\fill[popblue!30] (3,0) rectangle (4,1.5);   % Eff 3  -> h=1.5
% Valeurs au-dessus des barres
\node at (0.5, 5.8) {11};
\node at (1.5, 7.8) {15};
\node at (2.5, 5.8) {11};
\node at (3.5, 1.8) {3};
\end{tikzpicture}
\legende{Histogramme des dépenses hebdomadaires de transport (40 élèves).}
\end{popfigure}

\textbf{Diagramme circulaire :}
Angles calculés : $99^\circ$, $135^\circ$, $99^\circ$, $27^\circ$.
Construction au rapporteur en partant du rayon horizontal vers la droite (sens anti-horaire).
\begin{popfigure}
\begin{tikzpicture}[scale=1.2]
% Camembert
\draw[fill=popblue!60] (0,0) -- (0:2.5) arc (0:99:2.5) -- cycle;
\node at (49.5:1.5) {\textbf{27,5\%}};
\draw[fill=poporange!70] (0,0) -- (99:2.5) arc (99:234:2.5) -- cycle;
\node at (166.5:1.5) {\textbf{37,5\%}};
\draw[fill=popgreen!60] (0,0) -- (234:2.5) arc (234:333:2.5) -- cycle;
\node at (283.5:1.5) {\textbf{27,5\%}};
\draw[fill=poppink!70] (0,0) -- (333:2.5) arc (333:360:2.5) -- cycle;
\node at (346.5:1.5) {\textbf{7,5\%}};
% Légende
\begin{scope}[xshift=4cm, yshift=1cm]
\node[anchor=west] at (0,1.2) {\color{popblue!60}\rule{0.5cm}{0.3cm} $[800;1300[$ : 27,5\%};
\node[anchor=west] at (0,0.6) {\color{poporange!70}\rule{0.5cm}{0.3cm} $[1300;1800[$ : 37,5\%};
\node[anchor=west] at (0,0) {\color{popgreen!60}\rule{0.5cm}{0.3cm} $[1800;2300[$ : 27,5\%};
\node[anchor=west] at (0,-0.6) {\color{poppink!70}\rule{0.5cm}{0.3cm} $[2300;2800[$ : 7,5\%};
\end{scope}
\end{tikzpicture}
\legende{Diagramme circulaire des fréquences des dépenses de transport.}
\end{popfigure}

\underline{\textbf{5. Proposition de tranche de prix pour le ticket subventionné — Rapport au Proviseur}}

\textbf{Objet :} Proposition de tarification pour le ticket de transport scolaire subventionné — Résultats de l'enquête 3ème A.

Monsieur le Proviseur,

À l'issue de l'enquête statistique réalisée auprès de 40 élèves de la classe de 3ème A, nous vous soumettons les résultats suivants pour éclairer la fixation du ticket de transport scolaire subventionné.

La dépense moyenne hebdomadaire s'élève à \textbf{1625 FCFA}. La classe modale, qui regroupe 37,5\% des élèves, est l'intervalle \textbf{[1300 ; 1800[ FCFA}. L'histogramme montre une concentration nette des effectifs dans cette classe centrale (15 élèves sur 40), tandis que les dépenses supérieures à 2300 FCFA ne concernent que 7,5\% de l'échantillon (3 élèves, probablement domiciliés en périphérie lointaine type Nkolbisson-Ngoa Ekellé ou Essos).

Au vu de ces éléments, nous proposons \textbf{de fixer le ticket subventionné à 1500 FCFA par semaine} (soit 6000 FCFA par mois), montant arrondi situé au centre de la classe modale (1550 FCFA) et proche de la moyenne observée (1625 FCFA).

Cette proposition présente trois avantages :
\begin{enumerate}
    \item \textbf{Équité :} Elle couvre la dépense réelle de la majorité (37,5\% des élèves paieraient exactement leur coût réel ou légèrement moins). Les 27,5\% d'élèves de la classe $[800;1300[$ (dépense $\le$ 1300) bénéficieraient d'un léger surplus (subvention inverse), acceptable socialement.
    \item \textbf{Simplicité de gestion :} Un montant unique évite la complexité administrative d'une grille tarifaire par distance (vérification des adresses, zonage), source de contentieux.
    \item \textbf{Soutenabilité :} Les 27,5\% d'élèves de la classe $[1800;2300[$ et les 7,5\% au-delà restent à charge partielle ; l'établissement peut négocier un forfait avec les transporteurs (Taxi-Moto, Carrefour, Bus Yango) sur la base de ce volume garanti.
\end{enumerate}

Nous recommandons une réévaluation annuelle de ce tarif à partir d'une nouvelle enquête élargie à tout le cycle.

Respectueusement, \\
Le groupe d'élèves statisticiens de 3ème A.
\end{exoresolu}

\courssub{Bilan de l'activité}
Cette activité a permis de mettre en évidence les points suivants, essentiels pour l'examen du BEPC :
\begin{itemize}
    \item \textbf{La nécessité du regroupement en classes} face à une liste de 40 valeurs brutes dispersées : sans classes, aucun tableau lisible, aucun graphique exploitable, aucune moyenne facile.
    \item \textbf{Le choix de l'amplitude des classes (500 FCFA)} n'est pas anodin : trop grande, elle masque la classe modale ; trop petite, elle crée des classes vides. Ici, 500 FCFA révèle parfaitement la concentration centrale.
    \item \textbf{La distinction entre classe modale et moyenne :} La classe modale ($[1300;1800[$) décrit la situation \emph{la plus fréquente} (le « cas type »), tandis que la moyenne (1625 FCFA) donne une \emph{valeur de synthèse} pour le budget global. Leur proximité ici valide la pertinence d'un tarif unique.
    \item \textbf{L'histogramme vs le diagramme circulaire :} L'histogramme (aires proportionnelles aux effectifs) permet de \emph{voir} la forme de la distribution (unimodale, symétrique). Le diagramme circulaire permet de \emph{comparer les parts relatives} (pourcentages) pour un public non expert (le proviseur).
    \item \textbf{L'argumentation chiffrée :} Une décision (le prix du ticket) ne se prend pas « au feeling ». Elle s'appuie sur des indicateurs (moyenne, mode, fréquences) et s'exprime par un texte structuré : constats $\rightarrow$ proposition $\rightarrow$ justifications.
\end{itemize}
\begin{attention}[Pièges évités lors de l'activité]
\begin{itemize}
\item Ne pas oublier que la borne supérieure 2800 de la dernière classe est \emph{incluse} (dernière classe fermée à droite), contrairement aux autres.
\item Pour la moyenne, utiliser les \textbf{centres des classes} ($c_i$), et non les bornes.
\item À l'histogramme : l'amplitude est constante, donc la hauteur des rectangles est proportionnelle aux effectifs. Si l'amplitude variait, il aurait fallu calculer la \emph{densité} (hauteur = effectif / amplitude).
\item Au diagramme circulaire : vérifier que la somme des angles fait exactement $360^\circ$ (arrondis obligatoires si nécessaire, ici sommes exactes).
\end{itemize}
\end{attention}

\newpage

\coursec{Méthodes et exercices résolus}

\courssub{Fiches méthodes et exercices résolus}

\begin{definition}[Classe, amplitude et centre d'une classe]
Lorsqu'une série statistique contient de nombreuses valeurs, on les \cle{regroupe en classes}, c'est-à-dire en intervalles de la forme $[a\,;\,b[$.
\begin{itemize}
\item L'\cle{amplitude} d'une classe $[a\,;\,b[$ est le nombre $b - a$.
\item Le \cle{centre} d'une classe $[a\,;\,b[$ est le nombre $c = \dfrac{a+b}{2}$.
\item La borne de gauche $a$ est \cle{incluse}, la borne de droite $b$ est \cle{exclue} : une valeur égale à $b$ appartient à la classe suivante.
\end{itemize}
\end{definition}

\begin{methode}[Regrouper une série en classes d'amplitude donnée]
Pour organiser une série statistique brute en classes :
\begin{enumerate}
\item Repérer la \cle{valeur minimale} et la \cle{valeur maximale} de la série, puis calculer l'\cle{étendue} : $e = x_{\max} - x_{\min}$.
\item Choisir l'\cle{amplitude} $a$ des classes (souvent imposée par l'énoncé) et déterminer le nombre de classes.
\item Écrire les classes sous la forme $[a\,;\,b[$ : la borne de gauche est \cle{incluse}, la borne de droite est \cle{exclue} (sauf éventuellement la dernière classe).
\item Dépouiller les données une par une (barrer chaque valeur comptée) et remplir la colonne des \cle{effectifs} $n_i$.
\item \textbf{Vérification obligatoire} : la somme des effectifs doit être égale à l'effectif total $N$.
\end{enumerate}
\end{methode}

\begin{exoresolu}[Regroupement en classes]
Voici les notes sur 20 obtenues par 30 élèves de 3ème à un devoir de mathématiques :
\[8\;;\;12\;;\;7\;;\;15\;;\;10\;;\;9\;;\;13\;;\;11\;;\;6\;;\;14\;;\;10\;;\;12\;;\;8\;;\;16\;;\;9\]
\[11\;;\;13\;;\;7\;;\;10\;;\;12\;;\;15\;;\;9\;;\;14\;;\;8\;;\;11\;;\;13\;;\;10\;;\;12\;;\;6\;;\;17\]
Regrouper ces notes en classes d'amplitude 3, en commençant par la classe $[5\,;\,8[$.
\tcblower
\textbf{Étape 1 : étendue.} La note minimale est $6$, la note maximale est $17$. L'étendue est $e = 17 - 6 = 11$.

\textbf{Étape 2 : classes d'amplitude 3.} En partant de 5 : $[5\,;\,8[$, $[8\,;\,11[$, $[11\,;\,14[$, $[14\,;\,17[$, $[17\,;\,20[$.

\textbf{Étape 3 : dépouillement.} On compte les notes dans chaque classe (attention : la note $8$ appartient à $[8\,;\,11[$ et non à $[5\,;\,8[$, car 8 est exclu de la première classe) :
\begin{itemize}
\item $[5\,;\,8[$ : notes 7, 7, 6, 6 $\Rightarrow$ effectif $4$ ;
\item $[8\,;\,11[$ : notes 8, 10, 9, 10, 8, 9, 10, 9, 8, 10 $\Rightarrow$ effectif $10$ ;
\item $[11\,;\,14[$ : notes 12, 13, 11, 12, 11, 13, 12, 11, 13, 12 $\Rightarrow$ effectif $10$ ;
\item $[14\,;\,17[$ : notes 15, 14, 16, 15, 14 $\Rightarrow$ effectif $5$ ;
\item $[17\,;\,20[$ : note 17 $\Rightarrow$ effectif $1$.
\end{itemize}

\textbf{Étape 4 : tableau.}
\begin{center}
\begin{tabularx}{\linewidth}{|l|X|X|X|X|X|}\hline
\rowcolor{popblueL} Classe & $[5;8[$ & $[8;11[$ & $[11;14[$ & $[14;17[$ & $[17;20[$ \\ \hline
Effectif & 4 & 10 & 10 & 5 & 1 \\ \hline
\end{tabularx}
\end{center}

\textbf{Vérification} : $4 + 10 + 10 + 5 + 1 = 30 = N$. Le dépouillement est correct.
\end{exoresolu}

\begin{definition}[Classe modale et fréquence d'une classe]
\begin{itemize}
\item La \cle{classe modale} d'une série regroupée en classes de même amplitude est la classe dont l'\textbf{effectif est le plus grand}.
\item La \cle{fréquence} d'une classe d'effectif $n_i$ dans une série d'effectif total $N$ est le nombre :
\[f_i = \frac{n_i}{N}.\]
Elle peut s'exprimer en fraction, en nombre décimal ou en pourcentage ($f_i \times 100\,\%$). La somme de toutes les fréquences est toujours égale à $1$ (soit $100\,\%$).
\end{itemize}
\end{definition}

\begin{methode}[Déterminer la classe modale et calculer une fréquence]
\begin{enumerate}
\item La \cle{classe modale} est la classe dont l'\textbf{effectif est le plus grand} (pour des classes de même amplitude).
\item La \cle{fréquence} d'une classe se calcule par : $f_i = \dfrac{n_i}{N}$ où $n_i$ est l'effectif de la classe et $N$ l'effectif total.
\item La fréquence peut s'exprimer en \textbf{fraction}, en \textbf{nombre décimal} ou en \textbf{pourcentage} : $f_i \times 100\,\%$.
\item \textbf{Contrôle} : la somme de toutes les fréquences doit être égale à $1$ (ou $100\,\%$).
\end{enumerate}
\end{methode}

\begin{exoresolu}[Classe modale et fréquences]
Une coopérative de Njombé a pesé 40 sacs de cacao. Les masses (en kg) sont regroupées ainsi :
\begin{center}
\begin{tabularx}{\linewidth}{|l|X|X|X|X|}\hline
\rowcolor{popblueL} Masse (kg) & $[50;60[$ & $[60;70[$ & $[70;80[$ & $[80;90[$ \\ \hline
Effectif & 6 & 14 & 12 & 8 \\ \hline
\end{tabularx}
\end{center}
\begin{enumerate}
\item Déterminer la classe modale.
\item Calculer la fréquence de chaque classe en pourcentage.
\end{enumerate}
\tcblower
\textbf{1. Classe modale.} L'effectif le plus grand est $14$, correspondant à la classe $[60\,;\,70[$. La classe modale est donc $\mathbf{[60\,;\,70[}$ : la plupart des sacs pèsent entre 60 kg et 70 kg.

\textbf{2. Fréquences.} L'effectif total est $N = 6 + 14 + 12 + 8 = 40$.
\begin{align*}
f_1 &= \frac{6}{40} = 0{,}15 = 15\,\% &
f_2 &= \frac{14}{40} = 0{,}35 = 35\,\% \\
f_3 &= \frac{12}{40} = 0{,}30 = 30\,\% &
f_4 &= \frac{8}{40} = 0{,}20 = 20\,\%
\end{align*}
\textbf{Vérification} : $15 + 35 + 30 + 20 = 100\,\%$. Les calculs sont cohérents.

\textbf{Conclusion} : la classe modale est $[60\,;\,70[$ et les fréquences sont $15\,\%$, $35\,\%$, $30\,\%$ et $20\,\%$.
\end{exoresolu}

\begin{methode}[Calculer la moyenne d'une série regroupée en classes]
Quand les données sont regroupées en classes, on ne connaît plus les valeurs exactes : on remplace chaque classe par son \cle{centre}.
\begin{enumerate}
\item Calculer le \cle{centre} $c_i$ de chaque classe : $c_i = \dfrac{\text{borne inférieure} + \text{borne supérieure}}{2}$.
\item Calculer chaque produit $n_i \times c_i$ (effectif $\times$ centre).
\item Additionner tous les produits : $\sum n_i c_i$.
\item Diviser par l'effectif total : $\bar{x} = \dfrac{\sum n_i c_i}{N}$.
\item Conclure par une phrase avec l'unité, en précisant qu'il s'agit d'une moyenne \emph{approchée}.
\end{enumerate}
\end{methode}

\begin{exoresolu}[Moyenne d'une série regroupée en classes]
On a relevé les distances parcourues (en km) par 50 motos-taxis de Bafoussam en une journée :
\begin{center}
\begin{tabularx}{\linewidth}{|l|X|X|X|X|}\hline
\rowcolor{popblueL} Distance (km) & $[20;40[$ & $[40;60[$ & $[60;80[$ & $[80;100[$ \\ \hline
Effectif & 8 & 18 & 16 & 8 \\ \hline
\end{tabularx}
\end{center}
Calculer la distance moyenne parcourue par une moto-taxi.
\tcblower
\textbf{Étape 1 : centres des classes.}
\[c_1 = \frac{20+40}{2} = 30 \;;\quad c_2 = \frac{40+60}{2} = 50 \;;\quad c_3 = \frac{60+80}{2} = 70 \;;\quad c_4 = \frac{80+100}{2} = 90.\]

\textbf{Étape 2 : produits $n_i c_i$.}
\begin{center}
\begin{tabularx}{\linewidth}{|l|X|X|X|X|X|}\hline
\rowcolor{popblueL} Classe & $[20;40[$ & $[40;60[$ & $[60;80[$ & $[80;100[$ & Total \\ \hline
Effectif $n_i$ & 8 & 18 & 16 & 8 & 50 \\ \hline
Centre $c_i$ & 30 & 50 & 70 & 90 & \\ \hline
$n_i c_i$ & 240 & 900 & 1120 & 720 & 2980 \\ \hline
\end{tabularx}
\end{center}

\textbf{Étape 3 : moyenne.}
\[\bar{x} = \frac{\sum n_i c_i}{N} = \frac{2980}{50} = 59{,}6.\]

\textbf{Conclusion} : chaque moto-taxi parcourt en moyenne environ $\mathbf{59{,}6}$ \textbf{km} par jour. (C'est une valeur approchée car on a remplacé chaque classe par son centre.)
\end{exoresolu}

\begin{methode}[Construire un histogramme]
Pour représenter graphiquement une série regroupée en classes de \textbf{même amplitude} :
\begin{enumerate}
\item Tracer deux axes perpendiculaires : en \textbf{abscisse} les classes (bornes reportées à intervalles réguliers), en \textbf{ordonnée} les effectifs.
\item Pour chaque classe, tracer un \cle{rectangle} dont la base est l'amplitude de la classe et la hauteur l'effectif.
\item Les rectangles sont \textbf{juxtaposés} (collés les uns aux autres) : il n'y a pas d'espace entre eux.
\item Choisir des échelles adaptées et les indiquer clairement sur les axes ; donner un titre au graphique.
\end{enumerate}
\textbf{Réflexe} : vérifier que la somme des hauteurs (effectifs) redonne bien $N$.
\end{methode}

\begin{exoresolu}[Histogramme]
Les âges des 35 élèves d'une classe de 3ème sont regroupés ainsi :
\begin{center}
\begin{tabularx}{\linewidth}{|l|X|X|X|X|}\hline
\rowcolor{popblueL} Âge (ans) & $[13;14[$ & $[14;15[$ & $[15;16[$ & $[16;17[$ \\ \hline
Effectif & 5 & 14 & 11 & 5 \\ \hline
\end{tabularx}
\end{center}
Construire l'histogramme de cette série.
\tcblower
\textbf{Étape 1 : vérification.} $5 + 14 + 11 + 5 = 35 = N$. Les classes ont toutes la même amplitude (1 an) : on peut construire un histogramme avec les effectifs en hauteur.

\textbf{Étape 2 : construction.} En abscisse, on reporte les bornes 13, 14, 15, 16, 17 ; en ordonnée, les effectifs de 0 à 14. Chaque rectangle a pour base 1 an et pour hauteur l'effectif de la classe.

\begin{popfigure}
\begin{center}
\begin{tikzpicture}[x=1.6cm,y=0.28cm]
\draw[->,thick] (12.6,0) -- (17.4,0) node[below left]{Âge (ans)};
\draw[->,thick] (13,0) -- (13,16) node[above left]{Effectif};
\foreach \x in {13,14,15,16,17} \draw (\x,0.2) -- (\x,-0.2) node[below]{\x};
\foreach \y in {5,10,14} \draw (12.95,\y) -- (13.05,\y) node[left]{\y};
\fill[popblue!50,draw=popblue,thick] (13,0) rectangle (14,5);
\fill[poporange!50,draw=poporange,thick] (14,0) rectangle (15,14);
\fill[popgreen!50,draw=popgreen,thick] (15,0) rectangle (16,11);
\fill[poppurple!50,draw=poppurple,thick] (16,0) rectangle (17,5);
\node at (13.5,5)[above]{5};
\node at (14.5,14)[above]{14};
\node at (15.5,11)[above]{11};
\node at (16.5,5)[above]{5};
\end{tikzpicture}
\end{center}
\legende{Histogramme des âges des élèves de la classe}
\end{popfigure}

\textbf{Conclusion} : les rectangles sont juxtaposés, de même largeur, et leurs hauteurs sont 5, 14, 11 et 5. On lit immédiatement que la classe modale est $[14\,;\,15[$.
\end{exoresolu}

\begin{methode}[Construire un diagramme circulaire ou semi-circulaire]
\begin{enumerate}
\item Calculer la \cle{fréquence} de chaque classe : $f_i = \dfrac{n_i}{N}$.
\item Calculer l'\cle{angle} de chaque secteur :
\begin{itemize}
\item diagramme \textbf{circulaire} : $\alpha_i = f_i \times 360°$ ;
\item diagramme \textbf{semi-circulaire} : $\alpha_i = f_i \times 180°$.
\end{itemize}
\item \textbf{Vérifier} que la somme des angles vaut $360°$ (ou $180°$).
\item Tracer un cercle (ou un demi-cercle), puis reporter chaque angle au rapporteur en partant du rayon précédent.
\item Colorier les secteurs et ajouter une \textbf{légende} et un \textbf{titre}.
\end{enumerate}
\end{methode}

\begin{exoresolu}[Diagramme semi-circulaire]
Sur 180 élèves d'un collège de Yaoundé, on a relevé le moyen de transport utilisé pour venir en classe :
\begin{center}
\begin{tabularx}{\linewidth}{|l|X|X|X|X|}\hline
\rowcolor{popblueL} Transport & À pied & Moto-taxi & Taxi & Bus \\ \hline
Effectif & 72 & 54 & 36 & 18 \\ \hline
\end{tabularx}
\end{center}
Construire un diagramme semi-circulaire de cette répartition.
\tcblower
\textbf{Étape 1 : fréquences et angles.} $N = 72 + 54 + 36 + 18 = 180$. Pour un diagramme semi-circulaire, $\alpha_i = \dfrac{n_i}{N} \times 180°$.
\begin{align*}
\text{À pied : } & \frac{72}{180} \times 180° = 72° &
\text{Moto-taxi : } & \frac{54}{180} \times 180° = 54° \\
\text{Taxi : } & \frac{36}{180} \times 180° = 36° &
\text{Bus : } & \frac{18}{180} \times 180° = 18°
\end{align*}
\textbf{Vérification} : $72° + 54° + 36° + 18° = 180°$. Correct.

\textbf{Étape 2 : construction.} On trace un demi-cercle, puis on reporte successivement au rapporteur les angles $72°$, $54°$, $36°$ et $18°$.

\begin{popfigure}
\begin{center}
\begin{tikzpicture}[scale=1.4]
\fill[popblue!50] (0:2.6) arc (0:72:2.6) -- (0,0) -- cycle;
\fill[poporange!50] (72:2.6) arc (72:126:2.6) -- (0,0) -- cycle;
\fill[popgreen!50] (126:2.6) arc (126:162:2.6) -- (0,0) -- cycle;
\fill[poppurple!50] (162:2.6) arc (162:180:2.6) -- (0,0) -- cycle;
\draw[thick] (0:2.6) arc (0:180:2.6) -- cycle;
\draw[thick] (0,0) -- (0:2.6); \draw[thick] (0,0) -- (72:2.6);
\draw[thick] (0,0) -- (126:2.6); \draw[thick] (0,0) -- (162:2.6);
\draw[thick] (0,0) -- (180:2.6);
\node at (36:1.7){À pied $72°$};
\node at (99:1.8){Moto $54°$};
\node at (144:1.85){Taxi $36°$};
\node at (171:1.95){Bus};
\end{tikzpicture}
\end{center}
\legende{Diagramme semi-circulaire des moyens de transport}
\end{popfigure}

\textbf{Conclusion} : le secteur « À pied » ($72°$) est le plus grand : c'est le moyen de transport majoritaire dans ce collège.
\end{exoresolu}

\begin{methode}[Résoudre un problème complet de statistiques (type BEPC)]
Face à un problème de synthèse :
\begin{enumerate}
\item \textbf{Lire} attentivement le tableau ou l'énoncé ; identifier l'effectif total $N$.
\item \textbf{Compléter} le tableau si nécessaire (effectifs, fréquences, centres, angles) en utilisant les relations $f_i = \dfrac{n_i}{N}$, $n_i = f_i \times N$, et le fait que $\sum n_i = N$.
\item \textbf{Calculer} les indicateurs demandés : classe modale, moyenne (avec les centres), fréquence.
\item \textbf{Représenter} graphiquement si demandé (histogramme ou diagramme circulaire).
\item \textbf{Rédiger} chaque réponse par une phrase complète, avec l'unité, et \textbf{vérifier} : somme des effectifs $= N$, somme des fréquences $= 1$, somme des angles $= 360°$.
\end{enumerate}
\end{methode}

\begin{exoresolu}[Problème de synthèse type BEPC]
Un commerçant du marché central de Douala a noté le montant des achats (en milliers de FCFA) de 60 clients un samedi :
\begin{center}
\begin{tabularx}{\linewidth}{|l|X|X|X|X|}\hline
\rowcolor{popblueL} Montant (milliers FCFA) & $[0;5[$ & $[5;10[$ & $[10;15[$ & $[15;20[$ \\ \hline
Effectif & 12 & 24 & $a$ & 9 \\ \hline
Fréquence (\%) & & & & \\ \hline
\end{tabularx}
\end{center}
\begin{enumerate}
\item Déterminer l'effectif $a$ manquant.
\item Compléter la ligne des fréquences en pourcentage.
\item Quelle est la classe modale ?
\item Calculer le montant moyen des achats d'un client.
\end{enumerate}
\tcblower
\textbf{1. Effectif manquant.} La somme des effectifs est $N = 60$ :
\[12 + 24 + a + 9 = 60 \iff a = 60 - 45 = 15.\]
Donc $\mathbf{a = 15}$.

\textbf{2. Fréquences.} On applique $f_i = \dfrac{n_i}{60} \times 100$ :
\begin{align*}
f_1 &= \frac{12}{60}\times 100 = 20\,\% &
f_2 &= \frac{24}{60}\times 100 = 40\,\% \\
f_3 &= \frac{15}{60}\times 100 = 25\,\% &
f_4 &= \frac{9}{60}\times 100 = 15\,\%
\end{align*}
Vérification : $20 + 40 + 25 + 15 = 100\,\%$. Correct.

\textbf{3. Classe modale.} L'effectif le plus grand est $24$ : la classe modale est $\mathbf{[5\,;\,10[}$ (en milliers de FCFA). La majorité des clients dépensent entre 5 000 et 10 000 FCFA.

\textbf{4. Montant moyen.} Centres des classes : $c_1 = 2{,}5$ ; $c_2 = 7{,}5$ ; $c_3 = 12{,}5$ ; $c_4 = 17{,}5$.
\[\sum n_i c_i = 12 \times 2{,}5 + 24 \times 7{,}5 + 15 \times 12{,}5 + 9 \times 17{,}5 = 30 + 180 + 187{,}5 + 157{,}5 = 555.\]
\[\bar{x} = \frac{555}{60} = 9{,}25.\]

\textbf{Conclusion} : le montant moyen des achats d'un client est d'environ $\mathbf{9{,}25}$ \textbf{milliers de FCFA}, soit \textbf{9 250 FCFA}.
\end{exoresolu}

\begin{savaistu}[Les statistiques au quotidien au Cameroun]
Les statistiques sont partout autour de toi : l'Institut National de la Statistique (INS) recense la population camerounaise et publie des tableaux regroupés en classes d'âges ; les opérateurs MTN et Orange analysent la consommation de forfaits par tranches ; les résultats du BEPC sont présentés par classes de notes. Savoir lire un histogramme ou un diagramme circulaire, c'est savoir lire l'information qui circule chaque jour dans les journaux et à la télévision.
\end{savaistu}

\begin{attention}[Les 5 erreurs qui coûtent des points au BEPC]
\begin{enumerate}
\item \textbf{Confondre classe modale et effectif.} La classe modale est une \emph{classe} (un intervalle, par exemple $[60\,;\,70[$), pas un effectif ni un centre. Écrire « la classe modale est 14 » est faux : 14 est l'effectif de la classe modale.
\item \textbf{Calculer la moyenne avec les bornes au lieu des centres.} Pour une série regroupée en classes, la moyenne se calcule obligatoirement avec les \emph{centres} des classes : $\bar{x} = \dfrac{\sum n_i c_i}{N}$. Utiliser les bornes inférieures ou supérieures fausse tout le calcul.
\item \textbf{Oublier de vérifier les totaux.} Somme des effectifs $= N$, somme des fréquences $= 1$ (ou $100\,\%$), somme des angles $= 360°$ (ou $180°$ pour un semi-circulaire). Ces vérifications, écrites sur la copie, montrent la maîtrise et évitent les erreurs de dépouillement.
\item \textbf{Espacer les rectangles de l'histogramme.} Dans un histogramme, les rectangles sont \emph{juxtaposés} (collés) car le caractère est continu. Des rectangles séparés correspondent à un diagramme en bâtons (caractère discret) : ne pas confondre.
\item \textbf{Mal rédiger les classes et les conclusions.} Une classe s'écrit $[a\,;\,b[$ avec la borne de droite \emph{exclue} : une valeur égale à $b$ appartient à la classe suivante. Et toute réponse doit se terminer par une phrase avec l'unité (FCFA, km, kg, ans...), sans quoi le correcteur retire des points de présentation.
\end{enumerate}
\end{attention}

\newpage

\coursec{Exercices}

\courssub{Je vérifie mes connaissances}

\begin{exercice}[QCM : Vocabulaire et formules de base]
Parmi les affirmations suivantes, une seule est exacte. Laquelle ? Justifier brièvement.
\begin{enumerate}
    \item Dans une série statistique regroupée en classes, la fréquence d'une classe est égale à l'effectif de cette classe divisé par l'effectif total.
    \item La classe modale est la classe dont la borne inférieure est la plus petite.
    \item Pour calculer la moyenne d'une série regroupée en classes, on utilise les bornes inférieures des classes.
    \item La somme des angles d'un diagramme circulaire est toujours égale à $180^\circ$.
\end{enumerate}
\end{exercice}

\begin{exercice}[Vrai ou Faux ? Justifier]
Dire si les affirmations suivantes sont vraies ou fausses. Justifier chaque réponse par un calcul ou une propriété du cours.
\begin{enumerate}
    \item Dans un histogramme représentant une série regroupée en classes de même amplitude, la hauteur de chaque rectangle est proportionnelle à l'effectif de la classe.
    \item Si une série statistique comporte $30$ valeurs, la somme des fréquences (en pourcentage) est égale à $100$.
    \item La moyenne d'une série regroupée en classes est toujours égale à l'une des valeurs de la série.
    \item Le polygone des effectifs se construit en reliant les milieux des sommets des rectangles de l'histogramme.
\end{enumerate}
\end{exercice}

\begin{exercice}[Définitions et notations]
Compléter les phrases suivantes avec les termes mathématiques appropriés.
\begin{enumerate}
    \item L'intervalle $[a ; b[$ qui regroupe des valeurs voisines s'appelle une \ldots
    \item Le nombre de valeurs appartenant à une classe est l'\ldots de cette classe.
    \item Le quotient de l'effectif d'une classe par l'effectif total est la \ldots de cette classe.
    \item La classe qui possède le plus grand effectif est la \ldots
    \item Pour calculer la moyenne d'une série regroupée en classes, on remplace chaque valeur de la classe par le \ldots de cette classe.
\end{enumerate}
\end{exercice}

\begin{exercice}[Calculs directs sur un petit effectif]
On a relevé les âges (en années) de $10$ élèves d'un club de lecture : 
\[ 12;\ 13;\ 13;\ 14;\ 14;\ 14;\ 15;\ 15;\ 16;\ 17 \]
On regroupe cette série en classes d'amplitude $2$ : $[12 ; 14[$, $[14 ; 16[$, $[16 ; 18[$.
\begin{enumerate}
    \item Construire le tableau des effectifs et des fréquences (en \%) de cette série regroupée.
    \item Déterminer la classe modale.
    \item Calculer la moyenne de cette série regroupée. Arrondir au centième.
\end{enumerate}
\end{exercice}

\courssub{Je m'entraîne}

\begin{exercice}[Tailles d'élèves : tableau, classe modale et moyenne]
Dans un collège de Yaoundé, on a mesuré la taille (en cm) de $50$ élèves de 3ème. Les résultats sont regroupés dans le tableau ci-dessous :
\begin{center}
\begin{tabular}{|c|c|c|c|c|}\hline
Classes (en cm) & $[150 ; 160[$ & $[160 ; 170[$ & $[170 ; 180[$ & $[180 ; 190[$ \\ \hline
Effectifs & $8$ & $18$ & $16$ & $8$ \\ \hline
\end{tabular}
\end{center}
\begin{enumerate}
    \item Compléter le tableau en ajoutant la ligne des fréquences (en \%) et la ligne des valeurs caractéristiques (milieux des classes).
    \item Quelle est la classe modale ? Interpréter ce résultat dans le contexte.
    \item Calculer la taille moyenne des élèves de cet échantillon. Arrondir au dixième de cm.
\end{enumerate}
\end{exercice}

\begin{exercice}[Tableau incomplet et diagramme circulaire]
Un enquêteur a interrogé $200$ familles du quartier Mvog-Ada sur leur nombre d'enfants scolarisés. Le tableau ci-dessous présente les résultats, mais certaines cases sont manquantes.
\begin{center}
\begin{tabular}{|c|c|c|c|c|c|}\hline
Nombre d'enfants & $0$ & $1$ & $2$ & $3$ & $4$ et plus \\ \hline
Effectifs & $40$ & $60$ & \ldots & $30$ & $20$ \\ \hline
Fréquences (\%) & \ldots & $30\%$ & $25\%$ & \ldots & \ldots \\ \hline
Angles ($^\circ$) & $72^\circ$ & \ldots & \ldots & $54^\circ$ & \ldots \\ \hline
\end{tabular}
\end{center}
\begin{enumerate}
    \item Compléter intégralement le tableau.
    \item Construire le diagramme circulaire correspondant à cette série statistique. (Indiquer la méthode de construction avec le rapporteur).
    \item Quel pourcentage de familles a au moins $2$ enfants scolarisés ?
\end{enumerate}
\end{exercice}

\begin{exercice}[Notes de mathématiques : moyenne et pourcentage]
Les notes sur $20$ obtenues par les élèves d'une classe de 3ème à un devoir de mathématiques sont regroupées ainsi :
\begin{center}
\begin{tabular}{|c|c|c|c|c|c|}\hline
Classes & $[0 ; 4[$ & $[4 ; 8[$ & $[8 ; 12[$ & $[12 ; 16[$ & $[16 ; 20]$ \\ \hline
Effectifs & $3$ & $7$ & $12$ & $10$ & $3$ \\ \hline
\end{tabular}
\end{center}
\begin{enumerate}
    \item Calculer la moyenne de la classe. Arrondir au centième.
    \item Calculer le pourcentage d'élèves ayant obtenu une note supérieure ou égale à $8$. Arrondir au dixième de pourcent.
    \item Le professeur dit : « Plus de la moitié de la classe a au moins la moyenne ($10/20$) ». Vérifier cette affirmation à l'aide du tableau.
\end{enumerate}
\end{exercice}

\begin{exercice}[Histogramme et polygone des effectifs : distances domicile-collège]
Une enquête auprès de $60$ élèves d'un lycée de Douala donne les distances (en km) entre leur domicile et l'établissement :
\begin{center}
\begin{tabular}{|c|c|c|c|c|c|}\hline
Classes (km) & $[0 ; 2[$ & $[2 ; 4[$ & $[4 ; 6[$ & $[6 ; 8[$ & $[8 ; 10[$ \\ \hline
Effectifs & $15$ & $20$ & $12$ & $8$ & $5$ \\ \hline
\end{tabular}
\end{center}
\begin{enumerate}
    \item Construire l'histogramme de cette série statistique (choisir une unité de surface : 1 cm$^2$ pour 1 élève).
    \item Sur le même repère, tracer le polygone des effectifs.
    \item Identifier graphiquement la classe modale. Vérifier par le calcul.
    \item Quelle proportion d'élèves habite à moins de $4$ km du lycée ?
\end{enumerate}
\end{exercice}

\courssub{Je me prépare au Probatoire}

\begin{exercice}[Gestion de stock : ventes d'un commerçant au Marché Central]
Monsieur Kamga tient une boutique de vêtements au Marché Central de Yaoundé. Il a relevé ses recettes journalières (en milliers de FCFA) sur un mois de $30$ jours ouvrés :
\begin{center}
\begin{tabular}{|c|c|c|c|c|c|}\hline
Classes (kFCFA) & $[50 ; 70[$ & $[70 ; 90[$ & $[90 ; 110[$ & $[110 ; 130[$ & $[130 ; 150[$ \\ \hline
Effectifs (jours) & $4$ & $8$ & $10$ & $5$ & $3$ \\ \hline
\end{tabular}
\end{center}
\begin{enumerate}
    \item Calculer la recette moyenne journalière de Monsieur Kamga. Arrondir au millier de FCFA.
    \item Déterminer la classe modale. Que représente-t-elle concrètement pour le commerçant ?
    \item Monsieur Kamga souhaite commander du stock pour le mois prochain. Il estime qu'il doit avoir un stock correspondant à $1,2$ fois sa recette moyenne mensuelle prévisionnelle. Quelle valeur de stock (en FCFA) doit-il prévoir ? (On suppose $30$ jours ouvrés).
    \item Rédiger une conclusion argumentée pour l'aider à prendre sa décision de commande, en tenant compte de la dispersion des ventes (sans calcul d'écart-type, juste observation du tableau).
\end{enumerate}
\end{exercice}

\begin{exercice}[Lecture critique : Erreur dans un journal]
Le journal « Cameroon Tribune » a publié un diagramme circulaire représentant la répartition des abonnés mobile par opérateur dans une région (effectif total : $10\,000$ abonnés). Les angles indiqués sur le diagramme sont :
\begin{itemize}
    \item MTN : $130^\circ$
    \item Orange : $110^\circ$
    \item Nexttel : $80^\circ$
    \item Camtel : $50^\circ$
\end{itemize}
\begin{enumerate}
    \item Vérifier par le calcul si la somme des angles est correcte pour un diagramme circulaire.
    \item En supposant que l'angle de MTN est exact, recalculer les effectifs pour chaque opérateur.
    \item Déduire l'erreur commise par le graphiste du journal sur les angles. Proposer les angles corrects pour Orange, Nexttel et Camtel.
    \item Pourquoi est-il important de vérifier la somme des angles avant de publier un diagramme circulaire ?
\end{enumerate}
\end{exercice}

\begin{exercice}[Comparaison de deux classes : 3ème A vs 3ème B]
Deux classes de 3ème, la 3ème A ($35$ élèves) et la 3ème B ($30$ élèves), ont passé le même devoir de mathématiques. Les notes sont regroupées en classes :
\begin{center}
\begin{tabular}{|c|c|c|c|c|c|}\hline
Classes & $[0 ; 5[$ & $[5 ; 10[$ & $[10 ; 15[$ & $[15 ; 20]$ & Total \\ \hline
Effectifs 3ème A & $3$ & $10$ & $15$ & $7$ & $35$ \\ \hline
Effectifs 3ème B & $2$ & $8$ & $12$ & $8$ & $30$ \\ \hline
\end{tabular}
\end{center}
\begin{enumerate}
    \item Calculer la moyenne de chaque classe. Arrondir au centième.
    \item Déterminer la classe modale de chaque série.
    \item Calculer le pourcentage de réussite (note $\ge 10$) dans chaque classe.
    \item Rédiger une comparaison justifiée du niveau des deux classes en utilisant l'ensemble des résultats obtenus (moyenne, classe modale, taux de réussite). Quelle classe semble avoir un niveau globalement meilleur ? Argumenter.
\end{enumerate}
\end{exercice}

\newpage

\coursec{Corrigés des exercices}

\coursec{Corrigés des exercices d'application — Statistiques}

\begin{corrige}[Exercice 1 — QCM : Vocabulaire et formules de base]
La seule affirmation exacte est la \textbf{1)}.

\begin{enumerate}
    \item \textbf{Vraie.} Par définition, la fréquence d'une classe est $f = \dfrac{\text{effectif de la classe}}{\text{effectif total}}$.
    \item \textbf{Fausse.} La classe modale est la classe qui possède le \cle{plus grand effectif} (ou la plus grande fréquence), et non celle dont la borne inférieure est la plus petite.
    \item \textbf{Fausse.} Pour calculer la moyenne d'une série regroupée en classes, on utilise les \cle{centres} (milieux) des classes, pas les bornes inférieures.
    \item \textbf{Fausse.} La somme des angles d'un diagramme circulaire est égale à $360^\circ$ (c'est $180^\circ$ pour un diagramme semi-circulaire).
\end{enumerate}
\end{corrige}

\begin{corrige}[Exercice 2 — Vrai ou Faux ? Justifier]
\begin{enumerate}
    \item \textbf{Vrai.} Dans un histogramme, l'aire de chaque rectangle est proportionnelle à l'effectif. Si toutes les classes ont la même amplitude (même largeur de base), alors la hauteur est proportionnelle à l'effectif, puisque aire $=$ base $\times$ hauteur avec une base constante.
    \item \textbf{Vrai.} La somme des fréquences de toutes les classes est toujours égale à $1$, soit $100\%$, quel que soit l'effectif total ($30$ ou autre). En effet : $\sum f_i = \dfrac{\sum n_i}{N} = \dfrac{N}{N} = 1$.
    \item \textbf{Faux.} La moyenne est une valeur calculée qui n'appartient pas nécessairement à la série. Contre-exemple : la série $1\,; 2$ a pour moyenne $1,5$, qui n'est pas une valeur de la série.
    \item \textbf{Vrai.} Le polygone des effectifs s'obtient en joignant par des segments les milieux des côtés supérieurs (sommets) des rectangles de l'histogramme. Ces milieux correspondent aux centres des classes, munis de leurs effectifs.
\end{enumerate}
\end{corrige}

\begin{corrige}[Exercice 3 — Définitions et notations]
\begin{enumerate}
    \item L'intervalle $[a ; b[$ qui regroupe des valeurs voisines s'appelle une \cle{classe}.
    \item Le nombre de valeurs appartenant à une classe est l'\cle{effectif} de cette classe.
    \item Le quotient de l'effectif d'une classe par l'effectif total est la \cle{fréquence} de cette classe.
    \item La classe qui possède le plus grand effectif est la \cle{classe modale}.
    \item Pour calculer la moyenne d'une série regroupée en classes, on remplace chaque valeur de la classe par le \cle{centre} (ou milieu) de cette classe.
\end{enumerate}
\end{corrige}

\begin{corrige}[Exercice 4 — Calculs directs sur un petit effectif]
\begin{enumerate}
    \item Comptons les valeurs dans chaque classe :
    \begin{itemize}
        \item $[12 ; 14[$ : $12\,; 13\,; 13$ soit $3$ valeurs ;
        \item $[14 ; 16[$ : $14\,; 14\,; 14\,; 15\,; 15$ soit $5$ valeurs ;
        \item $[16 ; 18[$ : $16\,; 17$ soit $2$ valeurs.
    \end{itemize}
    Vérification : $3 + 5 + 2 = 10$. Les fréquences : $\dfrac{3}{10} = 30\%$ ; $\dfrac{5}{10} = 50\%$ ; $\dfrac{2}{10} = 20\%$.
    \begin{center}
    \begin{tabular}{|c|c|c|c|}\hline
    Classes & $[12 ; 14[$ & $[14 ; 16[$ & $[16 ; 18[$ \\ \hline
    Effectifs & $3$ & $5$ & $2$ \\ \hline
    Fréquences (\%) & $30$ & $50$ & $20$ \\ \hline
    \end{tabular}
    \end{center}
    \item La classe modale est $[14 ; 16[$ car c'est elle qui a le plus grand effectif ($5$).
    \item Les centres des classes sont : $\dfrac{12+14}{2} = 13$ ; $\dfrac{14+16}{2} = 15$ ; $\dfrac{16+18}{2} = 17$.
    \[ \bar{x} = \dfrac{3 \times 13 + 5 \times 15 + 2 \times 17}{10} = \dfrac{39 + 75 + 34}{10} = \dfrac{148}{10} = 14,8. \]
    La moyenne de la série regroupée est $\bar{x} = 14,80$ (au centième).
\end{enumerate}
\end{corrige}

\begin{corrige}[Exercice 5 — Tailles d'élèves : tableau, classe modale et moyenne]
\begin{enumerate}
    \item Fréquences : $\dfrac{8}{50} = 16\%$ ; $\dfrac{18}{50} = 36\%$ ; $\dfrac{16}{50} = 32\%$ ; $\dfrac{8}{50} = 16\%$. (Vérification : $16 + 36 + 32 + 16 = 100\%$.)
    Centres des classes : $155\,; 165\,; 175\,; 185$.
    \begin{center}
    \begin{tabular}{|c|c|c|c|c|}\hline
    Classes (cm) & $[150 ; 160[$ & $[160 ; 170[$ & $[170 ; 180[$ & $[180 ; 190[$ \\ \hline
    Effectifs & $8$ & $18$ & $16$ & $8$ \\ \hline
    Fréquences (\%) & $16$ & $36$ & $32$ & $16$ \\ \hline
    Centres & $155$ & $165$ & $175$ & $185$ \\ \hline
    \end{tabular}
    \end{center}
    \item La classe modale est $[160 ; 170[$ (effectif $18$, le plus grand). Interprétation : la plupart des élèves de 3ème de ce collège mesurent entre $160$ cm et $170$ cm.
    \item Moyenne :
    \[ \bar{x} = \dfrac{8 \times 155 + 18 \times 165 + 16 \times 175 + 8 \times 185}{50} = \dfrac{1240 + 2970 + 2800 + 1480}{50} = \dfrac{8490}{50} = 169,8. \]
    La taille moyenne des élèves est d'environ $169,8$ cm.
\end{enumerate}
\end{corrige}

\begin{corrige}[Exercice 6 — Tableau incomplet et diagramme circulaire]
\begin{enumerate}
    \item \textbf{Effectif de la modalité « 2 » :} sa fréquence est $25\%$, donc $200 \times 0,25 = 50$ familles. Vérification par l'effectif total : $40 + 60 + 50 + 30 + 20 = 200$. \checkmark
    
    \textbf{Fréquences :} $0$ enfant : $\dfrac{40}{200} = 20\%$ ; $3$ enfants : $\dfrac{30}{200} = 15\%$ ; $4$ et plus : $\dfrac{20}{200} = 10\%$. (Somme : $20 + 30 + 25 + 15 + 10 = 100\%$. \checkmark)
    
    \textbf{Angles} (formule : angle $= 360^\circ \times$ fréquence) :
    \begin{itemize}
        \item $0$ enfant : $360 \times 0,20 = 72^\circ$ (cohérent avec l'énoncé) ;
        \item $1$ enfant : $360 \times 0,30 = 108^\circ$ ;
        \item $2$ enfants : $360 \times 0,25 = 90^\circ$ ;
        \item $3$ enfants : $360 \times 0,15 = 54^\circ$ (cohérent) ;
        \item $4$ et plus : $360 \times 0,10 = 36^\circ$.
    \end{itemize}
    Vérification : $72 + 108 + 90 + 54 + 36 = 360^\circ$. \checkmark
    \begin{center}
    \begin{tabular}{|c|c|c|c|c|c|}\hline
    Nombre d'enfants & $0$ & $1$ & $2$ & $3$ & $4$ et plus \\ \hline
    Effectifs & $40$ & $60$ & $50$ & $30$ & $20$ \\ \hline
    Fréquences (\%) & $20$ & $30$ & $25$ & $15$ & $10$ \\ \hline
    Angles ($^\circ$) & $72$ & $108$ & $90$ & $54$ & $36$ \\ \hline
    \end{tabular}
    \end{center}
    \item \textbf{Construction :} tracer un cercle et un rayon de référence. À l'aide du rapporteur, mesurer successivement les secteurs : $72^\circ$ (0 enfant), puis $108^\circ$ (1 enfant), puis $90^\circ$ (2 enfants), puis $54^\circ$ (3 enfants) ; le dernier secteur mesure automatiquement $36^\circ$. Légender chaque secteur.
    \item « Au moins $2$ enfants » signifie $2$, $3$ ou $4$ et plus : $25\% + 15\% + 10\% = 50\%$. La moitié des familles du quartier a au moins $2$ enfants scolarisés.
\end{enumerate}
\end{corrige}

\begin{corrige}[Exercice 7 — Notes de mathématiques : moyenne et pourcentage]
Effectif total : $3 + 7 + 12 + 10 + 3 = 35$ élèves.
\begin{enumerate}
    \item Centres des classes : $2\,; 6\,; 10\,; 14\,; 18$.
    \[ \bar{x} = \dfrac{3 \times 2 + 7 \times 6 + 12 \times 10 + 10 \times 14 + 3 \times 18}{35} = \dfrac{6 + 42 + 120 + 140 + 54}{35} = \dfrac{362}{35} \approx 10,34. \]
    La moyenne de la classe est d'environ $10,34/20$.
    \item Les élèves ayant une note $\ge 8$ sont dans les classes $[8 ; 12[$, $[12 ; 16[$ et $[16 ; 20]$ : $12 + 10 + 3 = 25$ élèves.
    \[ \dfrac{25}{35} \times 100 \approx 71,4\%. \]
    Environ $71,4\%$ des élèves ont obtenu une note supérieure ou égale à $8$.
    \item Les notes $\ge 10$ se trouvent dans les classes $[12 ; 16[$ et $[16 ; 20]$, plus une partie inconnue de la classe $[8 ; 12[$. On sait seulement qu'au minimum $10 + 3 = 13$ élèves ont au moins $10$ (ceux des deux dernières classes), et au maximum $12 + 10 + 3 = 25$. La moitié de la classe vaut $17,5$ élèves. Comme le tableau regroupé ne précise pas la répartition à l'intérieur de $[8 ; 12[$, on ne peut pas affirmer avec certitude que plus de $17$ élèves ont au moins $10$ : l'affirmation du professeur est plausible mais \textbf{non vérifiable exactement} à partir du tableau regroupé. (C'est une limite du regroupement en classes : on perd l'information fine.)
\end{enumerate}
\end{corrige}

\begin{corrige}[Exercice 8 — Histogramme et polygone des effectifs : distances domicile-collège]
\begin{enumerate}
    \item Toutes les classes ont la même amplitude ($2$ km), donc les hauteurs des rectangles sont proportionnelles aux effectifs. Avec $1$ cm$^2$ pour $1$ élève et une base de $2$ cm par classe, les hauteurs sont : $7,5$ cm ; $10$ cm ; $6$ cm ; $4$ cm ; $2,5$ cm.
    \begin{popfigure}
    \begin{center}
    \begin{tikzpicture}[x=1cm,y=1cm]
        % Axes
        \draw[->] (0,0) -- (10.5,0) node[below right]{distance (km)};
        \draw[->] (0,0) -- (0,11) node[above]{effectif};
        % Graduations x
        \foreach \x in {0,2,4,6,8,10} \draw (\x,0.1) -- (\x,-0.1) node[below]{\x};
        % Graduations y
        \foreach \y in {2,4,6,8,10} \draw (0.1,\y) -- (-0.1,\y) node[left]{\y};
        % Histogramme (aire = effectif)
        \fill[popblueL,draw=popblue] (0,0) rectangle (2,7.5);
        \fill[popblueL,draw=popblue] (2,0) rectangle (4,10);
        \fill[popblueL,draw=popblue] (4,0) rectangle (6,6);
        \fill[popblueL,draw=popblue] (6,0) rectangle (8,4);
        \fill[popblueL,draw=popblue] (8,0) rectangle (10,2.5);
        % Polygone des effectifs (centres des classes, hauteurs = effectifs)
        \draw[poporange,thick,mark=*,mark size=1.5pt] plot coordinates {(1,15) (3,20) (5,12) (7,8) (9,5)};
    \end{tikzpicture}
    \end{center}
    \legende{Histogramme (bleu, aire proportionnelle aux effectifs) et polygone des effectifs (orange, sommets aux centres des classes).}
    \end{popfigure}
    \item Le polygone des effectifs relie les points de coordonnées $(1\,;15)$, $(3\,;20)$, $(5\,;12)$, $(7\,;8)$, $(9\,;5)$, c'est-à-dire les milieux des sommets des rectangles de l'histogramme usuel (hauteur = effectif), tracé en orange ci-dessus.
    \item Graphiquement, le rectangle le plus haut (ou de plus grande aire) correspond à la classe $[2 ; 4[$. Vérification par le calcul : l'effectif de cette classe est $20$, le plus grand des effectifs ($20 > 15 > 12 > 8 > 5$). La classe modale est donc $[2 ; 4[$.
    \item « Moins de $4$ km » regroupe les classes $[0 ; 2[$ et $[2 ; 4[$ : $15 + 20 = 35$ élèves.
    \[ \dfrac{35}{60} = \dfrac{7}{12} \approx 0,583 \quad \text{soit environ } 58,3\%. \]
    Près de $6$ élèves sur $10$ habitent à moins de $4$ km du collège.
\end{enumerate}
\end{corrige}

\begin{corrige}[Exercice 9 — Gestion de stock : ventes d'un commerçant au Marché Central]
\begin{enumerate}
    \item Centres des classes : $60\,; 80\,; 100\,; 120\,; 140$ (en kFCFA). Effectif total : $4 + 8 + 10 + 5 + 3 = 30$ jours.
    \[ \bar{x} = \dfrac{4 \times 60 + 8 \times 80 + 10 \times 100 + 5 \times 120 + 3 \times 140}{30} = \dfrac{240 + 640 + 1000 + 600 + 420}{30} = \dfrac{2900}{30} \approx 96,67. \]
    La recette moyenne journalière est d'environ $97\,000$ FCFA (arrondie au millier).
    \item La classe modale est $[90 ; 110[$ (effectif $10$, le plus grand). Concrètement, c'est la fourchette de recettes que Monsieur Kamga réalise le plus souvent : un tiers de ses journées (10 jours sur 30) lui rapportent entre $90\,000$ et $110\,000$ FCFA.
    \item Recette mensuelle prévisionnelle : $96,67 \times 30 \approx 2900$ kFCFA, soit $2\,900\,000$ FCFA. Stock à prévoir :
    \[ 1,2 \times 2\,900\,000 = 3\,480\,000 \text{ FCFA}. \]
    \item \textbf{Conclusion rédigée :} Monsieur Kamga peut tabler sur une recette mensuelle d'environ $2,9$ millions de FCFA, ce qui justifie un stock de $3\,480\,000$ FCFA avec une marge de sécurité de $20\%$. Toutefois, l'observation du tableau montre une dispersion notable : $7$ journées sur $30$ (classes $[50;70[$ et $[130;150[$) s'écartent fortement de la moyenne, avec des recettes allant de moins de $70\,000$ à près de $150\,000$ FCFA. La majorité des journées ($23$ sur $30$, soit environ $77\%$) se situe entre $70\,000$ et $130\,000$ FCFA, ce qui rend la prévision raisonnablement fiable, mais le commerçant doit garder une trésorerie de précaution pour absorber les journées creuses.
    
    \textbf{Barème indicatif (sur 10) :} centres des classes et moyenne (3 pts) ; classe modale et interprétation (2 pts) ; calcul du stock (2 pts) ; conclusion argumentée mentionnant la dispersion (3 pts).
    
    \textbf{Points de vigilance :} utiliser les \emph{centres} des classes et non les bornes ; convertir correctement les kFCFA en FCFA ; la conclusion doit s'appuyer sur le tableau et non sur des impressions.
\end{enumerate}
\end{corrige}

\begin{corrige}[Exercice 10 — Lecture critique : Erreur dans un journal]
\begin{enumerate}
    \item Somme des angles publiés : $130 + 110 + 80 + 50 = 370^\circ \neq 360^\circ$. La somme est \textbf{incorrecte} : un diagramme circulaire doit totaliser exactement $360^\circ$.
    \item Si l'angle de MTN ($130^\circ$) est exact, l'effectif MTN est :
    \[ \dfrac{130}{360} \times 10\,000 = \dfrac{1\,300\,000}{360} \approx 3\,611 \text{ abonnés}. \]
    Il reste $10\,000 - 3\,611 = 6\,389$ abonnés à répartir. Avec les angles publiés pour les autres : Orange : $\dfrac{110}{360} \times 10\,000 \approx 3\,056$ ; Nexttel : $\dfrac{80}{360} \times 10\,000 \approx 2\,222$ ; Camtel : $\dfrac{50}{360} \times 10\,000 \approx 1\,389$. Total : $3\,611 + 3\,056 + 2\,222 + 1\,389 = 10\,278 \neq 10\,000$ : les effectifs calculés à partir des angles publiés dépassent l'effectif total, ce qui confirme l'erreur.
    \item En supposant les effectifs issus des angles exacts en proportion mais la somme fausse, l'erreur du graphiste est un excès de $10^\circ$ ($370^\circ$ au lieu de $360^\circ$), probablement dû à un arrondi systématique à la dizaine supérieure. Une répartition cohérente consiste à réduire de $10^\circ$ l'un des secteurs, par exemple en corrigeant les arrondis : si les fréquences réelles sont MTN $36\%$, Orange $30\%$, Nexttel $22\%$, Camtel $12\%$, les angles corrects sont :
    \[ \text{Orange : } 0,30 \times 360 = 108^\circ \ ;\quad \text{Nexttel : } 0,22 \times 360 \approx 79^\circ \ ;\quad \text{Camtel : } 0,12 \times 360 \approx 43^\circ, \]
    avec MTN : $0,36 \times 360 \approx 130^\circ$, et $130 + 108 + 79 + 43 = 360^\circ$. \checkmark
    \item Vérifier la somme des angles est un contrôle de cohérence indispensable : un diagramme dont les angles ne totalisent pas $360^\circ$ déforme visuellement les parts de chaque catégorie et induit le lecteur en erreur sur la réalité des parts de marché. C'est une question d'honnêteté et de rigueur de l'information.
    
    \textbf{Barème indicatif (sur 10) :} somme et conclusion (2 pts) ; calculs d'effectifs (3 pts) ; identification de l'erreur et angles corrigés avec somme $360^\circ$ (3 pts) ; argumentation (2 pts).
    
    \textbf{Points de vigilance :} citer la propriété « la somme des angles d'un diagramme circulaire vaut $360^\circ$ » ; vérifier la cohérence effectifs/angles dans les deux sens ; les angles corrigés proposés doivent effectivement totaliser $360^\circ$.
\end{enumerate}
\end{corrige}

\begin{corrige}[Exercice 11 — Comparaison de deux classes : 3ème A vs 3ème B]
\begin{enumerate}
    \item Centres des classes : $2,5\,; 7,5\,; 12,5\,; 17,5$.
    \begin{align*}
    \bar{x}_A &= \dfrac{3 \times 2,5 + 10 \times 7,5 + 15 \times 12,5 + 7 \times 17,5}{35} = \dfrac{7,5 + 75 + 187,5 + 122,5}{35} = \dfrac{392,5}{35} \approx 11,21. \\
    \bar{x}_B &= \dfrac{2 \times 2,5 + 8 \times 7,5 + 12 \times 12,5 + 8 \times 17,5}{30} = \dfrac{5 + 60 + 150 + 140}{30} = \dfrac{355}{30} \approx 11,83.
    \end{align*}
    Moyenne de la 3ème A : $\approx 11,21/20$ ; moyenne de la 3ème B : $\approx 11,83/20$.
    \item Classe modale de la 3ème A : $[10 ; 15[$ (effectif $15$). Classe modale de la 3ème B : $[10 ; 15[$ (effectif $12$). Les deux classes ont la même classe modale.
    \item Taux de réussite (note $\ge 10$) :
    \begin{itemize}
        \item 3ème A : $15 + 7 = 22$ élèves, soit $\dfrac{22}{35} \times 100 \approx 62,9\%$ ;
        \item 3ème B : $12 + 8 = 20$ élèves, soit $\dfrac{20}{30} \times 100 \approx 66,7\%$.
    \end{itemize}
    \item \textbf{Comparaison rédigée :} Les deux classes ont la même classe modale $[10 ; 15[$, ce qui indique un profil central comparable. Cependant, la 3ème B présente une moyenne supérieure ($11,83$ contre $11,21$) et un taux de réussite plus élevé ($66,7\%$ contre $62,9\%$). De plus, la 3ème B comporte proportionnellement plus de très bonnes notes : $\dfrac{8}{30} \approx 26,7\%$ d'élèves dans $[15 ; 20]$ contre $\dfrac{7}{35} = 20\%$ en 3ème A, et moins d'élèves en difficulté dans $[0 ; 5[$ ($6,7\%$ contre $8,6\%$). Tous les indicateurs convergent : la \textbf{3ème B a un niveau globalement meilleur}, même si l'écart reste modéré (environ $0,6$ point de moyenne).
    
    \textbf{Barème indicatif (sur 10) :} moyennes avec centres de classes (3 pts) ; classes modales (1 pt) ; taux de réussite (2 pts) ; comparaison argumentée mobilisant les trois indicateurs (4 pts).
    
    \textbf{Points de vigilance :} ne pas comparer les effectifs bruts (les classes n'ont pas le même effectif total) mais des fréquences ou pourcentages ; la conclusion doit citer explicitement moyenne, classe modale et taux de réussite ; mentionner que les moyennes sont calculées à partir des centres des classes.
\end{enumerate}
\end{corrige}

\newpage

\coursec{Fiche bilan}

\begin{popfigure}
\begin{tikzpicture}[mindmap, grow cyclic, every node/.style={concept, font=\small},
  root concept/.append style={concept color=popblue, font=\bfseries},
  level 1/.append style={level distance=3.6cm, sibling angle=60},
  level 2/.append style={level distance=2.2cm, font=\scriptsize}]
\node[root concept]{Statistiques : regroupement en classes}
  child[concept color=poporange]{ node {Regroupement en classes}
    child { node {Amplitude $a$} }
    child { node {Centre $c_i=\dfrac{b_i+b_{i+1}}{2}$} }
  }
  child[concept color=popgreen]{ node {Classe modale}
    child { node {Plus grand effectif} }
  }
  child[concept color=poppurple]{ node {Fr\'equences}
    child { node {$f_i=\dfrac{n_i}{N}$} }
    child { node {$\sum f_i=1$} }
  }
  child[concept color=popgold]{ node {Moyenne}
    child { node {$\bar{x}=\dfrac{\sum n_i c_i}{N}$} }
  }
  child[concept color=popblue!70]{ node {Histogramme}
    child { node {Rectangles, aires $\propto$ effectifs} }
  }
  child[concept color=poppink]{ node {Diagrammes circulaires}
    child { node {Angle $=f_i\times 360^\circ$} }
    child { node {Semi-circulaire : $f_i\times 180^\circ$} }
  };
\end{tikzpicture}
\legende{Carte mentale de la le\c{c}on : les six id\'ees cl\'es des statistiques avec regroupement en classes.}
\end{popfigure}

\begin{bilan}
\textbf{D\'efinitions cl\'es}
\begin{itemize}
  \item \cle{Classe} : intervalle $[a\,;b[$ dans lequel on regroupe les valeurs du caract\`ere ; l'\cle{amplitude} est $b-a$.
  \item \cle{Centre de classe} : $c_i=\dfrac{\text{borne inf\'erieure}+\text{borne sup\'erieure}}{2}$.
  \item \cle{Effectif} $n_i$ : nombre de valeurs dans la classe ; \cle{effectif total} $N=\sum n_i$.
  \item \cle{Fr\'equence} d'une classe : $f_i=\dfrac{n_i}{N}$ (nombre entre $0$ et $1$, souvent exprim\'e en \% ; $\sum f_i=1$).
  \item \cle{Classe modale} : classe ayant le plus grand effectif (ou la plus grande fr\'equence).
\end{itemize}

\textbf{Formules \`a conna\^itre par c\oe ur}
\begin{center}
\begin{tabularx}{\linewidth}{|l|X|}
\hline
\rowcolor{popblueL} \textbf{Notion} & \textbf{Formule} \\
\hline
Centre de classe & $c_i=\dfrac{b_i+b_{i+1}}{2}$ \\
\hline
Fr\'equence d'une classe & $f_i=\dfrac{n_i}{N}$ \quad (en \% : $f_i\times 100$) \\
\hline
Moyenne d'une s\'erie regroup\'ee & $\bar{x}=\dfrac{n_1c_1+n_2c_2+\cdots+n_pc_p}{N}=\dfrac{\sum n_ic_i}{N}$ \\
\hline
Angle (diagramme circulaire) & $\alpha_i=f_i\times 360^\circ=\dfrac{n_i}{N}\times 360^\circ$ \\
\hline
Angle (diagramme semi-circulaire) & $\alpha_i=f_i\times 180^\circ$ \\
\hline
\end{tabularx}
\end{center}

\textbf{M\'ethodes express}
\begin{itemize}
  \item \emph{Calculer une moyenne} : ajouter une colonne << centre $c_i$ >> et une colonne << $n_i\times c_i$ >> au tableau, additionner, diviser par $N$.
  \item \emph{Tracer un histogramme} : classes contigu\"es en abscisse, rectangles dont les \cle{aires} sont proportionnelles aux effectifs (hauteur = effectif si les amplitudes sont \'egales).
  \item \emph{Tracer un diagramme circulaire} : calculer chaque angle $\alpha_i=f_i\times 360^\circ$, v\'erifier que la somme vaut $360^\circ$, reporter au rapporteur.
  \item \emph{Interpr\'eter} : toujours conclure par une phrase en lien avec la situation (ex. << la majorit\'e des \'el\`eves p\`ese entre ... >>).
\end{itemize}
\end{bilan}

\begin{aretenir}[Checklist avant le BEPC]
Avant l'\'epreuve, v\'erifie que tu sais :
\begin{itemize}
  \item[$\square$] regrouper des donn\'ees en classes d'amplitude donn\'ee et compter les effectifs ;
  \item[$\square$] calculer le centre $c_i$ de chaque classe ;
  \item[$\square$] calculer une fr\'equence $f_i=\dfrac{n_i}{N}$ et l'exprimer en pourcentage ;
  \item[$\square$] v\'erifier que la somme des fr\'equences vaut $1$ (ou $100\,\%$) ;
  \item[$\square$] identifier la classe modale et l'interpr\'eter par une phrase ;
  \item[$\square$] calculer la moyenne $\bar{x}=\dfrac{\sum n_ic_i}{N}$ avec les centres de classes ;
  \item[$\square$] compl\'eter un tableau statistique \`a trous (effectifs, fr\'equences, centres) ;
  \item[$\square$] construire un histogramme avec des axes gradu\'es et un titre ;
  \item[$\square$] calculer les angles d'un diagramme circulaire ($f_i\times 360^\circ$) ou semi-circulaire ($f_i\times 180^\circ$) ;
  \item[$\square$] reporter les angles au rapporteur proprement et l\'egender le diagramme ;
  \item[$\square$] r\'ediger une conclusion justifi\'ee \`a partir du tableau ou du graphique ;
  \item[$\square$] ne pas confondre effectif (un nombre entier) et fr\'equence (entre $0$ et $1$).
\end{itemize}
\end{aretenir}

\findecours

\end{document}
